
\Data\ID{1003019101}
\Updated{2021-04-08 07:04:56}
\Level{A}
\SubArea{Probability}
\LANGen{
\Question{
An urn contains chips with the numbers from \( 1 \) to \( 30 \).  What is the probability that the chip with a prime number is drawn?}
{\( \frac13\dot{=}\,0.3333 \)}{\( \frac{11}{30}\dot{=}\,0.3667 \)}{\( \frac25 =0.4 \)}{\( \frac3{10}=0.3 \)}{}{}}
\LANGpl{
\Question{
Pudełko zawiera żetony o numerach od \( 1 \) do \( 30 \). Jakie jest prawdopodobieństwo, że wylosujemy żeton z liczbą pierwszą?}
{\( \frac13\dot{=}\,0{,}3333 \)}{\( \frac{11}{30}\dot{=}\,0{,}3667 \)}{\( \frac25 =0{,}4 \)}{\( \frac3{10}=0{,}3 \)}{}{}}
\LANGcs{
\Question{
V krabici jsou žetony s čísly od \( 1 \) do \( 30 \).  S jakou pravděpodobností vytáhneme prvočíslo?}
{\( \frac13\dot{=}\,0{,}3333 \)}{\( \frac{11}{30}\dot{=}\,0{,}3667 \)}{\( \frac25 =0{,}4 \)}{\( \frac3{10}=0{,}3 \)}{}{}}
\LANGsk{
\Question{
V krabici sú žetóny s číslami od \( 1 \) do \( 30 \).  S akou pravdepodobnosťou vytiahneme prvočíslo?}
{\( \frac13\dot{=}\,0{,}3333 \)}{\( \frac{11}{30}\dot{=}\,0{,}3667 \)}{\( \frac25 =0{,}4 \)}{\( \frac3{10}=0{,}3 \)}{}{}}
\LANGes{
\Question{
En una caja hay chips con números del \( 1 \) al \( 30 \). ¿Cuál es la probabilidad de sacar un número primo?}
{\( \frac13\dot{=}\,0{,}3333 \)}{\( \frac{11}{30}\dot{=}\,0{,}3667 \)}{\( \frac25 =0{,}4 \)}{\( \frac3{10}=0{,}3 \)}{}{}}
\End

\Data\ID{1003019102}
\Updated{2021-01-27 03:01:29}
\Level{A}
\SubArea{Probability}
\LANGen{
\Question{
An urn contains \( 19 \) red balls and \( 9 \) blue balls. Find the least number of blue balls that need to be added to the urn to ensure higher probability of drawing a blue ball than \( 0.65 \).}
{\( 27 \)}{\( 26 \)}{\( 10 \)}{\( 0 \)}{}{}}
\LANGpl{
\Question{
W pudełku znajduje się  \( 19 \) czerwonych piłek i \( 9 \) niebieskich piłek. Określ minimalną ilość  niebieskich piłek, które  należy dodać do pudełka, aby prawdopodobieństwo wylosowania niebieskiej piłki  było większe niż \( 0{,}65 \).}
{\( 27 \)}{\( 26 \)}{\( 10 \)}{\( 0 \)}{}{}}
\LANGcs{
\Question{
V krabici je \( 19 \) červených a \( 9 \) modrých kuliček. Určete minimální počet modrých kuliček, které je třeba do krabice přidat, aby pravděpodobnost vytažení modré kuličky (při následujícím vytažení jedné kuličky) byla větší než \( 0{,}65 \).}
{\( 27 \)}{\( 26 \)}{\( 10 \)}{\( 0 \)}{}{}}
\LANGsk{
\Question{
V krabici je \( 19 \) červených a \( 9 \) modrých guliek. Určte, najmenej koľko modrých guliek treba ešte do krabice pridať, aby pri následnom vytiahnutí jednej guľky bola pravdepodobnosť vytiahnutia modrej guľky  väčšia ako \( 0{,}65 \).}
{\( 27 \)}{\( 26 \)}{\( 10 \)}{\( 0 \)}{}{}}
\LANGes{
\Question{
En una caja hay \( 19 \) bolas rojas y \( 9 \) bolas azules. Averigua la cantidad mínima de bolas azules que tenemos que añadir a la caja para que la probabilidad de sacar una bola azul sea mayor que \( 0{,}65 \).}
{\( 27 \)}{\( 26 \)}{\( 10 \)}{\( 0 \)}{}{}}
\End

\Data\ID{1003019103}
\Updated{2021-01-27 03:01:13}
\Level{A}
\SubArea{Probability}
\LANGen{
\Question{
There are \( 30 \) students in the class, one of them is Adam. The teacher picks randomly three students to be tested. What is the probability that Adam is among them?}
{\( \frac{\binom{29}2}{\binom{30}3}=0{.}1 \)}{\( \frac{\binom{29}2}{\binom{30}2}\dot{=}\,0{.}9333 \)}{\( \frac{\binom{29}3}{\binom{30}3}=0{.}9 \)}{\( \frac{\binom31\binom{27}2}{\binom{30}{3}}\dot{=}\,0{.}2594 \)}{}{}}
\LANGpl{
\Question{
W klasie jest  \( 30 \) uczniów, jednym z nich jest Adam. Nauczyciel wybiera losowo trzech uczniów do odpowiedzi. Jakie jest prawdopodobieństwo, że Adam jest wśród nich?}
{\( \frac{\binom{29}2}{\binom{30}3}=0{,}1 \)}{\( \frac{\binom{29}2}{\binom{30}2}\dot{=}\,0{,}9333 \)}{\( \frac{\binom{29}3}{\binom{30}3}=0{,}9 \)}{\( \frac{\binom31\binom{27}2}{\binom{30}{3}}\dot{=}\,0{,}2594 \)}{}{}}
\LANGcs{
\Question{
Ve třídě je \( 30 \) žáků, jedním z nich je i Adam. Učitel náhodně vyvolá k tabuli tři žáky. Jaká je pravděpodobnost, že mezi nimi bude i Adam?}
{\( \frac{\binom{29}2}{\binom{30}3}=0{,}1 \)}{\( \frac{\binom{29}2}{\binom{30}2}\dot{=}\,0{,}9333 \)}{\( \frac{\binom{29}3}{\binom{30}3}=0{,}9 \)}{\( \frac{\binom31\binom{27}2}{\binom{30}{3}}\dot{=}\,0{,}2594 \)}{}{}}
\LANGsk{
\Question{
V triede je \( 30 \) žiakov, jedným z nich je aj Adam. Učiteľ náhodne vyvolá k tabuli odpovedať troch žiakov. Aká je pravdepodobnosť, že Adam bude jedným z nich?}
{\( \frac{\binom{29}2}{\binom{30}3}=0{,}1 \)}{\( \frac{\binom{29}2}{\binom{30}2}\dot{=}\,0{,}9333 \)}{\( \frac{\binom{29}3}{\binom{30}3}=0{,}9 \)}{\( \frac{\binom31\binom{27}2}{\binom{30}{3}}\dot{=}\,0{,}2594 \)}{}{}}
\LANGes{
\Question{
En una clase hay \( 30 \) alumnos, uno de ellos de llama Adam. El profesor va a sacar tres alumnos para ir a la pizarra. ¿Qué es la probabilidad de que le saque a Adam?}
{\( \frac{\binom{29}2}{\binom{30}3}=0{,}1 \)}{\( \frac{\binom{29}2}{\binom{30}2}\dot{=}\,0{,}9333 \)}{\( \frac{\binom{29}3}{\binom{30}3}=0{,}9 \)}{\( \frac{\binom31\binom{27}2}{\binom{30}{3}}\dot{=}\,0{,}2594 \)}{}{}}
\End

\Data\ID{1003019104}
\Updated{2021-01-27 03:01:45}
\Level{A}
\SubArea{Probability}
\LANGen{
\Question{
Four dice are thrown. What is the probability that the sum is \( 6 \)?}
{\( \frac{10}{6^4}\dot{=}\,0{.}0077 \)}{\( \frac1{21}\dot{=}\,0{.}0476 \)}{\( \frac1{6^4}\dot{=}\,0{.}0008 \)}{\( \frac2{6^4}\dot{=}\,0{.}0015 \)}{}{}}
\LANGpl{
\Question{
Rzucono czterema kostkami.  Jakie jest prawdopodobieństwo, że suma wyrzuconych liczb  jest równa  \( 6 \)?}
{\( \frac{10}{6^4}\dot{=}\,0{,}0077 \)}{\( \frac1{21}\dot{=}\,0{,}0476 \)}{\( \frac1{6^4}\dot{=}\,0{,}0008 \)}{\( \frac2{6^4}\dot{=}\,0{,}0015 \)}{}{}}
\LANGcs{
\Question{
Jaká je pravděpodobnost, že při hodu čtyřmi hracími kostkami padne součet čísel \( 6 \)?}
{\( \frac{10}{6^4}\dot{=}\,0{,}0077 \)}{\( \frac1{21}\dot{=}\,0{,}0476 \)}{\( \frac1{6^4}\dot{=}\,0{,}0008 \)}{\( \frac2{6^4}\dot{=}\,0{,}0015 \)}{}{}}
\LANGsk{
\Question{
Aká je pravdepodobnosť, že pri hode štyrmi hracími kockami padne súčet čísel \( 6 \)?}
{\( \frac{10}{6^4}\dot{=}\,0{,}0077 \)}{\( \frac1{21}\dot{=}\,0{,}0476 \)}{\( \frac1{6^4}\dot{=}\,0{,}0008 \)}{\( \frac2{6^4}\dot{=}\,0{,}0015 \)}{}{}}
\LANGes{
\Question{
Vamos a tirar cuatro dados. ¿Qué es la probabilidad de que la suma de los cuatro números sea \( 6 \)?}
{\( \frac{10}{6^4}\dot{=}\,0{,}0077 \)}{\( \frac1{21}\dot{=}\,0{,}0476 \)}{\( \frac1{6^4}\dot{=}\,0{,}0008 \)}{\( \frac2{6^4}\dot{=}\,0{,}0015 \)}{}{}}
\End

\Data\ID{1003029201}
\Updated{2021-01-29 11:01:50}
\Level{A}
\SubArea{Probability}
\LANGen{
\Question{
Three dice are thrown together. Let \( A \) be the event “the sum is \( 5 \)” and \( B \) be the event “the sum is \( 16 \)”. 
Which of the following statements is true?}
{Events \( A \) and \( B \) have the same probability of occurrence.}{Event \( A \) is more likely to occur than event \( B \).}{Event \( B \) is more likely to occur than event \( A \).}{}{}{}}
\LANGpl{
\Question{
Trzy kostki zostały rzucone jednocześnie. Niech \( A \) będzie zdarzeniem “suma wynosi \( 5 \)” a \( B \) będzie zdarzeniem “suma wynosi \( 16 \)”. 
Które z poniższych stwierdzeń jest prawdziwe?}
{Zdarzenia \( A \) i \( B \) mają takie samo prawdopodobieństwo wystąpienia.}{Zdarzenie \( A \) jest bardziej prawdopodobne niż zdarzenie \( B \).}{Zdarzenie\( B \) jest bardziej prawdopodobne niż zdarzenie \( A \).}{}{}{}}
\LANGcs{
\Question{
Házíme 3 kostkami a pozorujeme součet bodů, které na kostkách padnou. Označme jev \( A \) to, že “součet bodů je \( 5 \)” a jev \( B \) to, že “součet bodů je \( 16 \)”. Vyberte pravdivé tvrzení.}
{Jevy \( A \) i \( B \) jsou stejně pravděpodobné.}{Jev \( A \) je pravděpodobnější než jev \( B \).}{Jev \( B \) je pravděpodobnější než jev \( A \).}{}{}{}}
\LANGsk{
\Question{
Hádžeme 3 kockami a pozorujeme súčet bodov, ktoré na kockách padnú. Označme jav \( A \) to, že “súčet bodov je \( 5 \)” a jav \( B \) to, že “súčet bodov je \( 16 \)”. Vyberte pravdivé tvrdenie.}
{Javy \( A \) aj \( B \) sú rovnako pravdepodobné.}{Jav \( A \) je pravdepodobnejší ako jav \( B \).}{Jav \( B \) je pravdepodobnejší ako jav \( A \).}{}{}{}}
\LANGes{
\Question{
Vamos a tirar tres dados y ver la suma de los tres números. El fenómeno "la suma de los puntos es \( 5 \)” vamos a llamarlo \( A \), y el fenómeno \( B \) es "la suma de los puntos es \( 16 \)”.  Elige la declaración correcta.}
{Los fenómenos \( A \) y \( B \) tienen la misma probabilidad.}{\( A \) es más probable que \( B \).}{\( B \) es más probable que \( A \).}{}{}{}}
\End

\Data\ID{1003029202}
\Updated{2021-02-03 11:02:28}
\Level{A}
\SubArea{Probability}
\LANGen{
\Question{
In a set of \( 100 \) items, \( 15 \) are defective. We pick randomly \( 10 \) items from this set. First eight items were not defective. Find the probability that the ninth item selected is not defective too. Results are rounded to two decimal places.}
{\( \frac{77}{92}=0{.}84 \)}{\( \frac{85}{92}=0{.}92\)}{\( \frac{15}{92}=0{.}16 \)}{\( \frac7{92}=0{.}08 \)}{}{}}
\LANGpl{
\Question{
W zbiorze  \( 100 \) rzeczy, \( 15 \) jest wadliwych. Ze zbioru losowo wybieramy \( 10 \) rzeczy. Pierwszych osiem rzeczy nie posiada wady. Jakie jest prawdopodobieństwo, że dziewiąta wybrana rzecz również nie jest wadliwa. Wyniki zaokrąglono do dwóch miejsc po przecinku.}
{\( \frac{77}{92}=0{,}84 \)}{\( \frac{85}{92}=0{,}92\)}{\( \frac{15}{92}=0{,}16 \)}{\( \frac7{92}=0{,}08 \)}{}{}}
\LANGcs{
\Question{
Mezi  \( 100 \) výrobky je \( 15 \) zmetků. Postupně z nich náhodně vybereme \( 10 \) ke kontrole.  Prvních osm vybraných výrobků bylo dobrých. Jaká je pravděpodobnost, že ani devátý vybraný výrobek nebude zmetek? Výsledky jsou zaokrouhlené na dvě desetinná místa.}
{\( \frac{77}{92}=0{,}84 \)}{\( \frac{85}{92}=0{,}92\)}{\( \frac{15}{92}=0{,}16 \)}{\( \frac7{92}=0{,}08 \)}{}{}}
\LANGsk{
\Question{
Medzi  \( 100 \) výrobkami je \( 15 \) nepodarkov. Postupne z nich náhodne vyberieme \( 10 \) na kontrolu.  Prvých osem vybraných výrobkov bolo dobrých. Aká je pravdepodobnosť, že deviaty vybratý výrobok nebude nepodarok? Výsledky sú zaokrúhlené na dve desatinné miesta.}
{\( \frac{77}{92}=0{,}84 \)}{\( \frac{85}{92}=0{,}92\)}{\( \frac{15}{92}=0{,}16 \)}{\( \frac7{92}=0{,}08 \)}{}{}}
\LANGes{
\Question{
Entre \( 100 \) productos hay \( 15 \) defectuosos. Vamos a elegir  \( 10 \) productos al azar para controlar.  Primeros ocho productos sacados fueros sin defecto. ¿Qué es la probabilidad de que el noveno producto sea sin defecto? Aproxima el resultado a dos cifras decimales.}
{\( \frac{77}{92}=0{,}84 \)}{\( \frac{85}{92}=0{,}92\)}{\( \frac{15}{92}=0{,}16 \)}{\( \frac7{92}=0{,}08 \)}{}{}}
\End

\Data\ID{1003041601}
\Updated{2021-02-10 09:02:35}
\Level{A}
\SubArea{Probability}
\LANGen{
\Question{
The wooden cube with the edges of length \( 5\,\mathrm{cm} \) has faces painted in blue.  Suppose we cut the cube into small unit cubes (the edge length is \( 1\,\mathrm{cm}\)) and select one of the unit cubes at random. What is the probability that the selected cube has at least two faces painted in blue?}
{\( 0{.}352 \)}{\( 0{.}288 \)}{\( 0{.}480 \)}{\( 0{.}432 \)}{}{}}
\LANGpl{
\Question{
Drewniany sześcian ma krawędzie o długości \( 5\,\mathrm{cm} \), ściany sześcianu mają niebieski kolor.  Dany sześcian  dzielimy na małe sześciany  (krawędź ma długość \( 1\,\mathrm{cm}\)) i losowo wybieramy jeden mały sześcian. Jakie jest prawdopodobieństwo, że wybierzemy sześcian, którego przynajmniej dwie ściany są niebieskie?}
{\( 0{,}352 \)}{\( 0{,}288 \)}{\( 0{,}480 \)}{\( 0{,}432 \)}{}{}}
\LANGcs{
\Question{
Dřevěná krychle s hranou o velikosti \( 5\,\mathrm{cm} \) je natřená na modro. Rozřežeme ji na jednotkové krychle (s hranou o velikosti \( 1\,\mathrm{cm}\)). Určete pravděpodobnost, že při náhodném výběru z nich vybereme krychli s alespoň dvěma modrými stěnami.}
{\( 0{,}352 \)}{\( 0{,}288 \)}{\( 0{,}480 \)}{\( 0{,}432 \)}{}{}}
\LANGsk{
\Question{
Drevená kocka s hranou veľkosti \( 5\,\mathrm{cm} \) je natretá na modro. Rozrežeme ju na jednotkové kocky (s hranou \( 1\,\mathrm{cm}\)). Určte pravdepodobnosť, že pri náhodnom výbere z nich vyberieme kocku s aspoň dvoma modrými stenami.}
{\( 0{,}352 \)}{\( 0{,}288 \)}{\( 0{,}480 \)}{\( 0{,}432 \)}{}{}}
\LANGes{
\Question{
La superficie de un cubo de madera es azul. El lado del cubo mide\( 5\,\mathrm{cm} \). Vamos a cortarlo en cubos unitarios (cuyo lado mide \( 1\,\mathrm{cm}\)). Calcula la probabilidad de que al eligir al azar vamos a coger un cubo con por lo mínimo dos lados azules.}
{\( 0{,}352 \)}{\( 0{,}288 \)}{\( 0{,}480 \)}{\( 0{,}432 \)}{}{}}
\End

\Data\ID{1003041701}
\Updated{2021-02-03 02:02:13}
\Level{A}
\SubArea{Probability}
\LANGen{
\Question{
Find the probability that a two-digit number selected at random is even or a power of \( 2 \).}
{\( \frac{45}{90}=0{.}5 \)}{\( \frac{48}{90}\,\dot{=}\,0{.}5333 \)}{\( \frac{42}{90}\,\dot{=}\,0{.}4667 \)}{\( \frac{45}{90}\cdot\frac3{90}\,\dot{=}\,0{.}1667 \)}{}{}}
\LANGpl{
\Question{
Jakie jest prawdopodobieństwo, że wybrana losowo dwucyfrowa liczba jest parzysta lub jest potęgą   \( 2 \)?}
{\( \frac{45}{90}=0{,}5 \)}{\( \frac{48}{90}\,\dot{=}\,0{,}5333 \)}{\( \frac{42}{90}\,\dot{=}\,0{,}4667 \)}{\( \frac{45}{90}\cdot\frac3{90}\,\dot{=}\,0{,}1667 \)}{}{}}
\LANGcs{
\Question{
Určete pravděpodobnost, že náhodně vybrané dvojciferné číslo bude sudé anebo mocninou čísla \( 2 \).}
{\( \frac{45}{90}=0{,}5 \)}{\( \frac{48}{90}\,\dot{=}\,0{,}5333 \)}{\( \frac{42}{90}\,\dot{=}\,0{,}4667 \)}{\( \frac{45}{90}\cdot\frac3{90}\,\dot{=}\,0{,}1667 \)}{}{}}
\LANGsk{
\Question{
Určte pravdepodobnosť, že náhodne vybrané dvojciferné číslo bude párne alebo mocnina čísla \( 2 \).}
{\( \frac{45}{90}=0{,}5 \)}{\( \frac{48}{90}\,\dot{=}\,0{,}5333 \)}{\( \frac{42}{90}\,\dot{=}\,0{,}4667 \)}{\( \frac{45}{90}\cdot\frac3{90}\,\dot{=}\,0{,}1667 \)}{}{}}
\LANGes{
\Question{
Calcula la probabilidad de que un número de dos cifras elegido al azar sea par o potencia del número \( 2 \).}
{\( \frac{45}{90}=0{,}5 \)}{\( \frac{48}{90}\,\dot{=}\,0{,}5333 \)}{\( \frac{42}{90}\,\dot{=}\,0{,}4667 \)}{\( \frac{45}{90}\cdot\frac3{90}\,\dot{=}\,0{,}1667 \)}{}{}}
\End

\Data\ID{9000154808}
\Updated{2021-02-17 10:02:53}
\Level{A}
\SubArea{Probability}
\IMG{9000154808_dices_0.png}
\LANGen{
\Question{
Merry Men, the outlaws who followed Robin Hood, play the dice game. Little John wins if he
gets the sum \(8\) 
by rolling two dices. Find the probability that he wins. Round your answer to three
decimal places.}
{\(0{.}139\)}{\(0{.}194\)}{\(0{.}806\)}{\(0{.}778\)}{}{}}
\LANGpl{
\Question{
Towarzysze Robin Hooda grają w kości. Mały John wygra, jeśli wyrzuci na obu kostkach sum licczb
 \(8\). Jakie jest prawdopodobieństwo jego wygranej?  Zaokrągli  odpowiedź do trzech miejsc po przecinku.}
{\(0{,}139\)}{\(0{,}194\)}{\(0{,}806\)}{\(0{,}778\)}{}{}}
\LANGcs{
\Question{
Malý John hraje kostky proti Robinu Hoodovi. K výhře mu chybí, aby
při hodu dvěma kostkami padl součet osm. Jaká je pravděpodobnost,
že porazí Robina již prvním hodem? Výsledek zaokrouhlete na
\(3\) 
desetinná místa.}
{\(0{,}139\)}{\(0{,}194\)}{\(0{,}806\)}{\(0{,}778\)}{}{}}
\LANGsk{
\Question{
Malý John hrá hru s kockami proti Robinovi Hoodovi. Vyhrá, ak pri hode dvoma kockami padne súčet \(8\). Aká je pravdepodobnosť, že malý John vyhrá? Výsledok zaokrúhlite na tri desatinné miesta.}
{\(0{,}139\)}{\(0{,}194\)}{\(0{,}806\)}{\(0{,}778\)}{}{}}
\LANGes{
\Question{
El pequeño Juan juega a dados contra Robin hood. Si tira una suma de ocho, puede ganar. ¿Qué probabilidad hay de que le vence a Robin Hood con su primer tiro? Aproxima el resultado a 
\(3\) cifras decimales.}
{\(0{,}139\)}{\(0{,}194\)}{\(0{,}806\)}{\(0{,}778\)}{}{}}
\End

\Data\ID{1003019201}
\Updated{2021-01-27 03:01:24}
\Level{B}
\SubArea{Probability}
\LANGen{
\Question{
Four dice are thrown. What is the probability that the sum is greater than \( 5 \)?}
{\( 1-\frac5{6^4}\dot{=}\,0{.}9961 \)}{\( \frac5{6^4}\dot{=}\,0{.}0039 \)}{\( \frac{19}{24}\dot{=}\,0{.}7917 \)}{\( 1-\frac1{6^4}\dot{=}\,0{.}9992 \)}{}{}}
\LANGpl{
\Question{
Rzucono czterema kostkami.  Jakie jest prawdopodobieństwo, że suma wyrzuconych liczb  jest większa od  \( 5 \)?}
{\( 1-\frac5{6^4}\dot{=}\,0{,}9961 \)}{\( \frac5{6^4}\dot{=}\,0{,}0039 \)}{\( \frac{19}{24}\dot{=}\,0{,}7917 \)}{\( 1-\frac1{6^4}\dot{=}\,0{,}9992 \)}{}{}}
\LANGcs{
\Question{
Jaká je pravděpodobnost, že při hodu čtyřmi hracími kostkami padne součet čísel větší než \( 5 \)?}
{\( 1-\frac5{6^4}\dot{=}\,0{,}9961 \)}{\( \frac5{6^4}\dot{=}\,0{,}0039 \)}{\( \frac{19}{24}\dot{=}\,0{,}7917 \)}{\( 1-\frac1{6^4}\dot{=}\,0{,}9992 \)}{}{}}
\LANGsk{
\Question{
Aká je pravdepodobnosť, že pri hode štyrmi hracími kockami bude súčet čísel väčší ako \( 5 \)?}
{\( 1-\frac5{6^4}\dot{=}\,0{,}9961 \)}{\( \frac5{6^4}\dot{=}\,0{,}0039 \)}{\( \frac{19}{24}\dot{=}\,0{,}7917 \)}{\( 1-\frac1{6^4}\dot{=}\,0{,}9992 \)}{}{}}
\LANGes{
\Question{
Vamos a tirar cuatro dados. ¿Qué es la probabilidad de que la suma de los cuatro números sea mayor que  \( 5 \)?}
{\( 1-\frac5{6^4}\dot{=}\,0{,}9961 \)}{\( \frac5{6^4}\dot{=}\,0{,}0039 \)}{\( \frac{19}{24}\dot{=}\,0{,}7917 \)}{\( 1-\frac1{6^4}\dot{=}\,0{,}9992 \)}{}{}}
\End

\Data\ID{1003019202}
\Updated{2021-01-29 11:01:50}
\Level{B}
\SubArea{Probability}
\LANGen{
\Question{
There are fifty apples left on a tree and ten of them have worms. We pick five apples at random. What is the probability that at least one of them is without a worm?}
{\( 1-\frac{\binom{10}{5}}{\binom{50}{5}}\dot{=}\,0{.}9999 \)}{\( 1-\frac{\binom{10}{1}}{\binom{50}{5}}\dot{=}\,1{.}0000 \)}{\( 1-\frac{\binom{10}{1}\binom{40}{4}}{\binom{50}{5}}\dot{=}\,0{.}5687 \)}{\( 1-\frac{\binom{40}{5}}{\binom{50}{5}}\dot{=}\,0{.}6894\)}{}{}}
\LANGpl{
\Question{
Na drzewie zostało pięćdziesiąt jabłek,  dziesięć jabłek jest robaczywych. Wybieramy losowo pięć jabłek. Jakie jest prawdopodobieństwo, że przynajmniej jeden z nich nie jest robaczywe?}
{\( 1-\frac{\binom{10}{5}}{\binom{50}{5}}\dot{=}\,0{,}9999 \)}{\( 1-\frac{\binom{10}{1}}{\binom{50}{5}}\dot{=}\,1{,}0000 \)}{\( 1-\frac{\binom{10}{1}\binom{40}{4}}{\binom{50}{5}}\dot{=}\,0{,}5687 \)}{\( 1-\frac{\binom{40}{5}}{\binom{50}{5}}\dot{=}\,0{,}6894 \)}{}{}}
\LANGcs{
\Question{
Na stromě zůstalo 50 jablek, ale v deseti z nich je už červík. Utrhneme ze stromu 5 náhodně vybraných jablek. Jaká je pravděpodobnost, že alespoň v jednom z nich nebude červík?}
{\( 1-\frac{\binom{10}{5}}{\binom{50}{5}}\dot{=}\,0{,}9999 \)}{\( 1-\frac{\binom{10}{1}}{\binom{50}{5}}\dot{=}\,1{,}0000 \)}{\( 1-\frac{\binom{10}{1}\binom{40}{4}}{\binom{50}{5}}\dot{=}\,0{,}5687 \)}{\( 1-\frac{\binom{40}{5}}{\binom{50}{5}}\dot{=}\,0{,}6894 \)}{}{}}
\LANGsk{
\Question{
Na strome zostalo ešte 50 jabĺk, pričom v desiatich z nich je už červík. Aká je pravdepodobnosť, že ak náhodne odtrhneme päť jabĺk, tak v aspoň jednom z nich nebude červík?}
{\( 1-\frac{\binom{10}{5}}{\binom{50}{5}}\dot{=}\,0{,}9999 \)}{\( 1-\frac{\binom{10}{1}}{\binom{50}{5}}\dot{=}\,1{,}0000 \)}{\( 1-\frac{\binom{10}{1}\binom{40}{4}}{\binom{50}{5}}\dot{=}\,0{,}5687 \)}{\( 1-\frac{\binom{40}{5}}{\binom{50}{5}}\dot{=}\,0{,}6894 \)}{}{}}
\LANGes{
\Question{
En un árbol quedaron 50 manzanas, pero en diez de ellos hay un gusano. Vamos a coger 5 manzanas al azar. ¿Qué es la probabilidad de coger por lo mínimo una manzana sin gusano?}
{\( 1-\frac{\binom{10}{5}}{\binom{50}{5}}\dot{=}\,0{,}9999 \)}{\( 1-\frac{\binom{10}{1}}{\binom{50}{5}}\dot{=}\,1{,}0000 \)}{\( 1-\frac{\binom{10}{1}\binom{40}{4}}{\binom{50}{5}}\dot{=}\,0{,}5687 \)}{\( 1-\frac{\binom{40}{5}}{\binom{50}{5}}\dot{=}\,0{,}6894 \)}{}{}}
\End

\Data\ID{1003019204}
\Updated{2021-01-29 11:01:15}
\Level{B}
\SubArea{Probability}
\LANGen{
\Question{
Inside a circle is inscribed a square. A point is chosen at random from inside the circle. What is the probability that this point is located also in the square?}
{\( \frac2{\pi}\dot{=}\,0{.}6366 \)}{\( \frac{\pi}4\dot{=}\,0{.}7854 \)}{\( \frac{\sqrt{2}}{\pi}\,\dot{=}\,0{.}4502 \)}{\( \frac{\sqrt{2}}{2\pi}\,\dot{=}\,0{.}2251 \)}{}{}}
\LANGpl{
\Question{
Kwadrat jest wpisany w okrąg.  Jakie jest prawdopodobieństwo, że losowo wybrany punkt okręgu znajduje się również w kwadracie?}
{\( \frac2{\pi}\dot{=}\,0{,}6366 \)}{\( \frac{\pi}4\dot{=}\,0{,}7854 \)}{\( \frac{\sqrt{2}}{\pi}\,\dot{=}\,0{,}4502 \)}{\( \frac{\sqrt{2}}{2\pi}\,\dot{=}\,0{,}2251 \)}{}{}}
\LANGcs{
\Question{
Do kruhu je vepsaný čtverec. Jaká je pravděpodobnost, že náhodně zvolený bod kruhu se nachází i ve čtverci?}
{\( \frac2{\pi}\dot{=}\,0{,}6366 \)}{\( \frac{\pi}4\dot{=}\,0{,}7854 \)}{\( \frac{\sqrt{2}}{\pi}\,\dot{=}\,0{,}4502 \)}{\( \frac{\sqrt{2}}{2\pi}\,\dot{=}\,0{,}2251 \)}{}{}}
\LANGsk{
\Question{
Do kruhu je vpísaný štvorec. Aká je pravdepodobnosť, že náhodne zvolený bod kruhu sa nachádza aj vo štvorci?}
{\( \frac2{\pi}\dot{=}\,0{,}6366 \)}{\( \frac{\pi}4\dot{=}\,0{,}7854 \)}{\( \frac{\sqrt{2}}{\pi}\,\dot{=}\,0{,}4502 \)}{\( \frac{\sqrt{2}}{2\pi}\,\dot{=}\,0{,}2251 \)}{}{}}
\LANGes{
\Question{
Un cuadro está inscrito dentro de una circunferencia. Vamos a eligir un punto del círculo. ¿Qué es la probabilidad de que este punto también pertenezca al cuadro?}
{\( \frac2{\pi}\dot{=}\,0{,}6366 \)}{\( \frac{\pi}4\dot{=}\,0{,}7854 \)}{\( \frac{\sqrt{2}}{\pi}\,\dot{=}\,0{,}4502 \)}{\( \frac{\sqrt{2}}{2\pi}\,\dot{=}\,0{,}2251 \)}{}{}}
\End

\Data\ID{1003019206}
\Updated{2021-02-16 02:02:53}
\Level{B}
\SubArea{Probability}
\LANGen{
\Question{
Adam and Eve met at the disco. They agreed to meet  the next day at the same location sometime between \( 1 \) p.m. and \( 2 \) p.m. Both of them will arrive independently at random times within the hour. Adam is greatly interested in the meeting, therefore he is willing to wait for Eve even up to half an hour, while Eve is willing to wait for Adam for \( 10 \) minutes. What is the probability that they will meet during that hour?}
{\( \frac{19}{36}\dot{=}\,0{.}5278 \)}{\( \frac{17}{36}\dot{=}\,0{.}4722 \)}{\( \frac{11}{36}\dot{=}\,0{.}3056 \)}{\( \frac{27}{36}=0{.}75 \)}{}{}}
\LANGpl{
\Question{
Adam i Ewa poznali się na dyskotece. Postanowili  spotkać się następnego dnia między godziną  \( 13 \) a \( 14 \). Oboje przyjeżdżają na miejsce spotkania w losowo wybranym czasie w ciągu jednej godziny. Adam jest bardziej zainteresowany spotkaniem i zamierza czekać pół godziny, natomiast Ewa zaczeka na Adama \( 10 \) minut. Jakie jest prawdopodobieństwo, że spotkają się w ciągu tej godziny?}
{\( \frac{19}{36}\dot{=}\,0{,}5278 \)}{\( \frac{17}{36}\dot{=}\,0{,}4722 \)}{\( \frac{11}{36}\dot{=}\,0{,}3056 \)}{\( \frac{27}{36}=0{,}75 \)}{}{}}
\LANGcs{
\Question{
Adam a Eva se setkali na diskotéce. Dohodli se, že se potkají druhý den mezi \( 13 \). a \( 14 \). hodinou. Adam má velký zájem o setkání, proto je ochotný čekat na Evu půl hodiny, Eva je ochotná čekat na Adama \( 10 \) minut. Jaká je pravděpodobnost, že se potkají, když jejich příchody na místo setkání jsou navzájem nezávislé a stejně pravděpodobné v průběhu celé dané hodiny?}
{\( \frac{19}{36}\dot{=}\,0{,}5278 \)}{\( \frac{17}{36}\dot{=}\,0{,}4722 \)}{\( \frac{11}{36}\dot{=}\,0{,}3056 \)}{\( \frac{27}{36}=0{,}75 \)}{}{}}
\LANGsk{
\Question{
Adam a Eva sa stretli na diskotéke. Dohodli sa, že sa stretnú na druhý deň medzi \( 13 \). a \( 14 \). hodinou. Adam má veľký záujem o stretnutie, preto je ochotný čakať na Evu pol hodiny, Eva je ochotná čakať Adama \( 10 \) minút. Aká je pravdepodobnosť, že sa stretnú, ak ich príchody na miesto stretnutia sú navzájom nezávislé a rovnako pravdepodobné v priebehu celej danej hodiny?}
{\( \frac{19}{36}\dot{=}\,0{,}5278 \)}{\( \frac{17}{36}\dot{=}\,0{,}4722 \)}{\( \frac{11}{36}\dot{=}\,0{,}3056 \)}{\( \frac{27}{36}=0{,}75 \)}{}{}}
\LANGes{
\Question{
Adam y Eva se han conocido en una discoteca. Han decidido quedar el siguiente día entre  \( 13 \). y \( 14 \). hora. Adam quiere verla mucho así que ha decidido esperar una media hora, Eva va a esperar \( 10 \) minutos. El tiempo de llegada de cada uno de ellos es aleatorio e independiente del otro. ¿Qué probabilidad hay de que Adam y Eva quedan?}
{\( \frac{19}{36}\dot{=}\,0{,}5278 \)}{\( \frac{17}{36}\dot{=}\,0{,}4722 \)}{\( \frac{11}{36}\dot{=}\,0{,}3056 \)}{\( \frac{27}{36}=0{,}75 \)}{}{}}
\End

\Data\ID{1003029301}
\Updated{2021-01-29 01:01:54}
\Level{B}
\SubArea{Probability}
\LANGen{
\Question{
Two regular dice are thrown. Find the probability that the sum is five or six. The results are rounded to two decimal places.}
{\( \frac14= 0{.}25 \)}{\( \frac5{36}\,\dot{=}\,0{.}14 \)}{\( \frac2{11}\,\dot{=}\,0{.}18 \)}{\( \frac{10}{36}\,\dot{=}\,0{.}28\)}{}{}}
\LANGpl{
\Question{
Rzucono dwoma kostkami. Jakie jest prawdopodobieństwo, że suma wyrzuconych liczb to pięć lub sześć? Zaokrągli wynik do dwóch miejsc po przecinku.}
{\( \frac14= 0{,}25 \)}{\( \frac5{36}\,\dot{=}\,0{,}14 \)}{\( \frac2{11}\,\dot{=}\,0{,}18 \)}{\( \frac{10}{36}\,\dot{=}\,0{,}28\)}{}{}}
\LANGcs{
\Question{
Hodíme dvěma kostkami. Jaká je pravděpodobnost, že výsledný součet padlých bodů bude \( 5 \) nebo \( 6 \)? Výsledek zaokrouhlete na dvě desetinná místa.}
{\( \frac14= 0{,}25 \)}{\( \frac5{36}\,\dot{=}\,0{,}14 \)}{\( \frac2{11}\,\dot{=}\,0{,}18 \)}{\( \frac{10}{36}\,\dot{=}\,0{,}28\)}{}{}}
\LANGsk{
\Question{
Hodíme dvoma kockami. Aká je pravdepodobnosť, že súčet padnutých bodov bude \( 5 \) alebo \( 6 \)? Výsledok zaokrúhlite na dve desatinné miesta.}
{\( \frac14= 0{,}25 \)}{\( \frac5{36}\,\dot{=}\,0{,}14 \)}{\( \frac2{11}\,\dot{=}\,0{,}18 \)}{\( \frac{10}{36}\,\dot{=}\,0{,}28\)}{}{}}
\LANGes{
\Question{
Vamos a tirar dos dados. ¿Qué es la probabilidad de que la suma de los puntos sea  \( 5 \) o \( 6 \)? Expresa el resultado con dos cifras decimales.}
{\( \frac14= 0{,}25 \)}{\( \frac5{36}\,\dot{=}\,0{,}14 \)}{\( \frac2{11}\,\dot{=}\,0{,}18 \)}{\( \frac{10}{36}\,\dot{=}\,0{,}28\)}{}{}}
\End

\Data\ID{1003029302}
\Updated{2021-02-03 02:02:51}
\Level{B}
\SubArea{Probability}
\LANGen{
\Question{
The quality inspection found that \( 85\% \) of the items are without defect, exactly one defect has \( 10\% \) of the items, and the other items have more than one defect. We pick one item randomly. What is the probability that this item has at least one defect?}
{\( 0{.}15 \)}{\( 0{.}10 \)}{\( 0{.}95 \)}{\( 0{.}01 \)}{}{}}
\LANGpl{
\Question{
Kontrola jakości wykazała, że \( 85\% \) pozycji nie ma wad, dokładnie jedną wadę ma \( 10\% \)  pozycji, a pozostałe mają więcej niż jedną wadę. Wybieramy jeden przedmiot losowo. Jakie jest prawdopodobieństwo, że przedmiot ten ma co najmniej jedną wadę?}
{\( 0{,}15 \)}{\( 0{,}10 \)}{\( 0{,}95 \)}{\( 0{,}01 \)}{}{}}
\LANGcs{
\Question{
Kontrola zjistila, že \( 85\:\% \) vyrobených součástek je bez vad, právě jednu vadu má \( 10\:\% \) součástek a ostatní součástky mají více než jednu vadu. Jaká je pravděpodobnost, že náhodně vybraná součástka bude mít alespoň jednu vadu?}
{\( 0{,}15 \)}{\( 0{,}10 \)}{\( 0{,}95 \)}{\( 0{,}01 \)}{}{}}
\LANGsk{
\Question{
Kontrola zistila, že z vyrobených súčiastok je \( 85\% \) bez chýb, práve jednu chybu má \( 10\% \) súčiastok a viac než jednu chybu majú ostatné súčiastky. Aká je pravdepodobnosť, že náhodne vybraný výrobok bude mať aspoň jednu chybu?}
{\( 0{,}15 \)}{\( 0{,}10 \)}{\( 0{,}95 \)}{\( 0{,}01 \)}{}{}}
\LANGes{
\Question{
Una verificación ha comprobado que \( 85\:\% \) de los productos es sin defecto,  \( 10\:\% \) de los productos tiene exactamente un defecto y el resto tiene más de un defecto. ¿Qué es la probabilidad de que un producto elegido al azar tenga por lo mínimo un defecto?}
{\( 0{,}15 \)}{\( 0{,}10 \)}{\( 0{,}95 \)}{\( 0{,}01 \)}{}{}}
\End

\Data\ID{1003029304}
\Updated{2021-02-03 03:02:58}
\Level{B}
\SubArea{Probability}
\LANGen{
\Question{
The probability of a man hitting a target is \( 0{.}9 \). What is the probability that he hits the target twice in a row? Round the result to two decimal places.}
{\( 0{.}81 \)}{\( 0{.}99 \)}{\( 0{.}95 \)}{\( 0{.}18 \)}{}{}}
\LANGpl{
\Question{
Prawdopodobieństwo trafienia w cel przez mężczyznę  wynosi  \( 0{,}9 \).  Jakie jest prawdopodobieństwo, że trafi w cel dwa razy z rzędu? Zaokrąglij wynik do dwóch miejsc po przecinku.}
{\( 0{,}81 \)}{\( 0{,}99 \)}{\( 0{,}95 \)}{\( 0{,}18 \)}{}{}}
\LANGcs{
\Question{
Pravděpodobnost, že střelec trefí terč je \( 0{,}9 \). Jaká je pravděpodobnost, že zasáhne cíl dvakrát za sebou? Výsledek zaokrouhlete na dvě desetinná místa.}
{\( 0{,}81 \)}{\( 0{,}99 \)}{\( 0{,}95 \)}{\( 0{,}18 \)}{}{}}
\LANGsk{
\Question{
Pravdepodobnosť, že strelec trafí terč je \( 0{,}9 \). Aká je pravdepodobnosť, že zasiahne cieľ dvakrát za sebou? Výsledok zaokrúhlite na dve desatinné miesta.}
{\( 0{,}81 \)}{\( 0{,}99 \)}{\( 0{,}95 \)}{\( 0{,}18 \)}{}{}}
\LANGes{
\Question{
La probabilidad de disparar al blanco con éxito es \( 0{,}9 \). ¿Qué es la probabilidad de disparar con éxito dos veces enseguida? Aproxima el resultado a dos cifras decimales.}
{\( 0{,}81 \)}{\( 0{,}99 \)}{\( 0{,}95 \)}{\( 0{,}18 \)}{}{}}
\End

\Data\ID{1003029305}
\Updated{2021-02-03 02:02:36}
\Level{B}
\SubArea{Probability}
\LANGen{
\Question{
The production process of a particular component consists of three independent phases.
Through long-term monitoring of production quality were found success rates of individual phases \( 90\% \), \( 80\% \) and \( 85\% \). If all three phases of the process are successful, the component is of good quality.  What is the probability of producing a good quality component?}
{\( 0{.}612 \)}{\( 0{.}003 \)}{\( 0{.}388 \)}{\( 0{.}997 \)}{}{}}
\LANGpl{
\Question{
Proces produkcji danego elementu składa się z trzech niezależnych faz. Poprzez długoterminowy monitoring jakości produkcji stwierdzono wskaźniki sukcesu poszczególnych faz \( 90\% \), \( 80\% \) i \( 85\% \). Jeśli wszystkie trzy fazy procesu zakończą się powodzeniem, element jest dobrej jakości. Jakie jest prawdopodobieństwo wytworzenia elementu dobrej jakości?}
{\( 0{,}612 \)}{\( 0{,}003 \)}{\( 0{,}388 \)}{\( 0{,}997 \)}{}{}}
\LANGcs{
\Question{
Výrobní proces určité součástky se skládá ze tří na sobě nezávislých operací. Dlouhodobým sledováním kvality výroby bylo zjištěno, že úspěšnost těchto operací je \( 90\:\% \), \( 80\:\% \) a \( 85\:\% \). Když se všechny tři operace vykonají úspěšně, je vyrobená součástka kvalitní. Jaká je pravděpodobnost výroby kvalitní součástky?}
{\( 0{,}612 \)}{\( 0{,}003 \)}{\( 0{,}388 \)}{\( 0{,}997 \)}{}{}}
\LANGsk{
\Question{
Výrobný proces určitej súčiastky sa skladá z troch na sebe nezávislých operácií. Dlhodobým sledovaním kvality výroby bolo zistené, že úspešnosť týchto operácií je \( 90\% \), \( 80\% \) a \( 85\% \). Ak sa všetky tri operácie vykonajú úspešne, je vyrobená súčiastka kvalitná. Aká je pravdepodobnosť výroby kvalitnej súčiastky?}
{\( 0{,}612 \)}{\( 0{,}003 \)}{\( 0{,}388 \)}{\( 0{,}997 \)}{}{}}
\LANGes{
\Question{
El proceso de construcción de un producto está compuesto por tres operaciones independientes. Sabemos que el éxito de estas operaciones es \( 90\:\% \), \( 80\:\% \) y \( 85\:\% \). 
Si todas las tres operaciones acaban con éxito, tenemos un producto de calidad. ¿Qué es la probabilidad de obtener un producto de calidad?}
{\( 0{,}612 \)}{\( 0{,}003 \)}{\( 0{,}388 \)}{\( 0{,}997 \)}{}{}}
\End

\Data\ID{1003041702}
\Updated{2021-02-04 01:02:01}
\Level{B}
\SubArea{Probability}
\LANGen{
\Question{
Product quality inspection was performed. Inspectors reported that \( 85\% \) of products have no defect, \( 8\% \) of products have exactly one defect, \( 5\% \) of products have exactly two defects and other products have more than two defects. What is the probability that a product selected at random has at least one defect?}
{\( 0{.}15 \)}{\( 0{.}07 \)}{\( 0{.}08 \)}{\( 0{.}13 \)}{}{}}
\LANGpl{
\Question{
Przeprowadzono kontrolę jakości produktu. Inspektorzy zgłosili, że \( 85\% \) produktów nie ma wad, \( 8\% \) produktów ma tylko jedną wadę, \( 5\% \) produktów ma dwie wady, a pozostałe produkty mają więcej niż dwie wady. Jakie jest prawdopodobieństwo, że wybrany losowo  produkt  ma co najmniej jedną wadę?}
{\( 0{,}15 \)}{\( 0{,}07 \)}{\( 0{,}08 \)}{\( 0{,}13 \)}{}{}}
\LANGcs{
\Question{
Kontrolou výrobků se zjistilo, že bez vad je \( 85\:\% \) z nich, nějakou jednu vadu má  \( 8\:\% \) z nich, nějaké dvě vady má \( 5\:\% \) z nich a ostatní výrobky mají více než dvě vady. Jaká je pravděpodobnost, že náhodně vybraný výrobek bude mít alespoň jednu vadu?}
{\( 0{,}15 \)}{\( 0{,}07 \)}{\( 0{,}08 \)}{\( 0{,}13 \)}{}{}}
\LANGsk{
\Question{
Kontrolou výrobkov sa zistilo, že bez chýb je \( 85\% \) z nich, nejakú jednu chybu má  \( 8\% \) z nich, nejaké dve chyby má \( 5\% \) z nich a ostatné výrobky majú viac ako dve chyby. Aká je pravdepodobnosť, že náhodne vybraný výrobok bude mať aspoň jednu chybu?}
{\( 0{,}15 \)}{\( 0{,}07 \)}{\( 0{,}08 \)}{\( 0{,}13 \)}{}{}}
\LANGes{
\Question{
Un control de calidad ha comprobado que \( 85\:\% \) de productos es sin defecto, \( 8\:\% \) de productos tiene un defecto, \( 5\:\% \) tiene dos defectos y el resto tiene más de dos defectos. ¿Qué es la probabilidad de que un producto elegido al azar tenga más de un defecto?}
{\( 0{,}15 \)}{\( 0{,}07 \)}{\( 0{,}08 \)}{\( 0{,}13 \)}{}{}}
\End

\Data\ID{1003041704}
\Updated{2021-02-17 03:02:18}
\Level{B}
\SubArea{Probability}
\LANGen{
\Question{
A set of Christmas tree lights contains \( 12 \) identical bulbs connected in parallel. Each of the bulbs has the reliability of \( 98\% \). What is the probability (expressed as a percentage) that all of the bulbs will glow? Round the result to one decimal place.

(Note: Reliability is the probability that the system will perform its intended function.)}
{\( 78{.}5\% \)}{\( 98{.}0\% \)}{\( 78{.}4\% \)}{\( 97{.}5\% \)}{}{}}
\LANGpl{
\Question{
Zestaw lampek choinkowych zawiera \( 12 \) identycznych żarówek połączonych równolegle. Niezawodność  każdej żarówki wynosi \( 98\% \). Jakie jest prawdopodobieństwo (wyrażone w procentach), że wszystkie żarówki będą świecić? Zaokrąglij wynik do jednego miejsca po przecinku. (Uwaga: niezawodność to prawdopodobieństwo, że system spełni zamierzoną funkcję.)}
{\( 78{,}5\% \)}{\( 98{,}0\% \)}{\( 78{,}4\% \)}{\( 97{,}5\% \)}{}{}}
\LANGcs{
\Question{
Vánoční osvětlení se skládá z \( 12 \) paralelně zapojených žárovek. Každá žárovka má spolehlivost \( 98\:\% \). Jaká je pravděpodobnost, že po připojení ke zdroji budou všechny žárovky svítit? Výsledek vyjádřete v procentech zaokrouhlených s přesností na desetiny. (Poznámka: Spolehlivost je pravděpodobnost, s jakou žárovka bude plnit svou funkci.)}
{\( 78{,}5\:\% \)}{\( 98{,}0\:\% \)}{\( 78{,}4\:\% \)}{\( 97{,}5\:\% \)}{}{}}
\LANGsk{
\Question{
Systém vianočných sviečok sa skladá z \( 12 \) žiaroviek zapojených paralelne. Každá žiarovka má spoľahlivosť \( 98\% \). Aká je pravdepodobnosť, že budú všetky vianočné sviečky svietiť? Výsledok vyjadrite v percentách zaokrúhlených s presnosťou na desatiny. (Poznámka: Spoľahlivosť je pravdepodobnosť, s akou žiarovka bude plniť svoju funkciu.)}
{\( 78{,}5\% \)}{\( 98{,}0\% \)}{\( 78{,}4\% \)}{\( 97{,}5\% \)}{}{}}
\LANGes{
\Question{
Una decoración de Navidad está compuesta por \( 12 \) bomlbillas conectadas paralelamente. Cada bombilla tiene eficacia de \( 98\:\% \). ¿Qué probabilidad hay de que después de conectarla a la fuente van a encender todas las bombillas? Expresa el resultado en porcentajes aproximados a una cifra decimal.}
{\( 78{,}5\:\% \)}{\( 98{,}0\:\% \)}{\( 78{,}4\:\% \)}{\( 97{,}5\:\% \)}{}{}}
\End

\Data\ID{1003041705}
\Updated{2021-02-04 03:02:56}
\Level{B}
\SubArea{Probability}
\LANGen{
\Question{
Product quality inspector checks two independent quality indicators, A and B. If the product does not comply with any of the indicators it is discarded otherwise it is accepted. The inspector reported that \( 95{.}4\% \) of the products was accepted. He also reported that \( 97{.}1\% \) of the products complied with the indicator A. How many of the products complied with the indicator B? Express the result in percentage and round it to two decimal places.}
{\( 98{.}25\% \)}{\( 98{.}24\% \)}{\( 92{.}63\% \)}{\( 92{.}64\% \)}{}{}}
\LANGpl{
\Question{
Inspektor jakości produktów sprawdza dwa niezależne wskaźniki jakości, A i B. Jeśli produkt nie spełnia któregokolwiek ze wskaźników, jest odrzucany, pozostałe są  akceptowane. Inspektor poinformował, że \( 95{,}4\% \) produktów zostało zaakceptowanych. Podał również, że  \( 97{,}1\% \) produktów spełniało wskaźnik A. Ile produktów spełniało wskaźnik B? Przedstaw wynik w procentach i zaokrąglij go do dwóch miejsc po przecinku.}
{\( 98{,}25\% \)}{\( 98{,}24\% \)}{\( 92{,}63\% \)}{\( 92{,}64\% \)}{}{}}
\LANGcs{
\Question{
Výstupní kontrola sleduje dva nezávislé ukazatela kvality součástek: A a B. Pokud součástka nevyhoví kterémukoliv ukazateli kvality, je na výstupní kontrole vyřazena (hodnocena jako nekvalitní). Výstupní kontrola vyhodnotila jako kvalitní 95,4\:\% součástek, přičemž ukazateli A vyhovělo 97,1\:\% součástek. Kolik součástek vyhovělo ukazateli B? Výsledek vyjádřete s přesností na setiny procent.}
{\( 98{,}25\:\% \)}{\( 98{,}24\:\% \)}{\( 92{,}63\:\% \)}{\( 92{,}64\:\% \)}{}{}}
\LANGsk{
\Question{
Po vyrobení súčiastok sleduje výstupná kontrola ich dva nezávislé ukazovatele kvality: A a B. Ak súčiastka nespĺňala nejaký ukazovateľ, bola vyradená. Kontrola schválila \( 95{,}4\% \) súčiastok, pričom ukazovateľ A spĺňalo \( 97{,}1\% \) súčiastok. Koľko súčiastok spĺňalo ukazovateľ B? Výsledok vyjadrite v percentách zaokrúhlených s presnosťou na stotiny.}
{\( 98{,}25\% \)}{\( 98{,}24\% \)}{\( 92{,}63\% \)}{\( 92{,}64\% \)}{}{}}
\LANGes{
\Question{
El control de salida tiene dos punteros: A y B. Si un producto no satisface cuálquier de los punteros, se descarta (no es de calidad). El control de salida ha evaluado  95,4\:\ % de los productos como de calidad y al puntero A lo ha satisfecho 97,1\:\ % de los productos. ¿Cuántos productos han satisfecho el puntero B? Aproxima el resultado a dos cifras decimales.}
{\( 98{,}25\:\% \)}{\( 98{,}24\:\% \)}{\( 92{,}63\:\% \)}{\( 92{,}64\:\% \)}{}{}}
\End

\Data\ID{1003041706}
\Updated{2021-02-04 03:02:46}
\Level{B}
\SubArea{Probability}
\LANGen{
\Question{
Four participants of a target shooting competition do hit the target with hitting probabilities: \( 0{.}80 \); \( 0{.}85 \); \( 0{.}90 \) and \( 0{.}95 \). What is the probability that just one of the participants does hit the target? Round the result to four decimal places.}
{\( 0{.}0057 \)}{\( 0{.}0056 \)}{\( 0{.}9999 \)}{\( 0{.}9998 \)}{}{}}
\LANGpl{
\Question{
Cztery osoby biorące udział w zawodach strzeleckich strzelają  w cel, prawdopodobieństwo trafienia wynosi: \( 0{,}80 \); \( 0{,}85 \); \( 0{,}90 \) i \( 0{,}95 \). Jakie jest prawdopodobieństwo, że tylko jeden z uczestników trafi w cel? Zaokrąglij wynik do czterech miejsc po przecinku.}
{\( 0{,}0057 \)}{\( 0{,}0056 \)}{\( 0{,}9999 \)}{\( 0{,}9998 \)}{}{}}
\LANGcs{
\Question{
Čtyři střelci střílejí na cíl. Trefí se s pravděpodobnostmi: \( 0{,}80 \); \( 0{,}85 \); \( 0{,}90 \) a \( 0{,}95 \). Jaká je pravděpodobnost, že právě jeden z těchto střelců zasáhne cíl? Výsledek zaokrouhlete s přesností na čtyři desetinná místa.}
{\( 0{,}0057 \)}{\( 0{,}0056 \)}{\( 0{,}9999 \)}{\( 0{,}9998 \)}{}{}}
\LANGsk{
\Question{
Štyria strelci strieľajú na cieľ a zasahujú ho s pravdepodobnosťami: \( 0{,}80 \); \( 0{,}85 \); \( 0{,}90 \) a \( 0{,}95 \). Aká je pravdepodobnosť, že práve jeden z nich zasiahne cieľ? Výsledok zaokrúhlite s presnosťou na štyri desatinné miesta.}
{\( 0{,}0057 \)}{\( 0{,}0056 \)}{\( 0{,}9999 \)}{\( 0{,}9998 \)}{}{}}
\LANGes{
\Question{
Cuatro personas disparan al blanco. Las probabilidades de disparar con écito son \( 0{,}80 \); \( 0{,}85 \); \( 0{,}90 \) y \( 0{,}95 \). ¿Qué es la probabilidad de que exactamente una persona dispara con éxito? Aproxima el resultado a cuatro cifras decimales.}
{\( 0{,}0057 \)}{\( 0{,}0056 \)}{\( 0{,}9999 \)}{\( 0{,}9998 \)}{}{}}
\End

\Data\ID{1003041707}
\Updated{2021-02-04 03:02:26}
\Level{B}
\SubArea{Probability}
\LANGen{
\Question{
Four participants of a target shooting competition do hit the target with hitting probabilities: \( 0{.}80 \); \( 0{.}85\); \( 0{.}90 \) and \( 0{.}95 \). What is the probability that at least one of the participants does hit the target? Round the result to four decimal places.}
{\( 0{.}9999 \)}{\( 0{.}9998 \)}{\( 0{.}0057 \)}{\( 0{.}0056 \)}{}{}}
\LANGpl{
\Question{
Cztery osoby biorące udział w zawodach strzeleckich strzelają  w cel,  prawdopodobieństwo  trafienia wynosi: \( 0{,}80 \); \( 0{,}85\); \( 0{,}90 \) i \( 0{,}95 \). Jakie jest prawdopodobieństwo, że co najmniej jeden z uczestników trafi w cel? Zaokrąglij wynik do czterech miejsc po przecinku.}
{\( 0{,}9999 \)}{\( 0{,}9998 \)}{\( 0{,}0057 \)}{\( 0{,}0056 \)}{}{}}
\LANGcs{
\Question{
Čtyři střelci střílejí na cíl. Trefí se s pravděpodobnostmi: \( 0{,}80 \); \( 0{,}85\); \( 0{,}90 \) a \( 0{,}95 \). Jaká je pravděpodobnost, že alespoň jeden z těchto střelců zasáhne cíl? Výsledek zaokrouhlete s přesností na čtyři desetinná místa.}
{\( 0{,}9999 \)}{\( 0{,}9998 \)}{\( 0{,}0057 \)}{\( 0{,}0056 \)}{}{}}
\LANGsk{
\Question{
Štyria strelci strieľajú na cieľ a zasahujú ho s pravdepodobnosťami: \( 0{,}80 \); \( 0{,}85\); \( 0{,}90 \) a \( 0{,}95 \). Aká je pravdepodobnosť, že aspoň jeden z nich zasiahne cieľ? Výsledok zaokrúhlite s presnosťou na štyri desatinné miesta.}
{\( 0{,}9999 \)}{\( 0{,}9998 \)}{\( 0{,}0057 \)}{\( 0{,}0056 \)}{}{}}
\LANGes{
\Question{
Cuatro personas disparan al blanco. Sus probabilidades de disparar con éxito son: \( 0{,}80 \); \( 0{,}85\); \( 0{,}90 \) y \( 0{,}95 \). ¿Qué es la probablidad de que por lo mínimo una de las personas dispara con éxito? Aproxima el resultado a cuatro cifras decimales.}
{\( 0{,}9999 \)}{\( 0{,}9998 \)}{\( 0{,}0057 \)}{\( 0{,}0056 \)}{}{}}
\End

\Data\ID{1003164501}
\Updated{2021-02-16 02:02:07}
\Level{B}
\SubArea{Probability}
\LANGen{
\Question{
In the house with a \( 7 \) metre high ground storey and \( 6 \) other storeys (\( 5 \) metres high each), there is a lift. On each storey, it is possible to enter this lift through a glass door that is \( 2 \) metres high. The lift malfunctioned and stopped somewhere on its way. What is the probability that (at the moment of stoppage) it would not be possible to see out of the lift only the wall of the lift shaft?}
{\( 0.7500 \)}{\( 0.7838 \)}{\( 0.7188 \)}{\( 0.7647 \)}{\( 0.7353 \)}{\( 0.7568 \)}}
\LANGpl{
\Question{
Dom posiada \( 7 \) metrowy parter oraz \( 6 \) pięter (każde o wysokości \( 5 \) metrów). W budynku znajduje się winda. Na każdym piętrze, można wejść do windy szklanymi drzwiami o wysokości  \( 2 \) metrów.  Winda uległa awarii i zatrzymała się. Jakie jest prawdopodobieństwo, że  (w momencie zatrzymania) nie jest możliwe zobaczenie z windy tylko ściany szybu windy?}
{\( 0{,}7500 \)}{\( 0{,}7838 \)}{\( 0{,}7188 \)}{\( 0{,}7647 \)}{\( 0{,}7353 \)}{\( 0{,}7568 \)}}
\LANGcs{
\Question{
V domě, který má  \( 7 \) metrů vysoké přízemí a  \( 6 \) poschodí (každé má výšku \( 5 \) metrů), je výtah. V každém poschodí i v přízemí se do něj vchází skleněnými dveřmi, které jsou vysoké \( 2 \) metry. Během poruchy zůstal výtah někde náhodně stát. Jaká je pravděpodobnost, že z něj (v okamžiku zastavení) nebude vidět jen stěnu šachty? Výsledek zaokrouhlete na \( 4 \) desetinná místa.}
{\( 0{,}7500 \)}{\( 0{,}7838 \)}{\( 0{,}7188 \)}{\( 0{,}7647 \)}{\( 0{,}7353 \)}{\( 0{,}7568 \)}}
\LANGsk{
\Question{
V dome, ktorý má  \( 7 \) metrov vysoké prízemie a  \( 6 \) poschodí (každé má výšku \( 5 \) metrov), je výťah. Na každom poschodí aj v prízemí sa do neho vchádza cez sklenené dvere, ktoré sú vysoké \( 2 \) metre. Pre poruchu zostal výťah niekde náhodne stáť. Aká je pravdepodobnosť, že z neho (v okamihu zastavenia) nebude vidieť len stenu šachty? Výsledok zaokrúhlite na \( 4 \) desatinné miesta.}
{\( 0{,}7500 \)}{\( 0{,}7838 \)}{\( 0{,}7188 \)}{\( 0{,}7647 \)}{\( 0{,}7353 \)}{\( 0{,}7568 \)}}
\LANGes{
\Question{
Una casa tiene planta baja de  \( 7 \) metros de altura y  \( 6 \) plantas (cada una tiene altura de \( 5 \) metros). En esta casa hay un ascensor al cuál se entra por puerta de vidrio de altura de dos metros en cada planta. Durante un trastorno el ascensor para al azar. ¿Qué probabilidad hay de que en ese momento no vamos a poder ver  solamente la pared? Aproxima el resultado a \( 4 \) cifras decimales.}
{\( 0{,}7500 \)}{\( 0{,}7838 \)}{\( 0{,}7188 \)}{\( 0{,}7647 \)}{\( 0{,}7353 \)}{\( 0{,}7568 \)}}
\End

\Data\ID{1003164502}
\Updated{2021-02-03 03:02:13}
\Level{B}
\SubArea{Probability}
\LANGen{
\Question{
Suppose the points \( A \) and \( B \) are randomly placed on a circle with a radius \( r \). What is the probability that the distance between \( A \) and \( B \) (length of the chord \( AB \)) is at least \( r \)?}
{\( \frac23 \)}{\( \frac13 \)}{\( \frac16 \)}{\( \frac56 \)}{\( \frac12 \)}{}}
\LANGpl{
\Question{
Punkty \( A \) i \( B \) zostały losowo umieszczone na okręgu o promieniu \( r \). Jakie jest prawdopodobieństwo, że odległość pomiędzy  \( A \) i \( B \) (długość cięciwy \( AB \)) jest równa co najmniej promieniowi  \( r \)?}
{\( \frac23 \)}{\( \frac13 \)}{\( \frac16 \)}{\( \frac56 \)}{\( \frac12 \)}{}}
\LANGcs{
\Question{
Máme body  \( A \) a \( B \) náhodně umístěné na kružnici s poloměrem \( r \). Jaká je pravděpodobnost, že vzdálenost bodů \( A \) a \( B \) (délka tětivy \( AB \)) bude alespoň \( r \)?}
{\( \frac23 \)}{\( \frac13 \)}{\( \frac16 \)}{\( \frac56 \)}{\( \frac12 \)}{}}
\LANGsk{
\Question{
Majme body  \( A \) a \( B \) náhodne umiestnené na kružnici s polomerom \( r \). Aká je pravdepodobnosť, že vzdialenosť bodov \( A \) a \( B \) (dĺžka tetivy \( AB \)) bude aspoň \( r \)?}
{\( \frac23 \)}{\( \frac13 \)}{\( \frac16 \)}{\( \frac56 \)}{\( \frac12 \)}{}}
\LANGes{
\Question{
Tenemos puntos \( A \) y \( B \) situados al azar en una circunferecia con radio \( r \). ¿Qué es la probabilidad de que la distancia entre los puntos \( A \) y \( B \) (la longitud de la cuerda \( AB \)) sea por lo mínimo \( r \)?}
{\( \frac23 \)}{\( \frac13 \)}{\( \frac16 \)}{\( \frac56 \)}{\( \frac12 \)}{}}
\End

\Data\ID{1103019203}
\Updated{2021-02-16 02:02:59}
\Level{B}
\SubArea{Probability}
\IMG{1103019203_0.pdf}
\LANGen{
\Question{
The dartboard in the picture is made up of three concentric circles with radii \( 4 \), \( 6 \) and \( 8 \) cm. Assuming that a dart thrown will land randomly on the dartboard, what is the probability that it lands in the red region?}
{\( \frac5{16}\,\dot{=}\,0{.}3125 \)}{\( \frac3{4}=0{.}75 \)}{\( \frac1{4}=0{.}25 \)}{\( \frac9{16}\,\dot{=}\,0{.}5625 \)}{}{}}
\LANGpl{
\Question{
Tarcza do gry w rzutki składa się z trzech okręgów współśrodkowych o promieniach  \( 4 \), \( 6 \) i \( 8 \) cm. Zakładając, że rzutka losowo wybierze obszar, jakie jest prawdopodobieństwo wylądowania rzutki na obszarze  czerwonym?}
{\( \frac5{16}\,\dot{=}\,0{,}3125 \)}{\( \frac3{4}=0{,}75 \)}{\( \frac1{4}=0{,}25 \)}{\( \frac9{16}\,\dot{=}\,0{,}5625 \)}{}{}}
\LANGcs{
\Question{
V terči na obrázku je poloměr tmavomodré kružnice \( 8 \) cm, poloměr tmavočervené kružnice \( 6 \) cm a poloměr žlutého kruhu \( 4 \) cm. Náhodně vystřelený šíp trefil terč. S jakou pravděpodobností zasáhl červené mezikruží?}
{\( \frac5{16}\,\dot{=}\,0{,}3125 \)}{\( \frac3{4}=0{,}75 \)}{\( \frac1{4}=0{,}25 \)}{\( \frac9{16}\,\dot{=}\,0{,}5625 \)}{}{}}
\LANGsk{
\Question{
V terči na obrázku je polomer tmavomodrej kružnice \( 8 \) cm, polomer tmavočervenej kružnice \( 6 \) cm a polomer žltého kruhu \( 4 \) cm. Náhodne vystrelený šíp trafil terč. S akou pravdepodobnosťou zasiahol červené medzikružie?}
{\( \frac5{16}\,\dot{=}\,0{,}3125 \)}{\( \frac3{4}=0{,}75 \)}{\( \frac1{4}=0{,}25 \)}{\( \frac9{16}\,\dot{=}\,0{,}5625 \)}{}{}}
\LANGes{
\Question{
El radio del círculo azul en la diana de dardos del dibujo es \( 8 \) cm, el radio del círculo rojo es \( 6 \) cm y el radio del círculo amarillo es \( 4 \) cm. Un dardo tirado al azar ha acertado a la diana. ¿Qué probabilidad hay de que está en el anillo rojo?}
{\( \frac5{16}\,\dot{=}\,0{,}3125 \)}{\( \frac3{4}=0{,}75 \)}{\( \frac1{4}=0{,}25 \)}{\( \frac9{16}\,\dot{=}\,0{,}5625 \)}{}{}}
\End

\Data\ID{1103041703}
\Updated{2021-02-17 03:02:08}
\Level{B}
\SubArea{Probability}
\IMG{1103041703.pdf}
\LANGen{
\Question{
Three identical light bulbs are connected to the battery as shown in the electrical circuit diagram.  The reliability of each of the bulbs is \( 0{.}95 \). What is the probability that the current flows through the circuit? Round the result to \( 4 \) decimal places.
(Note: Reliability is the probability that the component will perform its intended function.)}
{\( 0{.}9951 \)}{\( 0{.}8574 \)}{\( 0{.}9476 \)}{\( 0{.}9500 \)}{}{}}
\LANGpl{
\Question{
Trzy identyczne żarówki są podłączone do akumulatora, jak pokazano na schemacie. Niezawodność każdej z żarówek wynosi \( 0{,}95 \). Jakie jest prawdopodobieństwo, że prąd przepływa przez obwód? Zaokrąglij wynik do  \( 4 \) miejsc  po przecinku. (Uwaga: niezawodność to prawdopodobieństwo, że komponent spełni zamierzoną funkcję.)}
{\( 0{,}9951 \)}{\( 0{,}8574 \)}{\( 0{,}9476 \)}{\( 0{,}9500 \)}{}{}}
\LANGcs{
\Question{
Žárovky jsou připojeny ke zdroji napětí podle schématu na obrázku. Spolehlivost každé žárovky je \( 0{,}95 \). Jaká je pravděpodobnost, že obvodem bude procházet proud? Výsledek zaokrouhlete na \( 4 \) desetinné místa. (Poznámka: Spolehlivost je pravděpodobnost, s jakou žárovka bude plnit svou funkci.)}
{\( 0{,}9951 \)}{\( 0{,}8574 \)}{\( 0{,}9476 \)}{\( 0{,}9500 \)}{}{}}
\LANGsk{
\Question{
Žiarovky sú zapojené do siete podľa obrázka. Spoľahlivosť každej žiarovky je \( 0{,}95 \). Aká je pravdepodobnosť, že prúd prejde sieťou? Výsledok zaokrúhlite na \( 4 \) desatinné miesta. (Poznámka: Spoľahlivosť je pravdepodobnosť, s akou žiarovka bude plniť svoju funkciu.)}
{\( 0{,}9951 \)}{\( 0{,}8574 \)}{\( 0{,}9476 \)}{\( 0{,}9500 \)}{}{}}
\LANGes{
\Question{
Las bombillas están conectadas con la fuente de voltaje según el esquema en el dibujo. La eficacia de cada bombilla es \( 0{,}95 \). ¿Qué probabilidad hay de que la corriente va a pasar por el circuito? Aproxima el resultado a  \( 4 \) cifras decimales.}
{\( 0{,}9951 \)}{\( 0{,}8574 \)}{\( 0{,}9476 \)}{\( 0{,}9500 \)}{}{}}
\End

\Data\ID{1103164503}
\Updated{2021-02-16 11:02:17}
\Level{B}
\SubArea{Probability}
\IMG{1103164503.pdf}
\LANGen{
\Question{
An equilateral triangle with a side of 3 metres is drawn on the wall. Inside the triangle, there is a circle with the diameter of 1 metre.  If a fly sits by chance in the triangle, what is the probability that it does not sit inside the circle? Round the result to 4 decimal places.}
{\( 0.7985 \)}{\( 0.2015 \)}{\( 0.8061 \)}{\( 0.1939 \)}{}{}}
\LANGpl{
\Question{
Trójkąt równoboczny o boku równym 3 metry jest narysowany na ścianie.  W trójkącie znajduje się okrąg o średnicy 1 metra. Jeśli mucha usiądzie na trójkącie jakie jest prawdopodobieństwo, że nie usiądzie wewnątrz okręgu?  Zaokrągli wynik do 4 miejsc po przecinku.}
{\( 0{,}7985 \)}{\( 0{,}2015 \)}{\( 0{,}8061 \)}{\( 0{,}1939 \)}{}{}}
\LANGcs{
\Question{
Na stěně je namalovaný rovnostranný trojúhelník se stranou délky \( 3 \) metry. Uvnitř trojúhelníku je kruh s průměrem \( 1 \) metr. Jaká je pravděpodobnost, že moucha, která si sedne na náhodně vybrané místo uvnitř trojúhelníku, nebude sedět  uvnitř kruhu? Výsledek zaokrouhlete na \( 4 \) desetinná místa.}
{\( 0{,}7985 \)}{\( 0{,}2015 \)}{\( 0{,}8061 \)}{\( 0{,}1939 \)}{}{}}
\LANGsk{
\Question{
Na stene je namaľovaný rovnostranný trojuholník, ktorého strana má veľkosť \( 3 \) metre. Vo vnútri trojuholníka je kruh s priemerom \( 1 \) meter. Aká je pravdepodobnosť, že ak si mucha náhodne sadne do trojuholníka, nebude sedieť  vo vnútri kruhu? Výsledok zaokrúhlite na \( 4 \) desatinné miesta.}
{\( 0{,}7985 \)}{\( 0{,}2015 \)}{\( 0{,}8061 \)}{\( 0{,}1939 \)}{}{}}
\LANGes{
\Question{
Un triángulo equilátero con el lado de \( 3 \) metros está dibujado en la pared. Dentro del triángulo está un círculo con radio de \( 1 \) metro. Una mosca baja a la pared dentro del triángulo. ¿Qué probabilidad hay de que no está dentro del triángulo? Aproxima el resultado a  \( 4 \) cifras decimales.}
{\( 0{,}7985 \)}{\( 0{,}2015 \)}{\( 0{,}8061 \)}{\( 0{,}1939 \)}{}{}}
\End

\Data\ID{1103164504}
\Updated{2021-02-16 02:02:06}
\Level{B}
\SubArea{Probability}
\IMG{1103164504.pdf}
\LANGen{
\Question{
An equilateral triangle is drawn on the wall. Inside the triangle, there is a circle with the radius of \( 1 \) metre inscribed. If a fly sits by chance in the triangle, what is the probability that it does not sit  inside the circle? Round the result to \( 4 \) decimal places.}
{\( 0.3954 \)}{\( 0.6046 \)}{\( 0.3023 \)}{\( 0.6977 \)}{}{}}
\LANGpl{
\Question{
Trójkąt równoboczny jest narysowany na ścianie. Okrąg o promieniu  \( 1 \) metra jest wpisany w trójkąt. Jeśli mucha usiądzie przypadkowo na trójkącie, jakie jest prawdopodobieństwo, że nie znajduje się w okręgu? Zaokrągli wynik do  \( 4 \) miejsc po przecinku.}
{\( 0{,}3954 \)}{\( 0{,}6046 \)}{\( 0{,}3023 \)}{\( 0{,}6977 \)}{}{}}
\LANGcs{
\Question{
Na stěně je namalovaný rovnostranný trojúhelník, do kterého je vepsaná kružnice s poloměrem \( 1 \)  metr. Jaká je pravděpodobnost, že moucha, která si sedne na náhodně vybrané místo uvnitř trojúhelníku, nebude sedět  uvnitř kružnice? Výsledek zaokrouhlete na \( 4 \) desetinná místa.}
{\( 0{,}3954 \)}{\( 0{,}6046 \)}{\( 0{,}3023 \)}{\( 0{,}6977 \)}{}{}}
\LANGsk{
\Question{
Na stene je namaľovaný rovnostranný trojuholník, do ktorého je vpísaná kružnica s polomerom \( 1 \)  meter. Aká je pravdepodobnosť, že ak si mucha náhodne sadne do trojuholníka, nebude sedieť  vo vnútri kružnice? Výsledok zaokrúhlite na \( 4 \) desatinné miesta.}
{\( 0{,}3954 \)}{\( 0{,}6046 \)}{\( 0{,}3023 \)}{\( 0{,}6977 \)}{}{}}
\LANGes{
\Question{
Un triángulo equilátero con una circunferencia inscrita de radio \( 1 \) metro está dibujada en la pared. Una mosca baja en un lugar del triángulo al azar. ¿Qué probabilidad hay de que no estará dentro de la circunferencia?  Aproxima el resultado a \( 4 \) cifras decimales.}
{\( 0{,}3954 \)}{\( 0{,}6046 \)}{\( 0{,}3023 \)}{\( 0{,}6977 \)}{}{}}
\End

\Data\ID{1103164505}
\Updated{2021-02-17 03:02:45}
\Level{B}
\SubArea{Probability}
\IMG{1103164505.pdf}
\LANGen{
\Question{
Suppose we have a rectangular fish tank which is \( 4\,\mathrm{dm} \) long, \( 2\,\mathrm{dm} \) wide, and it is filled with water up to the height of \( 3\,\mathrm{dm} \). In its four bottom corners, there are jets through which a fresh air is driven into the water in specific intervals. The fresh air is driven to the distance of up to \( 5\,\mathrm{cm} \) from the tank corners.  If a fish swims inside the tank, what is the probability that the fish will not be hit by the stream of bubbles, at the moment when all four jets are acting? The fish dimensions can be neglected, round the result to \( 4 \) decimal places.}
{\( 0.9891 \)}{\( 0.0109 \)}{\( 0.9984 \)}{\( 0.0016 \)}{\( 0.9782 \)}{\( 0.0218 \)}}
\LANGpl{
\Question{
Załóżmy, że posiadamy akwarium w kształcie prostokąta o długości \( 4\,\mathrm{dm} \)  i szerokości \( 2\,\mathrm{dm} \), akwarium jest wypełnione wodą do wysokości \( 3\,\mathrm{dm} \). W czterech dolnych narożnikach znajdują się dysze, przez które świeże powietrze jest wpuszczane do wody w określonych odstępach czasu. Świeże powietrze jest kierowane na odległość do  \( 5\,\mathrm{cm} \) od narożników.  Jeśli ryba pływa w akwarium, jakie jest prawdopodobieństwo, że ryba nie zostanie uderzona strumieniem bąbelków, w momencie działania wszystkich czterech dysz? Wymiary ryb można pominąć. Zaokrągli wynik do \( 4 \) miejsc po przecinku.}
{\( 0{,}9891 \)}{\( 0{,}0109 \)}{\( 0{,}9984 \)}{\( 0{,}0016 \)}{\( 0{,}9782 \)}{\( 0{,}0218 \)}}
\LANGcs{
\Question{
Do akvária, které má tvar kvádru s rozměry podstavy \( 4\,\mathrm{dm} \) x \( 2\,\mathrm{dm} \), je napuštěná voda do výšky \( 3\,\mathrm{dm} \). Ve všech dolních rozích akvária (viz obrázek) jsou trysky, kterými je do vody vháněn v určitých intervalech čerstvý vzduch, a to  až do vzdálenosti \( 5\,\mathrm{cm} \) od rohů akvária. Jaká je pravděpodobnost, že ve chvíli, kdy trysky začnou do akvária vhánět vzduch, nebude rybka (jejíž rozměry zanedbáváme) proudem vzduchu zasažena? Výsledek zaokrouhlete na \( 4 \) desetinná místa.}
{\( 0{,}9891 \)}{\( 0{,}0109 \)}{\( 0{,}9984 \)}{\( 0{,}0016 \)}{\( 0{,}9782 \)}{\( 0{,}0218 \)}}
\LANGsk{
\Question{
Do akvária, ktoré má tvar kvádra s rozmermi podstavy \( 4\,\mathrm{dm} \) a \( 2\,\mathrm{dm} \), je naliata voda do výšky \( 3\,\mathrm{dm} \). V jeho štyroch spodných rohoch sú trysky, cez ktoré je do vody vháňaný v určitých intervaloch čerstvý vzduch, a to  až do vzdialenosti \( 5\,\mathrm{cm} \) od rohov akvária. Aká je pravdepodobnosť, že plávajúca rybka (ktorej rozmery zanedbávame) nebude zasiahnutá prúdom bublín v momente začiatku činnosti všetkých štyroch trysiek? Výsledok zaokrúhlite na \( 4 \) desatinné miesta.}
{\( 0{,}9891 \)}{\( 0{,}0109 \)}{\( 0{,}9984 \)}{\( 0{,}0016 \)}{\( 0{,}9782 \)}{\( 0{,}0218 \)}}
\LANGes{
\Question{
Un acuario tiene forma de un prisma con medidas de la base \( 4\,\mathrm{dm} \) x \( 2\,\mathrm{dm} \) y el agua denro tiene altura \( 3\,\mathrm{dm} \). En todas las esquinas de la base hay boquillas por cuáles entra aire el agua en el acuario hasta la distancia de \( 5\,\mathrm{cm} \) de las esquinas. (vea el dibujo) ¿Qué probabilidad hay de que en el momento cuando empieza a entrar el aire, no pega al pez (cuyo tamaño no nos interesa)? Aproxima el resultado a  \( 4 \)  cifras decimales.}
{\( 0{,}9891 \)}{\( 0{,}0109 \)}{\( 0{,}9984 \)}{\( 0{,}0016 \)}{\( 0{,}9782 \)}{\( 0{,}0218 \)}}
\End

\Data\ID{1103164506}
\Updated{2021-02-17 02:02:23}
\Level{B}
\SubArea{Probability}
\IMG{1103164506.pdf}
\LANGen{
\Question{
At night, a parachutist landed on the spot \( M \), which is \( 3\,\mathrm{km} \) and \( 4\,\mathrm{km} \) away from the two straight and mutually perpendicular roads \( p \) and \( q \) respectively (see the picture).   From the landing point, the parachutist walks straight in a random direction at the constant speed of \( 6\,\mathrm{km}/\mathrm{h} \). What is the probability that he reaches one of the roads in less than an hour? Round the result to \( 4 \) decimal places.
\[ \]
Hint: In the case of linear motion with constant speed, the speed is equal to the ratio of the displacement and the time of motion.}
{\( 0.5505 \)}{\( 0.4495 \)}{\( 0.6011 \)}{\( 0.3989 \)}{\( 0.3511 \)}{\( 0.6489 \)}}
\LANGpl{
\Question{
W nocy spadochroniarz wylądował na miejscu \( M \), oddalonym o \( 3\,\mathrm{km} \) i \( 4\,\mathrm{km} \) od dwóch prostych i wzajemnie prostopadłych dróg, odpowiednio, \( p \) i \( q \) (patrz zdjęcie). Z punktu lądowania spadochroniarz porusza się prosto w przypadkowym kierunku ze stałą prędkością \( 6\,\mathrm{km}/\mathrm{h} \).  Jakie jest prawdopodobieństwo, że dotrze do jednej z dróg w mniej niż godzinę? Zaokrąglij wynik do  \( 4 \) miejsc po przecinku.
\[ \]
Wskazówka: W przypadku ruchu liniowego ze stałą prędkością, prędkość jest równa stosunkowi przemieszczenia do czasu ruchu.}
{\( 0{,}5505 \)}{\( 0{,}4495 \)}{\( 0{,}6011 \)}{\( 0{,}3989 \)}{\( 0{,}3511 \)}{\( 0{,}6489 \)}}
\LANGcs{
\Question{
Výsadkář dopadl v noci na místo \( M \), které je od dvou na sebe kolmých cest \( p \) a \( q \) vzdálené \( 3\,\mathrm{km} \) a \( 4\,\mathrm{km} \) (viz obrázek). Z místa dopadu se výsadkář vydal rovnoměrným přímočarým pohybem rychlostí \( 6\,\mathrm{km}/\mathrm{h} \) náhodně zvoleným směrem. Jaká je pravděpodobnost, že výsadkář do hodiny narazí na některou z cest? Výsledek zaokrouhlete na \( 4 \) desetinná místa.
\[ \]
Tip: V případě rovnoměrného přímočarého pohybu se rychlost rovná poměru dráhy a času.}
{\( 0{,}5505 \)}{\( 0{,}4495 \)}{\( 0{,}6011 \)}{\( 0{,}3989 \)}{\( 0{,}3511 \)}{\( 0{,}6489 \)}}
\LANGsk{
\Question{
Výsadkár dopadol v noci na miesto \( M \), ktoré je od dvoch priamych a na seba kolmých ciest \( p \) a \( q \) vzdialené \( 3\,\mathrm{km} \) a \( 4\,\mathrm{km} \) (viď obrázok). Z miesta dopadu sa výsadkár vydal náhodným smerom priamočiaro rýchlosťou \( 6\,\mathrm{km}/\mathrm{h} \). Aká je pravdepodobnosť, že najneskôr za hodinu príde na niektorú z ciest? Výsledok zaokrúhlite na \( 4 \) desatinné miesta.
\[ \]
Tip: V prípade rovnomerného priamočiareho pohybu je rýchlosť rovná pomeru dráhy a času.}
{\( 0{,}5505 \)}{\( 0{,}4495 \)}{\( 0{,}6011 \)}{\( 0{,}3989 \)}{\( 0{,}3511 \)}{\( 0{,}6489 \)}}
\LANGes{
\Question{
Un paracaídista bajó al lugar \( M \), que dista \( 3\,\mathrm{km} \) y  \( 4\,\mathrm{km} \) de dos caminos \( p \) y \( q \) perpendiculares. (vea el dibujo)  Fue por un movimiento uniforme recto con velocidad \( 6\,\mathrm{km}/\mathrm{h} \) en una dirección aleatoria. ¿Qué probabilidad hay de que va a cruzar alguno de los caminos dentro de una hora? Aproxima el resultado a \( 4 \) cifras decimales.
Pista: En el caso del movimiento uniforme recto es la velocidad igual a la proporción entre la distancia y el tiempo.}
{\( 0{,}5505 \)}{\( 0{,}4495 \)}{\( 0{,}6011 \)}{\( 0{,}3989 \)}{\( 0{,}3511 \)}{\( 0{,}6489 \)}}
\End

\Data\ID{9000138301}
\Updated{2021-02-10 02:02:39}
\Level{B}
\SubArea{Probability}
\LANGen{
\Question{
Two dices are rolled. Find the probability that the sum of the numbers on both dices does not
exceed the number \(6\).}
{\(\frac{15}
{36}\,\dot{=}\,0{.}4167\)}{\(\frac{10}
{36}\,\dot{=}\,0{.}2778\)}{\(\frac{18}
{36}=0{.}5\)}{\(\frac{12}
{36}\,\dot{=}\,0{.}3333\)}{}{}}
\LANGpl{
\Question{
Rzucamy dwoma  kostkami. Jakie jest prawdopodobieństwo, że suma wyrzuconych liczb  wynosi co najwyżej \(6\)?}
{\(\frac{15}
{36}\,\dot{=}\,0{,}4167\)}{\(\frac{10}
{36}\,\dot{=}\,0{,}2778\)}{\(\frac{18}
{36}=0{,}5\)}{\(\frac{12}
{36}\,\dot{=}\,0{,}3333\)}{}{}}
\LANGcs{
\Question{
Hodíme dvěma kostkami. Jaká je pravděpodobnost, že součet bude nejvýše
\(6\)?}
{\(\frac{15}
{36}\,\dot{=}\,0{,}4167\)}{\(\frac{10}
{36}\,\dot{=}\,0{,}2778\)}{\(\frac{18}
{36}=0{,}5\)}{\(\frac{12}
{36}\,\dot{=}\,0{,}3333\)}{}{}}
\LANGsk{
\Question{
Hodíme dvoma kockami. Aká je pravdepodobnosť, že súčet bude najviac \(6\)?}
{\(\frac{15}
{36}\,\dot{=}\,0{,}4167\)}{\(\frac{10}
{36}\,\dot{=}\,0{,}2778\)}{\(\frac{18}
{36}=0{,}5\)}{\(\frac{12}
{36}\,\dot{=}\,0{,}3333\)}{}{}}
\LANGes{
\Question{
Vamos a tirar dos dados. ¿Qué probabilidad hay de que la suma de los puntos sea por lo máximo
\(6\)?}
{\(\frac{15}
{36}\,\dot{=}\,0{,}4167\)}{\(\frac{10}
{36}\,\dot{=}\,0{,}2778\)}{\(\frac{18}
{36}=0{,}5\)}{\(\frac{12}
{36}\,\dot{=}\,0{,}3333\)}{}{}}
\End

\Data\ID{9000138302}
\Updated{2021-02-16 02:02:48}
\Level{B}
\SubArea{Probability}
\LANGen{
\Question{
Two dices are rolled. Find the probability that we get either at least one number
\(6\) or the sum of the
numbers on both dices is \(8\).}
{\(\frac{14}
{36}\,\dot{=}\,0{.}3889\)}{\(\frac{16}
{36}\,\dot{=}\,0{.}4444\)}{\(\frac{11}
{36}\,\dot{=}\,0{.}3056\)}{\(\frac{5}
{36}\,\dot{=}\,0{.}1389\)}{}{}}
\LANGpl{
\Question{
Rzucamy dwoma  kostkami. Jakie jest prawdopodobieństwo, że przynajmniej na jednej kostce wyrzucimy
\(6\) lub  suma liczb na obu kostkach będzie równa  \(8\)?}
{\(\frac{14}
{36}\,\dot{=}\,0{,}3889\)}{\(\frac{16}
{36}\,\dot{=}\,0{,}4444\)}{\(\frac{11}
{36}\,\dot{=}\,0{,}3056\)}{\(\frac{5}
{36}\,\dot{=}\,0{,}1389\)}{}{}}
\LANGcs{
\Question{
Hodíme dvěma kostkami. Jaká je pravděpodobnost, že součet bude
\(8\) nebo že alespoň na
jedné kostce padne \(6\)?}
{\(\frac{14}
{36}\,\dot{=}\,0{,}3889\)}{\(\frac{16}
{36}\,\dot{=}\,0{,}4444\)}{\(\frac{11}
{36}\,\dot{=}\,0{,}3056\)}{\(\frac{5}
{36}\,\dot{=}\,0{,}1389\)}{}{}}
\LANGsk{
\Question{
Hodíme dvoma kockami. Aká je pravdepodobnosť, že aspoň na jednej kocke padne
\(6\) alebo súčet bude \(8\)?}
{\(\frac{14}
{36}\,\dot{=}\,0{,}3889\)}{\(\frac{16}
{36}\,\dot{=}\,0{,}4444\)}{\(\frac{11}
{36}\,\dot{=}\,0{,}3056\)}{\(\frac{5}
{36}\,\dot{=}\,0{,}1389\)}{}{}}
\LANGes{
\Question{
Vamos a tirar dos dados. ¿Qué probabilidad hay de que la suma sea 
\(8\) o que por lo mínimo un dado tenga \(6\)?}
{\(\frac{14}
{36}\,\dot{=}\,0{,}3889\)}{\(\frac{16}
{36}\,\dot{=}\,0{,}4444\)}{\(\frac{11}
{36}\,\dot{=}\,0{,}3056\)}{\(\frac{5}
{36}\,\dot{=}\,0{,}1389\)}{}{}}
\End

\Data\ID{9000138304}
\Updated{2021-02-16 02:02:29}
\Level{B}
\SubArea{Probability}
\IMG{001383_kostky2_0.png}
\LANGen{
\Question{
Two different dices (a white dice and a black dice) are rolled. Find the probability that we get the number
\(3\) on the black dice and a
number different from \(3\) 
on a white dice.}
{\(\frac{5}
{36}\,\dot{=}\,0{.}1389\)}{\(\frac{3}
{36}\,\dot{=}\,0{.}0833\)}{\(\frac{6}
{36}\,\dot{=}\,0{.}1667\)}{\(\frac{1}
{36}\,\dot{=}\,0{.}0278\)}{}{}}
\LANGpl{
\Question{
Rzucamy dwoma  kostkami  (białą i czarną). Jakie jest prawdopodobieństwo, że wyrzucimy
\(3\) na czarnej kostce i liczbę różną od 
 \(3\) na białej kostce?}
{\(\frac{5}
{36}\,\dot{=}\,0{,}1389\)}{\(\frac{3}
{36}\,\dot{=}\,0{,}0833\)}{\(\frac{6}
{36}\,\dot{=}\,0{,}1667\)}{\(\frac{1}
{36}\,\dot{=}\,0{,}0278\)}{}{}}
\LANGcs{
\Question{
Hodíme dvěma kostkami, bílou a černou. Jaká je pravděpodobnost, že na černé
kostce padne \(3\) a
na bílé kostce \(3\) 
nepadne?}
{\(\frac{5}
{36}\,\dot{=}\,0{,}1389\)}{\(\frac{3}
{36}\,\dot{=}\,0{,}0833\)}{\(\frac{6}
{36}\,\dot{=}\,0{,}1667\)}{\(\frac{1}
{36}\,\dot{=}\,0{,}0278\)}{}{}}
\LANGsk{
\Question{
Hodíme dvoma rozdielnymi kockami, bielou a čiernou. Aká je pravdepodobnosť, že na čiernej kocke padne
\(3\) a na bielej kocke 
\(3\) 
nepadne?}
{\(\frac{5}
{36}\,\dot{=}\,0{,}1389\)}{\(\frac{3}
{36}\,\dot{=}\,0{,}0833\)}{\(\frac{6}
{36}\,\dot{=}\,0{,}1667\)}{\(\frac{1}
{36}\,\dot{=}\,0{,}0278\)}{}{}}
\LANGes{
\Question{
Vamos a tirar dos dados, un negro y un blanco. ¿Qué probabilidad hay de obtener un \(3\) en el dado negro y no obtener un  \(3\) en el dado blanco?}
{\(\frac{5}
{36}\,\dot{=}\,0{,}1389\)}{\(\frac{3}
{36}\,\dot{=}\,0{,}0833\)}{\(\frac{6}
{36}\,\dot{=}\,0{,}1667\)}{\(\frac{1}
{36}\,\dot{=}\,0{,}0278\)}{}{}}
\End

\Data\ID{9000138306}
\Updated{2021-02-16 02:02:08}
\Level{B}
\SubArea{Probability}
\LANGen{
\Question{
Two dices are rolled. Find the probability that we get at least once the number
\(3\).}
{\(\frac{11}
{36}\,\dot{=}\,0{.}3056\)}{\(\frac{3}
{36}\,\dot{=}\,0{.}0833\)}{\(\frac{10}
{36}\,\dot{=}\,0{.}2778\)}{\(\frac{12}
{36}\,\dot{=}\,0{.}3333\)}{}{}}
\LANGpl{
\Question{
Rzucamy dwoma  kostkami. Jakie jest prawdopodobieństwo, że przynajmniej raz  wyrzucimy liczbę
\(3\)?}
{\(\frac{11}
{36}\,\dot{=}\,0{,}3056\)}{\(\frac{3}
{36}\,\dot{=}\,0{,}0833\)}{\(\frac{10}
{36}\,\dot{=}\,0{,}2778\)}{\(\frac{12}
{36}\,\dot{=}\,0{,}3333\)}{}{}}
\LANGcs{
\Question{
Hodíme dvěma kostkami. Jaká je pravděpodobnost, že alespoň jednou padne
\(3\)?}
{\(\frac{11}
{36}\,\dot{=}\,0{,}3056\)}{\(\frac{3}
{36}\,\dot{=}\,0{,}0833\)}{\(\frac{10}
{36}\,\dot{=}\,0{,}2778\)}{\(\frac{12}
{36}\,\dot{=}\,0{,}3333\)}{}{}}
\LANGsk{
\Question{
Hodíme dvoma kockami. Aká je pravdepodobnosť, že aspoň raz padne
\(3\)?}
{\(\frac{11}
{36}\,\dot{=}\,0{,}3056\)}{\(\frac{3}
{36}\,\dot{=}\,0{,}0833\)}{\(\frac{10}
{36}\,\dot{=}\,0{,}2778\)}{\(\frac{12}
{36}\,\dot{=}\,0{,}3333\)}{}{}}
\LANGes{
\Question{
Vamos a tirar dos dados. ¿Qué probabilidad hay de obtener por lo mínimo un
\(3\)?}
{\(\frac{11}
{36}\,\dot{=}\,0{,}3056\)}{\(\frac{3}
{36}\,\dot{=}\,0{,}0833\)}{\(\frac{10}
{36}\,\dot{=}\,0{,}2778\)}{\(\frac{12}
{36}\,\dot{=}\,0{,}3333\)}{}{}}
\End

\Data\ID{9000138307}
\Updated{2021-02-03 03:02:52}
\Level{B}
\SubArea{Probability}
\LANGen{
\Question{
Two dices are rolled. Find the probability that we get the product
\(10\).}
{\(\frac{2}
{36}\,\dot{=}\,0{.}0556\)}{\(\frac{10}
{36}\,\dot{=}\,0{.}2778\)}{\(\frac{1}
{36}\,\dot{=}\,0{.}0278\)}{\(\frac{5}
{36}\,\dot{=}\,0{.}1389\)}{}{}}
\LANGpl{
\Question{
Rzucamy dwoma kostkami. Jakie jest prawdopodobieństwo, że wyrzucimy iloczyn 
\(10\)?}
{\(\frac{2}
{36}\,\dot{=}\,0{,}0556\)}{\(\frac{10}
{36}\,\dot{=}\,0{,}2778\)}{\(\frac{1}
{36}\,\dot{=}\,0{,}0278\)}{\(\frac{5}
{36}0{,}1389\)}{}{}}
\LANGcs{
\Question{
Hodíme dvěma kostkami. Jaká je pravděpodobnost, že součin bude
\(10\)?}
{\(\frac{2}
{36}\,\dot{=}\,0{,}0556\)}{\(\frac{10}
{36}\,\dot{=}\,0{,}2778\)}{\(\frac{1}
{36}\,\dot{=}\,0{,}0278\)}{\(\frac{5}
{36}\,\dot{=}\,0{,}1389\)}{}{}}
\LANGsk{
\Question{
Hodíme dvoma kockami. Aká je pravdepodobnosť, že súčin bude 
\(10\)?}
{\(\frac{2}
{36}\,\dot{=}\,0{,}0556\)}{\(\frac{10}
{36}\,\dot{=}\,0{,}2778\)}{\(\frac{1}
{36}\,\dot{=}\,0{,}0278\)}{\(\frac{5}
{36}\,\dot{=}\,0{,}1389\)}{}{}}
\LANGes{
\Question{
Vamos a tirar dos dados. ¿Qué es la probabilidad de que la suma sea
\(10\)?}
{\(\frac{2}
{36}\,\dot{=}\,0{,}0556\)}{\(\frac{10}
{36}\,\dot{=}\,0{,}2778\)}{\(\frac{1}
{36}\,\dot{=}\,0{,}0278\)}{\(\frac{5}
{36}\,\dot{=}\,0{,}1389\)}{}{}}
\End

\Data\ID{9000138309}
\Updated{2021-02-03 03:02:32}
\Level{B}
\SubArea{Probability}
\LANGen{
\Question{
Two dices are rolled. Find the probability that we get either the
same number on both dices or the sum of the numbers on both dices is
\(6\).}
{\(\frac{10}
{36}\,\dot{=}\,0{.}2778\)}{\(\frac{11}
{36}\,\dot{=}\,0{.}3056\)}{\(\frac{6}
{36}\,\dot{=}\,0{.}1667\)}{\(\frac{5}
{36}\,\dot{=}\,0{.}1389\)}{}{}}
\LANGpl{
\Question{
Rzucamy dwoma  kostkami. Jakie jest prawdopodobieństwo, że wyrzucimy tę samą liczbę na obu kostkach lub suma liczb na obu kostkach jest równa
\(6\)?}
{\(\frac{10}
{36}\,\dot{=}\,0{,}2778\)}{\(\frac{11}
{36}\,\dot{=}\,0{,}3056\)}{\(\frac{6}
{36}\,\dot{=}\,0{,}1667\)}{\(\frac{5}
{36}\,\dot{=}\,0{,}1389\)}{}{}}
\LANGcs{
\Question{
Hodíme dvěma kostkami. Jaká je pravděpodobnost, že součet bude
\(6\) nebo
že na obou kostkách padne stejné číslo?}
{\(\frac{10}
{36}\,\dot{=}\,0{,}2778\)}{\(\frac{11}
{36}\,\dot{=}\,0{,}3056\)}{\(\frac{6}
{36}\,\dot{=}\,0{,}1667\)}{\(\frac{5}
{36}\,\dot{=}\,0{,}1389\)}{}{}}
\LANGsk{
\Question{
Hodíme dvoma kockami. Aká je pravdepodobnosť, že na obidvoch kockách padne rovnaké číslo
alebo súčet bude
\(6\)?}
{\(\frac{10}
{36}\,\dot{=}\,0{,}2778\)}{\(\frac{11}
{36}\,\dot{=}\,0{,}3056\)}{\(\frac{6}
{36}\,\dot{=}\,0{,}1667\)}{\(\frac{5}
{36}\,\dot{=}\,0{,}1389\)}{}{}}
\LANGes{
\Question{
Vamos a tirar dos dados. ¿Qué es la probabilidad de que la suma sea
\(6\) o los dos dados lleven el mismo número?}
{\(\frac{10}
{36}\,\dot{=}\,0{,}2778\)}{\(\frac{11}
{36}\,\dot{=}\,0{,}3056\)}{\(\frac{6}
{36}\,\dot{=}\,0{,}1667\)}{\(\frac{5}
{36}\,\dot{=}\,0{,}1389\)}{}{}}
\End

\Data\ID{9000154803}
\Updated{2021-02-17 02:02:10}
\Level{B}
\SubArea{Probability}
\LANGen{
\Question{
The probability that Robin Hood hits the target is
\(0.83\). The same
probability is \(0.61\) 
for Robin's fellow, Little John. Both Robin and John shot on the wolf. Find the
probability that they will hit the wolf. Round your answer to three decimal places.}
{\(0.934\)}{\(1.440\)}{\(0.506\)}{\(0.494\)}{}{}}
\LANGpl{
\Question{
Prawdopodobieństwo strzelenia do celu przez Robina wynosi
\(0{,}83\). Prawdopodobieństwo strzelenia do celu jego kolegi Małego Johna 
wynosi  \(0{,}61\).  Robin i John strzelili do wilka. Jakie jest prawdopodobieństwo trafienia wilka? Zaokrągli odpowiedź do trzech miejsc po przecinku.}
{\(0{,}934\)}{\(1{,}440\)}{\(0{,}506\)}{\(0{,}494\)}{}{}}
\LANGcs{
\Question{
Robin Hood zasáhne cíl s pravděpodobností
\(0{,}83\), Malý John s
pravděpodobností \(0{,}61\).
S jakou pravděpodobností skolí vlka, pokud střílí oba najednou? Výsledek
zaokrouhlete na \(3\) 
desetinná místa.}
{\(0{,}934\)}{\(1{,}440\)}{\(0{,}506\)}{\(0{,}494\)}{}{}}
\LANGsk{
\Question{
Robin Hood zasiahne cieľ s pravdepodobnosťou
\(0{,}83\), Malý John s pravdepodobnosťou \(0{,}61\). S akou pravdepodobnosťou zasiahnu vlka, ak strieľajú obidvaja naraz? Výsledok zaokrúhlite na tri desatinné miesta.}
{\(0{,}934\)}{\(1{,}440\)}{\(0{,}506\)}{\(0{,}494\)}{}{}}
\LANGes{
\Question{
Robin Hood acierta la diana con probabilidad de
\(0{,}83\), El Pequeño Juan a cierta con probabilidad de \(0{,}61\).
¿Qué probabilidad hay de que maten a un lobo si disparan en el mismmo momento? Aproxima el resultado a \(3\) 
cifras decimales.}
{\(0{,}934\)}{\(1{,}440\)}{\(0{,}506\)}{\(0{,}494\)}{}{}}
\End

\Data\ID{9000154806}
\Updated{2021-02-17 02:02:34}
\Level{B}
\SubArea{Probability}
\IMG{9000154808_dices_0_0.png}
\LANGen{
\Question{
A man plays a dice game. He rolls one dice three times and to win he needs at least
one number six. However, the dice is not fair. The probability of an even number is
two times bigger than the probability of an odd number. Find the probability that
the man wins. Suppose that all even numbers have equal probability and all odd
numbers also have equal probability. Round your answer to three decimal
places.}
{\(0.529\)}{\(0.471\)}{\(0.421\)}{\(0.579\)}{}{}}
\LANGpl{
\Question{
Mężczyzna gra w kości.  Rzuca jedną kostką trzy razy, aby wygrać musi przynajmniej raz wyrzucić szóstkę.
Kostki są jednak niesprawiedliwe.  Prawdopodobieństwo wyrzucenia liczby parzystej jest dwa razy większe, niż liczby nieparzystej.
Jakie jest prawdopodobieństwo, że mężczyzna wygra? Załóżmy, że prawdopodobieństwo wyrzucenia liczby parzystej jest takie same dla wszystkich liczb parzystych oraz prawdopodobieństwo wyrzucenia liczby nieparzystej jest takie same dla wszystkich liczb nieparzystych. 
Zaokrągli swoją odpowiedź do trzech miejsc po przecinku.}
{\(0{,}529\)}{\(0{,}471\)}{\(0{,}421\)}{\(0{,}579\)}{}{}}
\LANGcs{
\Question{
Malý John hraje v kasinu kostky a pro výhru musí hodit alespoň jednu
šestku ve třech hodech. (Hází vždy jen jednou kostkou.) Hraje ale s
cinknutými kostkami, u kterých sudá čísla padají dvakrát častěji
než lichá. Jakou má šanci, že uspěje? Výsledek zaokrouhlete na
\(3\) 
desetinná místa.}
{\(0{,}529\)}{\(0{,}471\)}{\(0{,}421\)}{\(0{,}579\)}{}{}}
\LANGsk{
\Question{
John hrá hru s kockami. Hádže jednou kockou tri krát a aby vyhral, potrebuje hodiť aspoň jednu šestku. Avšak kocka nie je obyčajná, pravdepodobnosť padnutia párneho čísla je dvakrát väčšia než pravdepodobnosť padnutia nepárneho čísla. 
Aká je pravdepodobnosť, že John vyhrá? Výsledok zaokrúhlite na 3 desatinné miesta.}
{\(0{,}529\)}{\(0{,}471\)}{\(0{,}421\)}{\(0{,}579\)}{}{}}
\LANGes{
\Question{
El Pequeño Juan juega a dados en un casino y tiene que tirar por lo mínimo un seis en tres tiros para ganar. (Cada vez tira un dado.) Está jugando con un dado de truco, en el cuál los número pares caen dos veces más que los impares. ¿Qué probabilidad de éxito tiene? Aproxima el resultado a
\(3\) 
cifras decimales.}
{\(0{,}529\)}{\(0{,}471\)}{\(0{,}421\)}{\(0{,}579\)}{}{}}
\End

\Data\ID{9000154807}
\Updated{2021-02-17 10:02:46}
\Level{B}
\SubArea{Probability}
\LANGen{
\Question{
The Robin's band has \(10\) 
men and \(5\) 
women. From this group they select randomly two fellows to negotiate with the
Sheriff of Nottingham. Find the probability that there will be one man and one
woman in the selected pair. Round your answer to three decimal places.}
{\(0.476\)}{\(0.952\)}{\(0.325\)}{\(0.675\)}{}{}}
\LANGpl{
\Question{
W grupie Robina jest  \(10\) 
mężczyzn i  \(5\) 
kobiet.  Z tej grupy wybierają losowo dwóch towarzyszy do negocjacji z
szeryfem  z Nottingham. Jakie jest  prawdopodobieństwo, że będzie to  jeden mężczyzna  i jedna
kobieta. Zaokrągli odpowiedź do trzech miejsc po przecinku.}
{\(0{,}476\)}{\(0{,}952\)}{\(0{,}325\)}{\(0{,}675\)}{}{}}
\LANGcs{
\Question{
Družina ze Sherwoodu, \(10\) 
mužů a \(5\) žen,
vybírá své \(2\) 
zástupce pro jednání s šerifem z Nottinghamu. Jaká je pravděpodobnost,
že budou vybráni právě muž a žena? Výsledek zaokrouhlete na
\(3\) 
desetinná místa.}
{\(0{,}476\)}{\(0{,}952\)}{\(0{,}325\)}{\(0{,}675\)}{}{}}
\LANGsk{
\Question{
Robinova družina má \(10\) 
mužov a \(5\) 
žien. Z tejto skupiny sa náhodne vyberú dvaja zástupcovia pre jednanie so šerifom z Nottinghamu. Aká je pravdepodobnosť, že budú vybraní práve muž a žena? Výsledok zaokrúhlite na tri desatinné miesta.}
{\(0{,}476\)}{\(0{,}952\)}{\(0{,}325\)}{\(0{,}675\)}{}{}}
\LANGes{
\Question{
La banda de Sherwood, \(10\) 
hombres y \(5\) mujeres,
está eligiendo \(2\) 
representantes para hablar con el sherif de Nottingham. ¿Qué probabilidad hay de elegir exactamente un hombre y una mujer? Aproxima el resultado a
\(3\) 
cifras decimales.}
{\(0{,}476\)}{\(0{,}952\)}{\(0{,}325\)}{\(0{,}675\)}{}{}}
\End

\Data\ID{1003019205}
\Updated{2021-02-16 02:02:37}
\Level{C}
\SubArea{Probability}
\LANGen{
\Question{
Adam and Eve met at the disco. They agreed  to meet the next  day at the same location sometime between \( 1 \) p.m. and \( 2 \) p.m. Both of them will arrive independently at random times within the hour and wait  ten minutes for the other. What is the probability that they will not meet during that hour?}
{\( \frac{25}{36}\dot{=}\,0{.}6944 \)}{\( \frac{11}{36}\dot{=}\,0{.}3056 \)}{\( \frac{35}{36}\dot{=}\,0{.}9722 \)}{\( \frac{24}{36}\dot{=}\,0{.}6667 \)}{}{}}
\LANGpl{
\Question{
Adam i Ewa poznali się na dyskotece. Postanowili  spotkać się następnego dnia między godziną \( 13 \)  a  \( 14 \). Oboje przyjeżdżają na miejsce spotkania w losowo wybranym czasie w ciągu jednej godziny i czekają na siebie dziesięć minut. Jakie jest prawdopodobieństwo, że nie spotkają się w ciągu tej godziny?}
{\( \frac{25}{36}\dot{=}\,0{,}6944 \)}{\( \frac{11}{36}\dot{=}\,0{,}3056 \)}{\( \frac{35}{36}\dot{=}\,0{,}9722 \)}{\( \frac{24}{36}\dot{=}\,0{,}6667 \)}{}{}}
\LANGcs{
\Question{
Adam a Eva se setkali na diskotéce. Dohodli se, že se potkají druhý den mezi \( 13 \). a \( 14 \). hodinou. Dohodli se, že na sebe budou čekat  \( 10 \) minut. Jejich příchody na dané místo jsou navzájem nezávislé a stejně pravděpodobné v průběhu dané hodiny. Jaká je pravděpodobnost, že se nepotkají?}
{\( \frac{25}{36}\dot{=}\,0{,}6944 \)}{\( \frac{11}{36}\dot{=}\,0{,}3056 \)}{\( \frac{35}{36}\dot{=}\,0{,}9722 \)}{\( \frac{24}{36}\dot{=}\,0{,}6667 \)}{}{}}
\LANGsk{
\Question{
Adam a Eva sa stretli na diskotéke. Dohodli sa, že sa stretnú na druhý deň medzi \( 13 \). a \( 14 \). hodinou. Každý z nich bude na druhého čakať  \( 10 \) minút. Ich príchody na dané miesto sú navzájom nezávislé a rovnako pravdepodobné v priebehu danej hodiny. Aká je pravdepodobnosť, že sa nestretnú?}
{\( \frac{25}{36}\dot{=}\,0{,}6944 \)}{\( \frac{11}{36}\dot{=}\,0{,}3056 \)}{\( \frac{35}{36}\dot{=}\,0{,}9722 \)}{\( \frac{24}{36}\dot{=}\,0{,}6667 \)}{}{}}
\LANGes{
\Question{
Adam y Eva se han conocido en una discoteca. Han decidido quedar el siguiente día entre  \( 13 \). y \( 14 \). hora. Han decidido esperar  \( 10 \) minutos. El tiempo de llegada de cada uno de ellos es aleatorio e independiente del otro. ¿Qué probabilidad hay de que Adam y Eva no quedan?}
{\( \frac{25}{36}\dot{=}\,0{,}6944 \)}{\( \frac{11}{36}\dot{=}\,0{,}3056 \)}{\( \frac{35}{36}\dot{=}\,0{,}9722 \)}{\( \frac{24}{36}\dot{=}\,0{,}6667 \)}{}{}}
\End

\Data\ID{1003029203}
\Updated{2021-01-29 01:01:59}
\Level{C}
\SubArea{Probability}
\LANGen{
\Question{
Three dice are thrown together. Find the probability that three different outcomes are rolled.  Results are rounded to two decimal places.}
{\( \frac{\binom61\cdot\binom51\cdot\binom41}{6^3}=0{.}56 \)}{\( \frac{\binom61+\binom51+\binom41}{6^3}=0{.}07 \)}{\( \frac{\binom66\cdot\binom65\cdot\binom64}{6^3}=0{.}42 \)}{\( \frac{\binom66+\binom65+\binom64}{6^3}=0{.}10 \)}{}{}}
\LANGpl{
\Question{
Trzy kostki zostały rzucone jednocześnie. Jakie jest prawdopodobieństwo, że zostaną wyrzucone trzy różne wyniki?  Wyniki są zaokrąglane do dwóch miejsc po przecinku.}
{\( \frac{\binom61\cdot\binom51\cdot\binom41}{6^3}=0{,}56 \)}{\( \frac{\binom61+\binom51+\binom41}{6^3}=0{,}07 \)}{\( \frac{\binom66\cdot\binom65\cdot\binom64}{6^3}=0{,}42 \)}{\( \frac{\binom66+\binom65+\binom64}{6^3}=0{,}10 \)}{}{}}
\LANGcs{
\Question{
Házíme třemi různými kostkami. Jaká je pravděpodobnost, že na těchto kostkách padnou navzájem různá čísla? Výsledky jsou zaokrouhlené na dvě desetinná místa.}
{\( \frac{\binom61\cdot\binom51\cdot\binom41}{6^3}=0{,}56 \)}{\( \frac{\binom61+\binom51+\binom41}{6^3}=0{,}07 \)}{\( \frac{\binom66\cdot\binom65\cdot\binom64}{6^3}=0{,}42 \)}{\( \frac{\binom66+\binom65+\binom64}{6^3}=0{,}10 \)}{}{}}
\LANGsk{
\Question{
Hádžeme tromi rôznymi kockami. Aká je pravdepodobnosť, že na týchto kockách padnú navzájom rôzne čísla? Výsledky sú zaokrúhlené na dve desatinné miesta.}
{\( \frac{\binom61\cdot\binom51\cdot\binom41}{6^3}=0{,}56 \)}{\( \frac{\binom61+\binom51+\binom41}{6^3}=0{,}07 \)}{\( \frac{\binom66\cdot\binom65\cdot\binom64}{6^3}=0{,}42 \)}{\( \frac{\binom66+\binom65+\binom64}{6^3}=0{,}10 \)}{}{}}
\LANGes{
\Question{
Vamos a tirar tres dados distintos. ¿Qué es la probabilidad de que saquemos tres números distintos? Aproxima los resultados a dos cifras decimales.}
{\( \frac{\binom61\cdot\binom51\cdot\binom41}{6^3}=0{,}56 \)}{\( \frac{\binom61+\binom51+\binom41}{6^3}=0{,}07 \)}{\( \frac{\binom66\cdot\binom65\cdot\binom64}{6^3}=0{,}42 \)}{\( \frac{\binom66+\binom65+\binom64}{6^3}=0{,}10 \)}{}{}}
\End

\Data\ID{1003029204}
\Updated{2021-02-04 03:02:13}
\Level{C}
\SubArea{Probability}
\LANGen{
\Question{
A class consists of \( 50 \) students including twins, Mark and Martin. For an exam students are randomly divided into two equally sized subgroups. Find the probability that Mark and Martin will be in the same  subgroup. The results are rounded to two decimal places.}
{\( \frac{\binom{48}{23}+\binom{48}{25}}{\binom{50}{25}}=0{.}49 \)}{\( \frac{\binom{48}{23}}{\binom{50}{25}}=0{.}24 \)}{\( \frac{2\cdot\binom{48}{24}}{\binom{50}{25}}=0{.}51 \)}{\( \frac{\binom{49}{24}}{\binom{50}{25}}=0{.}50 \)}{}{}}
\LANGpl{
\Question{
Klasa liczy \( 50 \) uczniów, w tym bliźniaków, Marka i Marcina. Do egzaminu uczniowie są losowo podzieleni na dwie podgrupy o jednakowej wielkości. Wyznacz prawdopodobieństwo, że Marek i Marcin będą w tej samej podgrupie. Wyniki są zaokrąglane do dwóch miejsc po przecinku.}
{\( \frac{\binom{48}{23}+\binom{48}{25}}{\binom{50}{25}}=0{,}49 \)}{\( \frac{\binom{48}{23}}{\binom{50}{25}}=0{,}24 \)}{\( \frac{2\cdot\binom{48}{24}}{\binom{50}{25}}=0{,}51 \)}{\( \frac{\binom{49}{24}}{\binom{50}{25}}=0{,}50 \)}{}{}}
\LANGcs{
\Question{
Padesát žáků 3. ročníku bude psát test z matematiky, proto je třeba je náhodně rozdělit do dvou stejně velkých skupin. Jaká je pravděpodobnost, že dvojčata Martin a Matěj, kteří jsou mezi nimi, budou zařazeni do stejné skupiny? Výsledky jsou zaokrouhlené na dvě desetinná místa.}
{\( \frac{\binom{48}{23}+\binom{48}{25}}{\binom{50}{25}}=0{,}49 \)}{\( \frac{\binom{48}{23}}{\binom{50}{25}}=0{,}24 \)}{\( \frac{2\cdot\binom{48}{24}}{\binom{50}{25}}=0{,}51 \)}{\( \frac{\binom{49}{24}}{\binom{50}{25}}=0{,}50 \)}{}{}}
\LANGsk{
\Question{
Päťdesiat žiakov 3. ročníka sa bude testovať z matematiky, preto ich treba náhodne rozdeliť do dvoch rovnako veľkých skupín. Aká je pravdepodobnosť, že dvojčatá Martin a Matej, ktorí sú medzi nimi, budú testovaní v tej istej skupine? Výsledky sú zaokrúhlené na dve desatinné miesta.}
{\( \frac{\binom{48}{23}+\binom{48}{25}}{\binom{50}{25}}=0{,}49 \)}{\( \frac{\binom{48}{23}}{\binom{50}{25}}=0{,}24 \)}{\( \frac{2\cdot\binom{48}{24}}{\binom{50}{25}}=0{,}51 \)}{\( \frac{\binom{49}{24}}{\binom{50}{25}}=0{,}50 \)}{}{}}
\LANGes{
\Question{
Cincuenta alumnos de tercera ESO van a hacer un examen de mates y por eso hay que dividirlos en dos grupos. ¿Qué es la probabilidad de que los gemelos Pedro y Pablo (que están en la clase) estén en el mismo grupo? Aproxima los resultados a dos cifras decimales.}
{\( \frac{\binom{48}{23}+\binom{48}{25}}{\binom{50}{25}}=0{,}49 \)}{\( \frac{\binom{48}{23}}{\binom{50}{25}}=0{,}24 \)}{\( \frac{2\cdot\binom{48}{24}}{\binom{50}{25}}=0{,}51 \)}{\( \frac{\binom{49}{24}}{\binom{50}{25}}=0{,}50 \)}{}{}}
\End

\Data\ID{1003029205}
\Updated{2021-01-29 02:01:00}
\Level{C}
\SubArea{Probability}
\LANGen{
\Question{
In the hospital, \( 22 \) boys and \( 18 \) girls were born in one month. Babies were listed in the register by their date of birth. Find the probability that there are two boys and three girls in the first five places of the  register. The results are rounded to four decimal places.}
{\( \frac{\binom{22}2\cdot\binom{18}3}{\binom{40}5}=0{.}2865 \)}{\( \frac{\binom{22}2\cdot\binom{18}3}{\frac{40!}{35!}}=0{.}0024 \)}{\( \frac{22^2\cdot18^3}{40^5} = 0{.}0276 \)}{\( \frac{\binom{22}3\cdot\binom{18}2}{\frac{40!}{35!}}=0{.}0030 \)}{}{}}
\LANGpl{
\Question{
W szpitalu w ciągu jednego miesiąca urodziło się \( 22 \) chłopców i \( 18 \) dziewczynek. Niemowlęta zostały zapisane w rejestrze według daty urodzenia. Wyznacz prawdopodobieństwo, że na pierwszych pięciu miejscach rejestru jest dwóch chłopców i trzy dziewczynki. Wyniki są zaokrąglane do czterech miejsc po przecinku.}
{\( \frac{\binom{22}2\cdot\binom{18}3}{\binom{40}5}=0{,}2865 \)}{\( \frac{\binom{22}2\cdot\binom{18}3}{\frac{40!}{35!}}=0{,}0024 \)}{\( \frac{22^2\cdot18^3}{40^5} = 0{,}0276 \)}{\( \frac{\binom{22}3\cdot\binom{18}2}{\frac{40!}{35!}}=0{,}0030 \)}{}{}}
\LANGcs{
\Question{
V určité porodnici se v jednom měsíci narodilo \( 22 \) chlapců a \( 18 \) děvčat. Děti jsou na seznam zapisovány podle svého data narození. Jaká je pravděpodobnost, že na prvních pěti místech tohoto seznamu jsou dva chlapci a tři děvčata? Výsledky jsou zaokrouhlené na čtyři desetinná místa.}
{\( \frac{\binom{22}2\cdot\binom{18}3}{\binom{40}5}=0{,}2865 \)}{\( \frac{\binom{22}2\cdot\binom{18}3}{\frac{40!}{35!}}=0{,}0024 \)}{\( \frac{22^2\cdot18^3}{40^5} = 0{,}0276 \)}{\( \frac{\binom{22}3\cdot\binom{18}2}{\frac{40!}{35!}}=0{,}0030 \)}{}{}}
\LANGsk{
\Question{
V určitej pôrodnici bolo v jednom mesiaci narodených \( 22 \) chlapcov a \( 18 \) dievčat. Deti sú do zoznamu zapisované podľa dátumu ich narodenia. Aká je pravdepodobnosť, že na prvých piatich miestach tohto zoznamu budú dvaja chlapci a tri dievčatá? Výsledky sú zaokrúhlené na štyri desatinné miesta.}
{\( \frac{\binom{22}2\cdot\binom{18}3}{\binom{40}5}=0{,}2865 \)}{\( \frac{\binom{22}2\cdot\binom{18}3}{\frac{40!}{35!}}=0{,}0024 \)}{\( \frac{22^2\cdot18^3}{40^5} = 0{,}0276 \)}{\( \frac{\binom{22}3\cdot\binom{18}2}{\frac{40!}{35!}}=0{,}0030 \)}{}{}}
\LANGes{
\Question{
En un hospital han nacido \( 22 \) chicos y \( 18 \) chicas este mes. Vamos a hacer una lista de los niños según el día del nacimiento. ¿Qué probabilidad hay de que en los primeros cinco lugares de la lista estén dos chicas y tres chicos?  Aproxima el resultado a cuatro cifras decimales.}
{\( \frac{\binom{22}2\cdot\binom{18}3}{\binom{40}5}=0{,}2865 \)}{\( \frac{\binom{22}2\cdot\binom{18}3}{\frac{40!}{35!}}=0{,}0024 \)}{\( \frac{22^2\cdot18^3}{40^5} = 0{,}0276 \)}{\( \frac{\binom{22}3\cdot\binom{18}2}{\frac{40!}{35!}}=0{,}0030 \)}{}{}}
\End

\Data\ID{1003029206}
\Updated{2021-02-03 02:02:58}
\Level{C}
\SubArea{Probability}
\LANGen{
\Question{
In the hospital, \( 22 \) boys and \( 18 \) girls were born in one month. Babies were listed in the register by their date of birth. Find the probability that there are at least three boys in the first five places of the register. The results are rounded to four decimal places.}
{\( \frac{\binom{22}3\cdot\binom{18}2+\binom{22}4\cdot\binom{18}1+\binom{22}5\cdot\binom{18}0}{\binom{40}5} = 0{.}5982 \)}{\( \frac{\binom{22}3+\binom{22}4+\binom{22}5}{\binom{40}5} = 0{.}0535 \)}{\( \frac{22^3\cdot18^2+22^4\cdot18^1+22^5\cdot18^0}{40^5}=0{.}1252 \)}{\( \frac{\binom{22}3\cdot\binom{18}2+\binom{22}4\cdot\binom{18}1+\binom{22}5\cdot\binom{18}0}{40^5} = 0{.}0038 \)}{}{}}
\LANGpl{
\Question{
W szpitalu w ciągu jednego miesiąca urodziło się \( 22 \) chłopców i  \( 18 \) dziewczynek. Niemowlęta zostały zapisane  w rejestrze według daty urodzenia. Jakie jest  prawdopodobieństwo, że na pierwszych pięciu miejscach rejestru znajduje się co najmniej trzech chłopców. Wyniki są zaokrąglane do czterech miejsc po przecinku.}
{\( \frac{\binom{22}3\cdot\binom{18}2+\binom{22}4\cdot\binom{18}1+\binom{22}5\cdot\binom{18}0}{\binom{40}5} = 0{,}5982 \)}{\( \frac{\binom{22}3+\binom{22}4+\binom{22}5}{\binom{40}5} = 0{,}0535 \)}{\( \frac{22^3\cdot18^2+22^4\cdot18^1+22^5\cdot18^0}{40^5}=0{,}1252 \)}{\( \frac{\binom{22}3\cdot\binom{18}2+\binom{22}4\cdot\binom{18}1+\binom{22}5\cdot\binom{18}0}{40^5} = 0{,}0038 \)}{}{}}
\LANGcs{
\Question{
V určité porodnici se během jednoho měsíce narodilo \( 22 \) chlapců a \( 18 \) děvčat. Děti jsou na seznam zapisovány podle svého data narození. Jaká je pravděpodobnost, že na prvních pěti místech tohoto seznamu jsou alespoň tři chlapci? Výsledky jsou zaokrouhlené na čtyři desetinná místa.}
{\( \frac{\binom{22}3\cdot\binom{18}2+\binom{22}4\cdot\binom{18}1+\binom{22}5\cdot\binom{18}0}{\binom{40}5} = 0{,}5982 \)}{\( \frac{\binom{22}3+\binom{22}4+\binom{22}5}{\binom{40}5} = 0{,}0535 \)}{\( \frac{22^3\cdot18^2+22^4\cdot18^1+22^5\cdot18^0}{40^5}=0{,}1252 \)}{\( \frac{\binom{22}3\cdot\binom{18}2+\binom{22}4\cdot\binom{18}1+\binom{22}5\cdot\binom{18}0}{40^5} = 0{,}0038 \)}{}{}}
\LANGsk{
\Question{
V určitej pôrodnici bolo v jednom mesiaci narodených \( 22 \) chlapcov a \( 18 \) dievčat. Deti sú do zoznamu zapisované podľa dátumu ich narodenia. Aká je pravdepodobnosť, že na prvých piatich miestach tohto zoznamu budú aspoň traja chlapci? Výsledky sú zaokrúhlené na štyri desatinné miesta.}
{\( \frac{\binom{22}3\cdot\binom{18}2+\binom{22}4\cdot\binom{18}1+\binom{22}5\cdot\binom{18}0}{\binom{40}5} = 0{,}5982 \)}{\( \frac{\binom{22}3+\binom{22}4+\binom{22}5}{\binom{40}5} = 0{,}0535 \)}{\( \frac{22^3\cdot18^2+22^4\cdot18^1+22^5\cdot18^0}{40^5}=0{,}1252 \)}{\( \frac{\binom{22}3\cdot\binom{18}2+\binom{22}4\cdot\binom{18}1+\binom{22}5\cdot\binom{18}0}{40^5} = 0{,}0038 \)}{}{}}
\LANGes{
\Question{
En un hospital han nacido  \( 22 \) chicos y \( 18 \) chicas durante un mes. Vamos a escribir los nombres de los bebés según su fecha de nacimiento. ¿Qué es la probabilidad de que en los primero 5 lugares de la lista hay por lo menos tres chicos? Aproxima los resultados a cuatro cifras decimales.}
{\( \frac{\binom{22}3\cdot\binom{18}2+\binom{22}4\cdot\binom{18}1+\binom{22}5\cdot\binom{18}0}{\binom{40}5} = 0{,}5982 \)}{\( \frac{\binom{22}3+\binom{22}4+\binom{22}5}{\binom{40}5} = 0{,}0535 \)}{\( \frac{22^3\cdot18^2+22^4\cdot18^1+22^5\cdot18^0}{40^5}=0{,}1252 \)}{\( \frac{\binom{22}3\cdot\binom{18}2+\binom{22}4\cdot\binom{18}1+\binom{22}5\cdot\binom{18}0}{40^5} = 0{,}0038 \)}{}{}}
\End

\Data\ID{1003029303}
\Updated{2021-02-03 03:02:39}
\Level{C}
\SubArea{Probability}
\LANGen{
\Question{
The probability of a man hitting a target is \( 0{.}9 \). If the man shoots three times, what is the probability that he hits the target at least once? Round the result to three decimal places.}
{\( 0{.}999 \)}{\( 0{.}729 \)}{\( 0{.}027 \)}{\( 0{.}243 \)}{}{}}
\LANGpl{
\Question{
Prawdopodobieństwo trafienia w cel przez mężczyznę wynosi \( 0{,}9 \). Jeśli mężczyzna strzela trzykrotnie, jakie jest prawdopodobieństwo, że trafi w cel przynajmniej raz?  Zaokrągli wynik do trzech miejsc po przecinku.}
{\( 0{,}999 \)}{\( 0{,}729 \)}{\( 0{,}027 \)}{\( 0{,}243 \)}{}{}}
\LANGcs{
\Question{
Pravděpodobnost, že střelec trefí terč je \( 0{,}9 \). Jaká je pravděpodobnost, že když střelec vystřelí třikrát, trefí terč alespoň jednou? Výsledek zaokrouhlete na tři desetinná místa.}
{\( 0{,}999 \)}{\( 0{,}729 \)}{\( 0{,}027 \)}{\( 0{,}243 \)}{}{}}
\LANGsk{
\Question{
Pravdepodobnosť, že strelec trafí terč je \( 0{,}9 \). Aká je pravdepodobnosť, že ak strelec vystrelí tri razy, trafí terč aspoň raz? Výsledok zaokrúhlite na tri desatinné miesta.}
{\( 0{,}999 \)}{\( 0{,}729 \)}{\( 0{,}027 \)}{\( 0{,}243 \)}{}{}}
\LANGes{
\Question{
La probabilidad de disparar al blanco con éxito es \( 0{,}9 \). Si disparamos tres veces, ¿qué es la probabilidad de disparar con éxito por lo mínimo una vez? Aproxima el resultado a tres cifras decimales.}
{\( 0{,}999 \)}{\( 0{,}729 \)}{\( 0{,}027 \)}{\( 0{,}243 \)}{}{}}
\End

\Data\ID{1003041602}
\Updated{2021-02-03 03:02:29}
\Level{C}
\SubArea{Probability}
\LANGen{
\Question{
A box contains \( 50 \) transistors and \( 4 \) of them are lower quality. From all transistors \( 5 \) are selected at random for inspection. What is the probability that of lower quality is at most one of the selected transistors?  Round the result to two decimal places.}
{\( \frac{\binom{46}5 + \binom{46}4\cdot\binom41}{\binom{50}5}\,\dot{=}\,0{.}96 \)}{\( \frac{\frac{46!}{41!}+\frac{46!}{42!}}{\frac{50!}{45!}}\,\dot{=}\,0{.}66 \)}{\( \frac{\binom{46}5 + \binom{46}4}{\binom{50}5}\,\dot{=}\,0{.}72 \)}{\( \frac{\frac{46!}{41!}+\frac{46!}{42!}\cdot \frac{4!}{3!}}{\frac{50!}{45!}}\,\dot{=}\,0{.}71 \)}{}{}}
\LANGpl{
\Question{
W pudełku jest  \( 50 \) tranzystorów,  \( 4 \) z nich są gorszej jakości. Inspektor losowo wybiera \( 5 \).  Jakie jest prawdopodobieństwo, że niższej jakości jest co najwyżej jeden z wybranych tranzystorów? Zaokrąglij wynik do dwóch miejsc po przecinku.}
{\( \frac{\binom{46}5 + \binom{46}4\cdot\binom41}{\binom{50}5}\,\dot{=}\,0{,}96 \)}{\( \frac{\frac{46!}{41!}+\frac{46!}{42!}}{\frac{50!}{45!}}\,\dot{=}\,0{,}66 \)}{\( \frac{\binom{46}5 + \binom{46}4}{\binom{50}5}\,\dot{=}\,0{,}72 \)}{\( \frac{\frac{46!}{41!}+\frac{46!}{42!}\cdot \frac{4!}{3!}}{\frac{50!}{45!}}\,\dot{=}\,0{,}71 \)}{}{}}
\LANGcs{
\Question{
V kontejneru je \( 50 \) výrobků, z nichž \( 4 \) jsou 2. jakosti. Jaká je pravděpodobnost, že mezi \( 5 \) náhodně vybranými výrobky je nejvýše jeden 2. jakosti? Výsledek zaokrouhlete s přesností na setiny.}
{\( \frac{\binom{46}5 + \binom{46}4\cdot\binom41}{\binom{50}5}\,\dot{=}\,0{,}96 \)}{\( \frac{\frac{46!}{41!}+\frac{46!}{42!}}{\frac{50!}{45!}}\,\dot{=}\,0{,}66 \)}{\( \frac{\binom{46}5 + \binom{46}4}{\binom{50}5}\,\dot{=}\,0{,}72 \)}{\( \frac{\frac{46!}{41!}+\frac{46!}{42!}\cdot \frac{4!}{3!}}{\frac{50!}{45!}}\,\dot{=}\,0{,}71 \)}{}{}}
\LANGsk{
\Question{
V kontajneri je \( 50 \) výrobkov, z nich \( 4 \) sú 2. kvality. Aká je pravdepodobnosť, že medzi \( 5 \) náhodne vybranými výrobkami je najviac jeden 2. kvality? Výsledok zaokrúhlite s presnosťou na stotiny.}
{\( \frac{\binom{46}5 + \binom{46}4\cdot\binom41}{\binom{50}5}\,\dot{=}\,0{,}96 \)}{\( \frac{\frac{46!}{41!}+\frac{46!}{42!}}{\frac{50!}{45!}}\,\dot{=}\,0{,}66 \)}{\( \frac{\binom{46}5 + \binom{46}4}{\binom{50}5}\,\dot{=}\,0{,}72 \)}{\( \frac{\frac{46!}{41!}+\frac{46!}{42!}\cdot \frac{4!}{3!}}{\frac{50!}{45!}}\,\dot{=}\,0{,}71 \)}{}{}}
\LANGes{
\Question{
En una recipiente hay \( 50 \) productos, de cuales \( 4 \) son de peor calidad. Vamos a sacar \( 5 \)  productos al azar. ¿Qué es la probabilidad de sacar por lo máximo un producto de peor calidad?  Aproxima el resultados a dos cifras decimales.}
{\( \frac{\binom{46}5 + \binom{46}4\cdot\binom41}{\binom{50}5}\,\dot{=}\,0{,}96 \)}{\( \frac{\frac{46!}{41!}+\frac{46!}{42!}}{\frac{50!}{45!}}\,\dot{=}\,0{,}66 \)}{\( \frac{\binom{46}5 + \binom{46}4}{\binom{50}5}\,\dot{=}\,0{,}72 \)}{\( \frac{\frac{46!}{41!}+\frac{46!}{42!}\cdot \frac{4!}{3!}}{\frac{50!}{45!}}\,\dot{=}\,0{,}71 \)}{}{}}
\End

\Data\ID{1003041603}
\Updated{2021-02-03 03:02:38}
\Level{C}
\SubArea{Probability}
\LANGen{
\Question{
There are \( 30 \) students in a class, of them are \( 14 \) girls and the rest are boys. The teacher selects two students for weekly routine help. If selection is done at random, what is the probability that these students are not two girls? Round the result to two decimal places.}
{\( \frac{\binom{16}2+\binom{16}1\cdot\binom{14}1}{\binom{30}2}\,\dot{=}\,0{.}79 \)}{\( \frac{\binom{16}2}{\binom{30}2}\,\dot{=}\,0{.}28 \)}{\( \frac{\binom{14}2}{\binom{30}2}\,\dot{=}\,0{.}21 \)}{\( \frac{\binom{16}1\cdot\binom{14}1}{\binom{30}2}\,\dot{=}\,0{.}51 \)}{}{}}
\LANGpl{
\Question{
W klasie jest  \( 30 \) uczniów, w tym jest  \( 14 \) dziewczyn.  Nauczyciel wybiera dwóch uczniów do odpowiedzi.  Wybór  odbywa się losowo, jakie jest prawdopodobieństwo, że ci uczniowie  nie są dwiema dziewczynami? Zaokrąglij wynik do dwóch miejsc po przecinku.}
{\( \frac{\binom{16}2+\binom{16}1\cdot\binom{14}1}{\binom{30}2}\,\dot{=}\,0{,}79 \)}{\( \frac{\binom{16}2}{\binom{30}2}\,\dot{=}\,0{,}28 \)}{\( \frac{\binom{14}2}{\binom{30}2}\,\dot{=}\,0{,}21 \)}{\( \frac{\binom{16}1\cdot\binom{14}1}{\binom{30}2}\,\dot{=}\,0{,}51 \)}{}{}}
\LANGcs{
\Question{
Ve třídě je \( 30 \) žáků, \( 14 \) děvčat a \( 16 \) chlapců. Učitel z nich náhodně vybere dva na týdenní službu. Jaká je pravděpodobnost, že to nebudou dvě děvčata? Výsledek zaokrouhlete na dvě desetinná místa.}
{\( \frac{\binom{16}2+\binom{16}1\cdot\binom{14}1}{\binom{30}2}\,\dot{=}\,0{,}79 \)}{\( \frac{\binom{16}2}{\binom{30}2}\,\dot{=}\,0{,}28 \)}{\( \frac{\binom{14}2}{\binom{30}2}\,\dot{=}\,0{,}21 \)}{\( \frac{\binom{16}1\cdot\binom{14}1}{\binom{30}2}\,\dot{=}\,0{,}51 \)}{}{}}
\LANGsk{
\Question{
V triede je \( 30 \) žiakov, \( 14 \) dievčat a \( 16 \) chlapcov. Učiteľ z nich náhodne vyberá dvoch na službu týždenníkov. Aká je pravdepodobnosť, že to nebudú dve dievčatá? Výsledok zaokrúhlite na dve desatinné miesta.}
{\( \frac{\binom{16}2+\binom{16}1\cdot\binom{14}1}{\binom{30}2}\,\dot{=}\,0{,}79 \)}{\( \frac{\binom{16}2}{\binom{30}2}\,\dot{=}\,0{,}28 \)}{\( \frac{\binom{14}2}{\binom{30}2}\,\dot{=}\,0{,}21 \)}{\( \frac{\binom{16}1\cdot\binom{14}1}{\binom{30}2}\,\dot{=}\,0{,}51 \)}{}{}}
\LANGes{
\Question{
En una clase hay \( 30 \) alumnos, \( 14 \) chicas y \( 16 \) chicos. El profesor elige dos estudiantes al azar para ser la ayuda de la semana. ¿Qué es la probabilidad de que no sean dos chicas? Aproxima el resultado a dos cifras decimales.}
{\( \frac{\binom{16}2+\binom{16}1\cdot\binom{14}1}{\binom{30}2}\,\dot{=}\,0{,}79 \)}{\( \frac{\binom{16}2}{\binom{30}2}\,\dot{=}\,0{,}28 \)}{\( \frac{\binom{14}2}{\binom{30}2}\,\dot{=}\,0{,}21 \)}{\( \frac{\binom{16}1\cdot\binom{14}1}{\binom{30}2}\,\dot{=}\,0{,}51 \)}{}{}}
\End

\Data\ID{1003158301}
\Updated{2021-02-04 03:02:10}
\Level{C}
\SubArea{Probability}
\LANGen{
\Question{
A deck of playing cards contains \( 4 \) aces, \( 12 \) face cards and \( 16 \) number cards. If you draw two cards, what is the probability that exactly one of the drawn cards is an ace or exactly one of the drawn cards is a face? Round the result to four decimal places.}
{\( 0{.}6129 \)}{\( 0{.}7097 \)}{\( 0{.}3065 \)}{\( 0{.}3548 \)}{}{}}
\LANGpl{
\Question{
Talia kart zawiera \( 4 \) asy, \( 12 \) kart z figurami oraz  \( 16 \) kart z cyframi. Jeśli wybierzesz dwie karty, jakie jest prawdopodobieństwo, że dokładnie jedna z wylosowanych kart to as lub karta z figurą? Zaokrąglij wynik do czterech miejsc po przecinku.}
{\( 0{,}6129 \)}{\( 0{,}7097 \)}{\( 0{,}3065 \)}{\( 0{,}3548 \)}{}{}}
\LANGcs{
\Question{
Balíček karet obsahuje  \( 4 \) esa, \( 12 \) karet s figurami a \( 16 \) karet s čísly. Z balíčku vytáhneme najednou dvě karty. Určete pravděpodobnost, že mezi vybranými kartami bude právě jedno eso nebo právě jedna karta s figurou. Výsledek zaokrouhlete na čtyři desetinná místa.}
{\( 0{,}6129 \)}{\( 0{,}7097 \)}{\( 0{,}3065 \)}{\( 0{,}3548 \)}{}{}}
\LANGsk{
\Question{
Balíček kariet obsahuje  \( 4 \) esá, \( 12 \) figúr a \( 16 \) kariet s číslami. Z balíčka vytiahneme naraz dve karty. Určte pravdepodobnosť, že medzi vybratými kartami bude práve jedno eso alebo práve jedna figúra. Výsledok zaokrúhlite na štyri desatinné miesta.}
{\( 0{,}6129 \)}{\( 0{,}7097 \)}{\( 0{,}3065 \)}{\( 0{,}3548 \)}{}{}}
\LANGes{
\Question{
Un paquete de cartas contiene \( 4 \) as, \( 12 \) figuras y \( 16 \) números. Vamos a sacar dos cartas del paquete. Averigua la probabilidad de sacar exactamente un as o exactamente una figura. Aproxima el resultado a cuatro cifras decimales.}
{\( 0{,}6129 \)}{\( 0{,}7097 \)}{\( 0{,}3065 \)}{\( 0{,}3548 \)}{}{}}
\End

\Data\ID{1003158302}
\Updated{2021-02-04 01:02:50}
\Level{C}
\SubArea{Probability}
\LANGen{
\Question{
What is the probability that among \( 10 \) boys from a class born in the same year (\( 365 \) days) there are at least two having the same birthday? Round the result to four decimal places.}
{\( 0{.}1169 \)}{\( 0{.}1619 \)}{\( 0{.}1961 \)}{\( 0{.}1916 \)}{\( 0{.}1196 \)}{\( 0{.}1691 \)}}
\LANGpl{
\Question{
Jakie jest prawdopodobieństwo, że wśród \( 10 \) chłopców urodzonych w tym samym roku (\( 365 \) dni) są co najmniej dwie osoby, które mają urodziny w tym samym dniu? Zaokrąglij wynik do czterech miejsc po przecinku.}
{\( 0{,}1169 \)}{\( 0{,}1619 \)}{\( 0{,}1961 \)}{\( 0{,}1916 \)}{\( 0{,}1196 \)}{\( 0{,}1691 \)}}
\LANGcs{
\Question{
Jaká je pravděpodobnost, že z \( 10 \) chlapců jedné třídy, narozených v témže roce (\( 365 \) dní), mají alespoň dva narozeniny ve stejný den? Výsledek zaokrouhlete na čtyři desetinná místa.}
{\( 0{,}1169 \)}{\( 0{,}1619 \)}{\( 0{,}1961 \)}{\( 0{,}1916 \)}{\( 0{,}1196 \)}{\( 0{,}1691 \)}}
\LANGsk{
\Question{
Aká je pravdepodobnosť, že z \( 10 \) chlapcov jednej triedy, narodených v tom istom roku (\( 365 \) dní), majú aspoň dvaja narodeniny v ten istý deň? Výsledok zaokrúhlite na štyri desatinné miesta.}
{\( 0{,}1169 \)}{\( 0{,}1619 \)}{\( 0{,}1961 \)}{\( 0{,}1916 \)}{\( 0{,}1196 \)}{\( 0{,}1691 \)}}
\LANGes{
\Question{
¿Qué es la probabilidad de que de los \( 10 \) chicos de una clase nacidos en el mismo año (\( 365 \) días), por lo mínimo dos celebran cumpleaños en el mismo día? Aproxima el resultado a cuatro cifras decimales.}
{\( 0{,}1169 \)}{\( 0{,}1619 \)}{\( 0{,}1961 \)}{\( 0{,}1916 \)}{\( 0{,}1196 \)}{\( 0{,}1691 \)}}
\End

\Data\ID{1003158303}
\Updated{2021-02-08 12:02:00}
\Level{C}
\SubArea{Probability}
\LANGen{
\Question{
If you roll a dice ten times, what is the probability of getting one exactly five times? Round the result to four decimal places.}
{\( 0{.}0130 \)}{\( 0{.}0033 \)}{\( 0{.}5000 \)}{\( 0{.}1984 \)}{\( 0{.}0833 \)}{}}
\LANGpl{
\Question{
Jeśli rzucisz  kostką  dziesięć razy, jakie jest prawdopodobieństwo wyrzucenia jedynki  dokładnie pięć razy? Zaokrąglij wynik do czterech miejsc po przecinku.}
{\( 0{,}0130 \)}{\( 0{,}0033 \)}{\( 0{,}5000 \)}{\( 0{,}1984 \)}{\( 0{,}0833 \)}{}}
\LANGcs{
\Question{
Hodíme \( 10 \) krát kostkou. Jaká je pravděpodobnost, že jednička padne právě \( 5 \) krát? Výsledek zapište s přesností na čtyři desetinná místa.}
{\( 0{,}0130 \)}{\( 0{,}0033 \)}{\( 0{,}5000 \)}{\( 0{,}1984 \)}{\( 0{,}0833 \)}{}}
\LANGsk{
\Question{
Aká je pravdepodobnosť, že ak \( 10 \) krát hodíme kockou, padne jednotka práve \( 5 \) krát? Výsledok zapíšte s presnosťou na štyri desatinné miesta.}
{\( 0{,}0130 \)}{\( 0{,}0033 \)}{\( 0{,}5000 \)}{\( 0{,}1984 \)}{\( 0{,}0833 \)}{}}
\LANGes{
\Question{
Vamos a tirar el dado \( 10 \) veces. ¿Qué es la probabilidad de que saquemos uno exactamente \( 5 \) veces? Aproxima el resultado a cuatro cifras decimales.}
{\( 0{,}0130 \)}{\( 0{,}0033 \)}{\( 0{,}5000 \)}{\( 0{,}1984 \)}{\( 0{,}0833 \)}{}}
\End

\Data\ID{1003158304}
\Updated{2021-02-10 09:02:59}
\Level{C}
\SubArea{Probability}
\LANGen{
\Question{
If eight coins are tossed together, what is the probability of getting at least two tails? Round the result to four decimal places.}
{\( 0{.}9648 \)}{\( 0{.}0352 \)}{\( 0{.}0273 \)}{\( 0{.}9727 \)}{\( 0{.}8750 \)}{}}
\LANGpl{
\Question{
Osiem monet  rzucono jednocześnie, jakie jest prawdopodobieństwo wyrzucenia co najmniej dwóch reszek? Zaokrąglij wynik do czterech miejsc po przecinku.}
{\( 0{,}9648 \)}{\( 0{,}0352 \)}{\( 0{,}0273 \)}{\( 0{,}9727 \)}{\( 0{,}8750 \)}{}}
\LANGcs{
\Question{
Najednou hodíme \( 8 \) mincemi. Jaká je pravděpodobnosti, že alespoň na dvou z nich padne rub? Výsledek zaokrouhlete na čtyři desetinná místa.}
{\( 0{,}9648 \)}{\( 0{,}0352 \)}{\( 0{,}0273 \)}{\( 0{,}9727 \)}{\( 0{,}8750 \)}{}}
\LANGsk{
\Question{
Aká je pravdepodobnosť, že ak naraz hodíme \( 8 \) mincí, padne aspoň na dvoch z nich rub mince? Výsledok zaokrúhlite na štyri desatinné miesta.}
{\( 0{,}9648 \)}{\( 0{,}0352 \)}{\( 0{,}0273 \)}{\( 0{,}9727 \)}{\( 0{,}8750 \)}{}}
\LANGes{
\Question{
Vamos lanzar \( 8 \) monedas a la vez. ¿Qué probabilidad hay de que por lo mínimo dos de ellas tengan cruz? Aproxima el resultado a cuatro cifras decimales.}
{\( 0{,}9648 \)}{\( 0{,}0352 \)}{\( 0{,}0273 \)}{\( 0{,}9727 \)}{\( 0{,}8750 \)}{}}
\End

\Data\ID{1003158305}
\Updated{2021-02-10 12:02:53}
\Level{C}
\SubArea{Probability}
\LANGen{
\Question{
If you roll a dice ten times, what is the probability of getting a prime number at most two times? Round the result to four decimal places.}
{\( 0{.}0547 \)}{\( 0{.}0030 \)}{\( 0{.}0439 \)}{\( 0{.}0034 \)}{\( 0{.}1000 \)}{}}
\LANGpl{
\Question{
Jeśli rzucisz kostką  dziesięć razy, jakie jest prawdopodobieństwo wyrzucenia liczby pierwszej najwyżej dwa razy? Zaokrąglij wynik do czterech miejsc po przecinku.}
{\( 0{,}0547 \)}{\( 0{,}0030 \)}{\( 0{,}0439 \)}{\( 0{,}0034 \)}{\( 0{,}1000 \)}{}}
\LANGcs{
\Question{
Hodíme \( 10 \) krát kostkou. Jaká je pravděpodobnost, že prvočíslo padne nejvýše \( 2 \) krát? Výsledek zapište s přesností na čtyři desetinná místa.}
{\( 0{,}0547 \)}{\( 0{,}0030 \)}{\( 0{,}0439 \)}{\( 0{,}0034 \)}{\( 0{,}1000 \)}{}}
\LANGsk{
\Question{
Aká je pravdepodobnosť, že ak \( 10 \) krát hodíme kockou, padne na nej prvočíslo najviac \( 2 \) krát? Výsledok zapíšte s presnosťou na štyri desatinné miesta.}
{\( 0{,}0547 \)}{\( 0{,}0030 \)}{\( 0{,}0439 \)}{\( 0{,}0034 \)}{\( 0{,}1000 \)}{}}
\LANGes{
\Question{
Vamos a lanzar una moneda \( 10 \) veces. ¿Qué probabilidad hay de un número primo cae por lo máximo \( 2 \) veces? Aproxima el resultado a cuatro cifras decimales.}
{\( 0{,}0547 \)}{\( 0{,}0030 \)}{\( 0{,}0439 \)}{\( 0{,}0034 \)}{\( 0{,}1000 \)}{}}
\End

\Data\ID{1003158306}
\Updated{2021-02-08 12:02:49}
\Level{C}
\SubArea{Probability}
\LANGen{
\Question{
What is the probability that at least three of $6$ dices show \( 6 \)? Round the result to four decimal places.}
{\( 0{.}0623 \)}{\( 0{.}0536 \)}{\( 0{.}0226 \)}{\( 0{.}0356 \)}{\( 0{.}0262 \)}{\( 0{.}0635 \)}}
\LANGpl{
\Question{
Jakie jest prawdopodobieństwo, że co najmniej na  trzech kostkach zostanie wyrzucona \( 6 \)? Zaokrągli wynik do czterech miejsc po przecinku.}
{\( 0{,}0623 \)}{\( 0{,}0536 \)}{\( 0{,}0226 \)}{\( 0{,}0356 \)}{\( 0{,}0262 \)}{\( 0{,}0635 \)}}
\LANGcs{
\Question{
Hodíme šesti kostkami. Jaká je pravděpodobnost, že na nich padnou alespoň \( 3 \) šestky? Výsledek zaokrouhlete na čtyři desetinná místa.}
{\( 0{,}0623 \)}{\( 0{,}0536 \)}{\( 0{,}0226 \)}{\( 0{,}0356 \)}{\( 0{,}0262 \)}{\( 0{,}0635 \)}}
\LANGsk{
\Question{
Aká je pravdepodobnosť, že ak hodíme šiestimi kockami, budú na nich aspoň \( 3 \) šestky? Výsledok zaokrúhlite na štyri desatinné miesta.}
{\( 0{,}0623 \)}{\( 0{,}0536 \)}{\( 0{,}0226 \)}{\( 0{,}0356 \)}{\( 0{,}0262 \)}{\( 0{,}0635 \)}}
\LANGes{
\Question{
Vamos a tirar seis dados. ¿Qué es la probabilidad de que saquemos por lo mínimo \( 3 \) seis? Aproxima el resultado a cuatro cifras decimales.}
{\( 0{,}0623 \)}{\( 0{,}0536 \)}{\( 0{,}0226 \)}{\( 0{,}0356 \)}{\( 0{,}0262 \)}{\( 0{,}0635 \)}}
\End

\Data\ID{1003158307}
\Updated{2021-02-10 02:02:56}
\Level{C}
\SubArea{Probability}
\LANGen{
\Question{
Suppose that the success rate of one specific medical treatment is \( 90\,\% \). If the treatment is given to \( 20 \) new patients, what is the probability that it is effective in at least \( 18 \) of them? Round the result to four decimal places.}
{\( 0{.}6769 \)}{\( 0{.}9000 \)}{\( 0{.}2852 \)}{\( 0{.}7148 \)}{\( 0{.}8100 \)}{}}
\LANGpl{
\Question{
Załóżmy, że wskaźnik sukcesu jednego konkretnego leczenia medycznego wynosi \( 90\,\% \). Jeśli leczeniu poddano \( 20 \) nowych pacjentów,  jakie jest prawdopodobieństwo jego skuteczności
co najmniej u  \( 18 \) z nich? Zaokrągli wynik do czterech miejsc po przecinku.}
{\( 0{,}6769 \)}{\( 0{,}9000 \)}{\( 0{,}2852 \)}{\( 0{,}7148 \)}{\( 0{,}8100 \)}{}}
\LANGcs{
\Question{
Předpokládejme, že určitý lék má účinnost \( 90\,\% \), tj. vyléčí \( 90\,\% \) pacientů. Jaká je pravděpodobnost, že když lék podáme \( 20 \) pacientům, tak alespoň \( 18 \) z nich se vyléčí? Výsledek zapište s přesností na čtyři desetinná místa.}
{\( 0{,}6769 \)}{\( 0{,}9000 \)}{\( 0{,}2852 \)}{\( 0{,}7148 \)}{\( 0{,}8100 \)}{}}
\LANGsk{
\Question{
Určitý liek úspešne lieči \( 90\,\% \) prípadov. Aká je pravdepodobnosť, že ak ho podáme \( 20 \) pacientom, tak sa aspoň \( 18 \) z nich sa vylieči? Výsledok zapíšte s presnosťou na štyri desatinné miesta.}
{\( 0{,}6769 \)}{\( 0{,}9000 \)}{\( 0{,}2852 \)}{\( 0{,}7148 \)}{\( 0{,}8100 \)}{}}
\LANGes{
\Question{
Vamos a suponer que una medicina tiene efectividad de \( 90\,\% \), ósea cura \( 90\,\% \) de enfermos, y vamos a darla a \( 20 \) enfemos. ¿Qué probabilidad hay de que curamos por lo mínimo \( 18 \)? Aproxima el resultado a cuatro cifras decimales.}
{\( 0{,}6769 \)}{\( 0{,}9000 \)}{\( 0{,}2852 \)}{\( 0{,}7148 \)}{\( 0{,}8100 \)}{}}
\End

\Data\ID{1003158308}
\Updated{2021-02-10 02:02:06}
\Level{C}
\SubArea{Probability}
\LANGen{
\Question{
The probability that the randomly selected product is first-class quality is \( 0{.}12 \). Determine the probability that at least \( 2 \) out of \( 50 \) randomly selected products are first-class quality. Round the result to four decimal places.}
{\( 0{.}9869 \)}{\( 0{.}9689 \)}{\( 0{.}8969 \)}{\( 0{.}8699 \)}{\( 0{.}9896 \)}{\( 0{.}8996 \)}}
\LANGpl{
\Question{
Prawdopodobieństwo, że losowo wybrany produkt jest najwyższej jakości wynosi  \( 0{,}12 \). Określ prawdopodobieństwo, że co najmniej \( 2 \) z \( 50 \) losowo wybranych produktów są najwyższej jakości. Zaokrąglij wynik do czterech miejsc po przecinku.}
{\( 0{,}9869 \)}{\( 0{,}9689 \)}{\( 0{,}8969 \)}{\( 0{,}8699 \)}{\( 0{,}9896 \)}{\( 0{,}8996 \)}}
\LANGcs{
\Question{
Pravděpodobnost toho, že náhodně vybraný výrobek bude první jakosti je \( 0{,}12 \). Náhodně vybereme \( 50 \) výrobků. Jaká je pravděpodobnost, že alespoň \( 2 \) z nich budou první jakosti? Výsledek zaokrouhlete na čtyři desetinná místa.}
{\( 0{,}9869 \)}{\( 0{,}9689 \)}{\( 0{,}8969 \)}{\( 0{,}8699 \)}{\( 0{,}9896 \)}{\( 0{,}8996 \)}}
\LANGsk{
\Question{
Pravdepodobnosť toho, že náhodne vybratý výrobok bude prvej akosti je \( 0{,}12 \). Náhodne vyberieme \( 50 \) výrobkov. Aká je pravdepodobnosť toho, že aspoň \( 2 \) z nich budú prvej akosti? Výsledok zaokrúhlite na štyri desatinné miesta.}
{\( 0{,}9869 \)}{\( 0{,}9689 \)}{\( 0{,}8969 \)}{\( 0{,}8699 \)}{\( 0{,}9896 \)}{\( 0{,}8996 \)}}
\LANGes{
\Question{
La probabilidad de que un producto elegido al azar sea de alta calidad es \( 0{,}12 \). Vamos a elegir \( 50 \) productos. ¿Qué probabilidad hay de que por lo mínimo \( 2 \) de ellos sean de alta calidad? Aproxima el resultado a cuatro cifras decimales.}
{\( 0{,}9869 \)}{\( 0{,}9689 \)}{\( 0{,}8969 \)}{\( 0{,}8699 \)}{\( 0{,}9896 \)}{\( 0{,}8996 \)}}
\End

\Data\ID{1003158309}
\Updated{2021-02-16 12:02:45}
\Level{C}
\SubArea{Probability}
\LANGen{
\Question{
Students of a class take a multiple choice test consisting of \( 10 \) tasks. There are \( 5 \) optional answers to each of the tasks, while only one answer is correct. One student, however, did not study for the test at all. Therefore, he circles his answers randomly, without performing any calculations. Find the probability that he will circle at least \( 3 \) correct answers. Round the result to four decimal places.}
{\( 0{.}3222 \)}{\( 0{.}8591 \)}{\( 0{.}1409 \)}{\( 0{.}6778 \)}{}{}}
\LANGpl{
\Question{
Uczniowie  przystępują do testu wielokrotnego wyboru składającego się z \( 10 \) zadań. W każdym zadaniu jest  \( 5 \) odpowiedzi do wyboru, tylko jedna jest poprawna. Jeden  z uczniów nie przygotował się do testu dlatego zakreśla odpowiedzi losowo, nie wykonuje żadnych obliczeń. Jakie jest prawdopodobieństwo, że wybierze co najmniej  \( 3 \) poprawne odpowiedzi?  Zaokrągli wynik do czterech miejsc po przecinku.}
{\( 0{,}3222 \)}{\( 0{,}8591 \)}{\( 0{,}1409 \)}{\( 0{,}6778 \)}{}{}}
\LANGcs{
\Question{
Ve třídě byl žákům zadán test s \( 10 \) otázkami. U každé z nich bylo uvedených \( 5 \) možných variant odpovědí, přičemž vždy právě jedna byla správná. Petr se na test vůbec nepřipravoval a odpovědi na otázky volil zcela náhodně. Jaká je pravděpodobnost, že Petr označil alespoň \( 3 \) správné odpovědi? Výsledek zaokrouhlete na čtyři desetinná místa.}
{\( 0{,}3222 \)}{\( 0{,}8591 \)}{\( 0{,}1409 \)}{\( 0{,}6778 \)}{}{}}
\LANGsk{
\Question{
V triede bol žiakom zadaný test s \( 10 \) úlohami a ku každej z nich bolo uvedených \( 5 \) možných výsledkov, pričom vždy len jeden bol správny. Jeden žiak sa však na test vôbec nepripravoval, preto krúžkuje svoje výsledky náhodne a bez výpočtov. Aká je pravdepodobnosť, že bude mať správne aspoň \( 3 \) výsledky? Výsledok zaokrúhlite na štyri desatinné miesta.}
{\( 0{,}3222 \)}{\( 0{,}8591 \)}{\( 0{,}1409 \)}{\( 0{,}6778 \)}{}{}}
\LANGes{
\Question{
Los alumnos de una clase hacen un examen de \( 10 \) preguntas. Cada una de las preguntas tiene \( 5 \) posibilidades de respuestas pero solamente una es correcta. Pedro no ha estudiado nada y ha elegido las respuestas al azar. ¿Qué probabilidad hay de que Pedro marcó por lo mínimo \( 3 \) respuestas correctas? Aproxima el resultado a cuatro cifras decimales.}
{\( 0{,}3222 \)}{\( 0{,}8591 \)}{\( 0{,}1409 \)}{\( 0{,}6778 \)}{}{}}
\End

\Data\ID{1003158406}
\Updated{2021-02-16 10:02:20}
\Level{C}
\SubArea{Probability}
\LANGen{
\Question{
A bag contains \( 10 \) white balls and \( 5 \) black balls. Two balls are drawn at random one by one without replacement. Find the probability that two black balls are drawn.}
{\( \frac2{21} \)}{\( \frac2{15} \)}{\( \frac14 \)}{\( \frac19 \)}{}{}}
\LANGpl{
\Question{
W pudełku znajduje się  \( 10 \) białych kul i  \( 5 \) czarnych. Z pudełka losowo wyciągnięto dwie kule jedna po drugiej, pierwszej nie włożono z powrotem do pudełka. Wskaż prawdopodobieństwo wylosowania dwóch czarnych kul.}
{\( \frac2{21} \)}{\( \frac2{15} \)}{\( \frac14 \)}{\( \frac19 \)}{}{}}
\LANGcs{
\Question{
V urně je \( 10 \) bílých a \( 5 \) černých koulí. Z urny vybereme postupně dvě koule, přičemž vybrané koule do urny nevracíme. Určete pravděpodobnost, že jsme vybrali dvě černé koule.}
{\( \frac2{21} \)}{\( \frac2{15} \)}{\( \frac14 \)}{\( \frac19 \)}{}{}}
\LANGsk{
\Question{
V urne je \( 10 \) bielych a \( 5 \) čiernych gulí. Z urny postupne vyberieme dve gule, pričom prvú vybratú guľu do urny už nevrátime. Určte pravdepodobnosť, že sme vybrali dve čierne gule.}
{\( \frac2{21} \)}{\( \frac2{15} \)}{\( \frac14 \)}{\( \frac19 \)}{}{}}
\LANGes{
\Question{
En una urna hay \( 10 \) bolas blancas y \( 5 \) negras. Vamos a sacar dos bolas sin devolverlas a la urna. Calcula la probabilidad de sacar dos bolas negras.}
{\( \frac2{21} \)}{\( \frac2{15} \)}{\( \frac14 \)}{\( \frac19 \)}{}{}}
\End

\Data\ID{1003158407}
\Updated{2021-02-17 03:02:35}
\Level{C}
\SubArea{Probability}
\LANGen{
\Question{
The long-term records of a car dealer show that a customer getting a new car buys within special equipment offers a parking assistant system (PAS) with the probability \( 50\,\% \) and that he buys a xenon flashlight with the probability \( 20\,\% \). Both items of special equipment (PAS and xenon lamps) are purchased by the customer with the probability \( 10\,\% \). What is the probability that the customer bought xenon lamps if you know he bought PAS?}
{\( 20\,\% \)}{\( 60\,\% \)}{\( 10\,\% \)}{\( 80\,\% \)}{}{}}
\LANGpl{
\Question{
Długoterminowy rejestr sprzedawcy samochodów  wskazuje, że klient kupujący nowy samochód w ramach dodatkowego wyposażenia zakupi również  system wspomagania  parkowania   (PAS) z prawdopodobieństwem  \( 50\,\% \) natomiast lampy ksenonowe z prawdopodobieństwem \( 20\,\% \). Prawdopodobieństwo zakupu obu elementów wyposażenia   (PAS i lamp ksenonowych) wynosi \( 10\,\% \).  Jakie jest prawdopodobieństwo, że klient zakupi lampy ksenonowe, jeśli wiemy, że kupił PAS?}
{\( 20\,\% \)}{\( 60\,\% \)}{\( 10\,\% \)}{\( 80\,\% \)}{}{}}
\LANGcs{
\Question{
Z dlouhodobých záznamů prodejce automobilů vyplývá, že zákazník, který si kupuje nový automobil, si v rámci zvláštní výbavy koupí s \( 50\% \)  pravděpodobností parkovacího asistenta (PAS) a s \( 20\% \) pravděpodobností xenonové svítilny. Obě položky zvláštní výbavy (PAS i xenonové svítilny) si zákazník koupí s pravděpodobností \( 10\,\% \). S jakou pravděpodobností si zákazník koupil xenonové svítilny, víte-li, že  si koupil PAS?}
{\( 20\,\% \)}{\( 60\,\% \)}{\( 10\,\% \)}{\( 80\,\% \)}{}{}}
\LANGsk{
\Question{
Z dlhodobých záznamov predajcu automobilov vyplýva, že zákazník, ktorý si kupuje nový automobil, si v rámci zvláštnej výbavy kúpi s \( 50\,\%\)-nou pravdepodobnosťou parkovacieho asistenta (PAS) a s \( 20\,\%\)-nou pravdepodobnosťou xenónové svetlá. Obidve položky zvláštnej výbavy (PAS i xenónové svetlá) si zákazník kúpi s pravdepodobnosťou \( 10\,\% \). S akou pravdepodobnosťou si zákazník kúpil xenónové svetlá, ak viete, že si kúpil PAS?}
{\( 20\,\% \)}{\( 60\,\% \)}{\( 10\,\% \)}{\( 80\,\% \)}{}{}}
\LANGes{
\Question{
Un vendedor de coches tiene registro de ventas y deuce que un cliente que compra un coche nuevo compra un sistema asistente de parking (PAS) con una probabilidad de \( 50\% \)  y una luz xenon con probabilidad de \( 20\% \). Un cliente compra ambos (PAS y luz xenon) con probabilidad de \( 10\,\% \). ¿Qué probabilidad hay de que un cliente compró la luz xenon sabiendo que ha comprado PAS?}
{\( 20\,\% \)}{\( 60\,\% \)}{\( 10\,\% \)}{\( 80\,\% \)}{}{}}
\End

\Data\ID{1003158408}
\Updated{2021-02-16 11:02:40}
\Level{C}
\SubArea{Probability}
\LANGen{
\Question{
The class consists of \( 10\,\% \) long haired boys, \( 30\,\% \) short haired boys, \( 50\,\% \) long haired girls and \( 10\,\% \) short haired girls. We choose one of the students from the class randomly. Find the probability that the student has long hair if you know he is a boy.}
{\( 0{.}25 \)}{\( 0{.}40 \)}{\( 0{.}10 \)}{\( 0{.}03 \)}{}{}}
\LANGpl{
\Question{
W klasie jest  \( 10\,\% \) chłopców z długimi włosami, \( 30\,\% \) chłopców z krótkimi włosami, \( 50\,\% \) dziewczyn z długimi włosami oraz  \( 10\,\% \) dziewczyn z krótkimi włosami. Losowo wybieramy jednego ucznia z klasy. Jakie jest prawdopodobieństwo wyboru ucznia z długimi włosami jeśli wiemy, że jest to chłopak?}
{\( 0{,}25 \)}{\( 0{,}40 \)}{\( 0{,}10 \)}{\( 0{,}03 \)}{}{}}
\LANGcs{
\Question{
Třída je složena z \( 10\,\% \) dlouhovlasých chlapců, \( 30\,\% \) krátkovlasých chlapců, \( 50\,\% \) dlouhovlasých dívek a \( 10\,\% \) krátkovlasých dívek. Náhodně vybereme ze třídy jednu osobu. Určete pravděpodobnost, že tato osoba má dlouhé vlasy, víte-li, že je to chlapec.}
{\( 0{,}25 \)}{\( 0{,}40 \)}{\( 0{,}10 \)}{\( 0{,}03 \)}{}{}}
\LANGsk{
\Question{
V jednej triede je \( 10\,\% \) dlhovlasých chlapcov, \( 30\,\% \) krátkovlasých chlapcov, \( 50\,\% \) dlhovlasých dievčat a \( 10\,\% \) krátkovlasých dievčat. Náhodne z triedy vyberieme jednu osobu. Určte pravdepodobnosť, že táto osoba má dlhé vlasy, pričom viete, že je to chlapec.}
{\( 0{,}25 \)}{\( 0{,}40 \)}{\( 0{,}10 \)}{\( 0{,}03 \)}{}{}}
\LANGes{
\Question{
Una clase está compuesta por \( 10\,\% \) de chicos de pelo largo,  \( 30\,\% \) de chicos de pelo corto, \( 50\,\% \) de chicas de pelo largo y \( 10\,\% \) de chicas de pelo corto. Vamos a eligir una persona de la clase al azar. Calcula la probabilidad de que esta persona tenga pelo largo sabiendo que es un chico.}
{\( 0{,}25 \)}{\( 0{,}40 \)}{\( 0{,}10 \)}{\( 0{,}03 \)}{}{}}
\End

\Data\ID{1103158401}
\Updated{2021-02-10 02:02:06}
\Level{C}
\SubArea{Probability}
\IMG{1103158401.pdf}
\LANGen{
\Question{
We roll red and yellow dice. Find the probability that the yellow die is two given that the sum of both dice is eight. (Note: In the following table, the sums of points on both dice are listed.)}
{\( \frac15 \)}{\( \frac16 \)}{\( \frac1{36} \)}{\( \frac5{36} \)}{}{}}
\LANGpl{
\Question{
Rzucamy żółtą i czerwoną kostką. Jakie jest prawdopodobieństwo, że na żółtej kostce wypadnie dwa wiedząc, że suma wyrzuconych liczb na obu kostkach jest równa osiem?  (Uwaga: w poniższej tabeli podano sumy punktów na obu kostkach.)}
{\( \frac15 \)}{\( \frac16 \)}{\( \frac1{36} \)}{\( \frac5{36} \)}{}{}}
\LANGcs{
\Question{
Hodíme červenou a žlutou kostkou. Určete pravděpodobnost, že na žluté kostce padly dva body, víte-li, že součet počtů bodů, které na obou kostkách padly, je osm.  (Poznámka: Pro výpočet můžete použít následující tabulku, v níž jsou uvedeny součty počtu bodů na obou kostkách.)}
{\( \frac15 \)}{\( \frac16 \)}{\( \frac1{36} \)}{\( \frac5{36} \)}{}{}}
\LANGsk{
\Question{
Hodíme červenú a žltú kocku. Určte pravdepodobnosť, že na žltej kocke padnú dva body, ak viete, že súčet počtu bodov, ktoré padli na oboch kockách, je osem.
\[ \]
Poznámka: Pre výpočet môžete použiť následujúcu tabuľku, v ktorej sú uvedené všetky možné súčty počtu bodov padnutých na dvoch kockách.}
{\( \frac15 \)}{\( \frac16 \)}{\( \frac1{36} \)}{\( \frac5{36} \)}{}{}}
\LANGes{
\Question{
Vamos a tirar un dado rojo y un amarillo. Calcula la probabilidad de que en el dado amarillo estén dos puntos si sabemos que la suma de los puntos de ambos dados es ocho.  (Pista: Para calcular puedes usar la siguiente tabla, donde hay las sumas de los puntos de ambos dados.)}
{\( \frac15 \)}{\( \frac16 \)}{\( \frac1{36} \)}{\( \frac5{36} \)}{}{}}
\End

\Data\ID{1103158402}
\Updated{2021-02-10 02:02:11}
\Level{C}
\SubArea{Probability}
\IMG{1103158402.pdf}
\LANGen{
\Question{
We roll red and yellow dice. Let the event $A$ be: On the red die a number greater than $2$  is rolled. And let the event $B$ be: The sum of points rolled on both  dice is greater than $6$. Find \( P(A|B) \). (Note: In the table the sums of points on both dice are listed.)}
{\( \frac34 \)}{\( \frac12 \)}{\( \frac67 \)}{\( \frac14 \)}{}{}}
\LANGpl{
\Question{
Rzucamy żółtą i czerwoną kostką. Zdarzenie $A$: Na czerwonej kostce wyrzucono liczbę większą niż $2$. Zdarzenie $B$: Suma liczb wyrzuconych na obu kostkach jest większa niż $6$.
Wyznacz \( P(A|B) \). (Uwaga: W tabeli podano sumy punktów na obu kostkach.)}
{\( \frac34 \)}{\( \frac12 \)}{\( \frac67 \)}{\( \frac14 \)}{}{}}
\LANGcs{
\Question{
Hodíme červenou a žlutou kostkou. Označme jako jev $A$ tvrzení, že na červené kostce padly více než $2$ body a jako jev $B$ tvrzení, že součet počtu bodů na obou kostkách je větší než $6$.
Určete \( P(A|B) \). Pro výpočet můžete použít tabulku, v níž jsou uvedeny součty počtu bodů na obou kostkách.}
{\( \frac34 \)}{\( \frac12 \)}{\( \frac67 \)}{\( \frac14 \)}{}{}}
\LANGsk{
\Question{
Hodíme červenú a žltú kocku. Označme ako $A$, že na červenej kocke padnú viac ako $2$ a ako $B$ súčet počtu bodov na oboch kockách je väčší ako $6$. Určte \( P(A|B) \).\[\]Poznámka: Pre výpočet môžete použiť tabuľku, v ktorej sú uvedené všetky možné súčty počtu bodov padnutých na dvoch kockách.}
{\( \frac34 \)}{\( \frac12 \)}{\( \frac67 \)}{\( \frac14 \)}{}{}}
\LANGes{
\Question{
Vamos a tirar un dado rojo y un amarillo.  Vamos a marcar el evento $A$ como la declaración que en el dado rojo tenemos más de $2$ puntos y el evento $B$ la declaración que la suma de los puntos en ambos dados es mayor de $6$.
Averigua \( P(A|B) \).  (Pista: Para calcular puedes usar la siguiente tabla, donde hay las sumas de los puntos de ambos dados.)}
{\( \frac34 \)}{\( \frac12 \)}{\( \frac67 \)}{\( \frac14 \)}{}{}}
\End

\Data\ID{1103158403}
\Updated{2021-02-16 12:02:59}
\Level{C}
\SubArea{Probability}
\IMG{1103158403_0.pdf}
\LANGen{
\Question{
A box (see the picture) contains \( 5 \) red balls and \( 7 \) green balls which are marked with numbers. A ball is drawn uniformly at random from the box. Find the probability that the ball is marked with even number given that the ball is red. (Round the result to two decimal places.)}
{\( 0{.}60 \)}{\( 0{.}50 \)}{\( 0{.}83 \)}{\( 0{.}25 \)}{}{}}
\LANGpl{
\Question{
W pudełku (spójrz na rysunek) znajduje się  \( 5 \) czerwonych kul oraz \( 7 \) zielonych kul. Kule oznaczone są liczbami. Losowo wybieramy jedną kulę z pudełka. Jakie jest prawdopodobieństwo, że kula jest oznaczona liczbą parzystą wiedząc, że kula jest czerwona? (Zaokrągli wynik do dwóch miejsc po przecinku.)}
{\( 0{,}60 \)}{\( 0{,}50 \)}{\( 0{,}83 \)}{\( 0{,}25 \)}{}{}}
\LANGcs{
\Question{
V urně (viz obrázek) je  \( 5 \) červených a \( 7 \) zelených koulí s čísly. Z urny náhodně vybereme jednu kouli. Určete pravděpodobnost, že na kouli je sudé číslo, víte-li, že koule má červenou barvu. Výsledek zaokrouhlete na dvě desetinná místa.}
{\( 0{,}60 \)}{\( 0{,}50 \)}{\( 0{,}83 \)}{\( 0{,}25 \)}{}{}}
\LANGsk{
\Question{
V urne (viď obrázok) je  \( 5 \) červených a \( 7 \) zelených gulí s číslami. Z urny náhodne vyberieme jednu guľu. Určte pravdepodobnosť, že na guli je párne číslo, pričom viete, že guľa má červenú farbu. Výsledok zaokrúhlite na dve desatinné miesta.}
{\( 0{,}60 \)}{\( 0{,}50 \)}{\( 0{,}83 \)}{\( 0{,}25 \)}{}{}}
\LANGes{
\Question{
En una urna hay \( 5 \) bolas rojas y \( 7 \) bolas verdes con números. Vamos a sacar una bola al azar. Calcula la probabilidad de sacar una bola con número par sabiendo que la bola es roja. Aproxima el resultado a dos cifras decimales.}
{\( 0{,}60 \)}{\( 0{,}50 \)}{\( 0{,}83 \)}{\( 0{,}25 \)}{}{}}
\End

\Data\ID{1103158404}
\Updated{2021-02-16 12:02:05}
\Level{C}
\SubArea{Probability}
\IMG{1103158404_0.pdf}
\LANGen{
\Question{
A box (see the picture) contains \( 5 \) red balls and \( 7 \) green balls which are marked with numbers. A ball is drawn uniformly at random from the box. Let the event $A$ be: The drawn ball is green. And let the event $B$ be: The drawn ball is marked with the number which is greater than $6$. Find \( P(A|B) \). (Round the result to two decimal places.)}
{\( 0{.}50 \)}{\( 0{.}43 \)}{\( 0{.}25 \)}{\( 0{.}83 \)}{}{}}
\LANGpl{
\Question{
Pudełko (spójrz na rysunek) zawiera \( 5 \) czerwonych kul oraz \( 7 \) zielonych kul. Kule oznaczone są liczbami. Z pudełka wyciągamy losowo jedną kulę. Zdarzenie $A$: Wylosowana kula jest zielona. Zdarzenie $B$: Wylosowana kula jest oznaczona liczbą większą niż $6$. Wskaż \( P(A|B) \). (Zaokrągli wynik do dwóch miejsc po przecinku.)}
{\( 0{,}50 \)}{\( 0{,}43 \)}{\( 0{,}25 \)}{\( 0{,}83 \)}{}{}}
\LANGcs{
\Question{
V urně (viz obrázek) je \( 5 \) červených koulí a \( 7 \) zelených koulí s čísly. Z urny náhodně vybereme jednu kouli. Označme jako jev $A$ tvrzení, že náhodně vybraná koule má zelenou barvu, a jako jev $B$ tvrzení, že na náhodně vybrané kouli je číslo větší než $6$. Určete \( P(A|B) \). Výsledek zaokrouhlete na dvě desetinná místa.}
{\( 0{,}50 \)}{\( 0{,}43 \)}{\( 0{,}25 \)}{\( 0{,}83 \)}{}{}}
\LANGsk{
\Question{
V urne (viď obrázok) je  \( 5 \) červených a \( 7 \) zelených gulí s číslami. Z urny náhodne vyberieme jednu guľu. Označme ako $A$, že náhodne vybratá guľa má zelenú farbu, a ako $B$, že na náhodne vybratej guli je číslo väčšie ako $6$. Určte \( P(A|B) \). Výsledok zaokrúhlite na dve desatinné miesta.}
{\( 0{,}50 \)}{\( 0{,}43 \)}{\( 0{,}25 \)}{\( 0{,}83 \)}{}{}}
\LANGes{
\Question{
En una urna hay \( 5 \) bolas rojas y \( 7 \) bolas verdes con números (vea el dibujo). Vamos a sacar una bola al azar. Vamos a marcar la declaración que la bola es verde como el evento $A$ y el evento $B$ es la declaración que el número sacado es mayor que $6$. Calcula \( P(A|B) \). Aproxima el resultado a dos cifras decimales.}
{\( 0{,}50 \)}{\( 0{,}43 \)}{\( 0{,}25 \)}{\( 0{,}83 \)}{}{}}
\End

\Data\ID{1103158405}
\Updated{2021-02-16 11:02:17}
\Level{C}
\SubArea{Probability}
\IMG{1103158405.pdf}
\LANGen{
\Question{
Let's roll a special die with the numbers placed according to the given scheme. Find the probability that the result is not \( 6 \) given that the result is greater than \( 2 \).}
{\( \frac34 \)}{\( \frac56 \)}{\( \frac15 \)}{\( \frac14 \)}{}{}}
\LANGpl{
\Question{
Rzućmy specjalną kostkę z cyframi umieszczonymi zgodnie z danym schematem. Jakie jest  prawdopodobieństwo, że wynik nie wynosi  \( 6 \) wiedząc, że jest większy niż \( 2 \)?}
{\( \frac34 \)}{\( \frac56 \)}{\( \frac15 \)}{\( \frac14 \)}{}{}}
\LANGcs{
\Question{
Mějme hrací kostku se speciálním pláštěm (viz obrázek). Určete pravděpodobnost, že na kostce nepadla šestka, víte-li, že padlo číslo větší než dvě.}
{\( \frac34 \)}{\( \frac56 \)}{\( \frac15 \)}{\( \frac14 \)}{}{}}
\LANGsk{
\Question{
Majme hraciu kocku s následujúcim plášťom. Určte pravdepodobnosť, že na kocke nepadla šestka, ak viete, že padlo číslo väčšie ako \( 2 \).}
{\( \frac34 \)}{\( \frac56 \)}{\( \frac15 \)}{\( \frac14 \)}{}{}}
\LANGes{
\Question{
Tenemos un dado con una composición especial (vea el dibujo). Calcula la probabilidad de que no ha caído un seis sabiendo que ha caído un número mayor que dos.}
{\( \frac34 \)}{\( \frac56 \)}{\( \frac15 \)}{\( \frac14 \)}{}{}}
\End

\Data\ID{9000138303}
\Updated{2021-02-16 02:02:58}
\Level{C}
\SubArea{Probability}
\LANGen{
\Question{
Two dices are rolled. Find the probability that we get the number
\(6\) 
on just one of the dices and the sum of the numbers on both dices is
\(8\).}
{\(\frac{2}
{36}\,\dot{=}\,0{.}0556\)}{\(\frac{5}
{36}\,\dot{=}\,0{.}1389\)}{\(\frac{11}
{36}\,\dot{=}\,0{.}3056\)}{\(\frac{14}
{36}\,\dot{=}\,0{.}3889\)}{}{}}
\LANGpl{
\Question{
Rzucamy dwoma  kostkami. Jakie jest prawdopodobieństwo, że tylko na jednej kostce wyrzucimy
\(6\) i
suma liczb na obu kostkach będzie równa 
\(8\)?}
{\(\frac{2}
{36}\,\dot{=}\,0{,}0556\)}{\(\frac{5}
{36}\,\dot{=}\,0{,}1389\)}{\(\frac{11}
{36}\,\dot{=}\,0{,}3056\)}{\(\frac{14}
{36}\,\dot{=}\,0{,}3889\)}{}{}}
\LANGcs{
\Question{
Hodíme dvěma kostkami. Jaká je pravděpodobnost, že součet bude
\(8\) a právě na jedné
kostce padne \(6\)?}
{\(\frac{2}
{36}\,\dot{=}\,0{,}0556\)}{\(\frac{5}
{36}\,\dot{=}\,0{,}1389\)}{\(\frac{11}
{36}\,\dot{=}\,0{,}3056\)}{\(\frac{14}
{36}\,\dot{=}\,0{,}3889\)}{}{}}
\LANGsk{
\Question{
Hodíme dvoma kockami. Aká je pravdepodobnosť, že práve na jednej kocke padne 
\(6\) 
a súčet bude
\(8\)?}
{\(\frac{2}
{36}\,\dot{=}\,0{,}0556\)}{\(\frac{5}
{36}\,\dot{=}\,0{,}1389\)}{\(\frac{11}
{36}\,\dot{=}\,0{,}3056\)}{\(\frac{14}
{36}\,\dot{=}\,0{,}3889\)}{}{}}
\LANGes{
\Question{
Vamos a tirar dos dados. ¿Qué probabilidad hay de que la suma sea 
\(8\) y exactamente un dado tenga \(6\)?}
{\(\frac{2}
{36}\,\dot{=}\,0{,}0556\)}{\(\frac{5}
{36}\,\dot{=}\,0{,}1389\)}{\(\frac{11}
{36}\,\dot{=}\,0{,}3056\)}{\(\frac{14}
{36}\,\dot{=}\,0{,}3889\)}{}{}}
\End

\Data\ID{9000138305}
\Updated{2021-02-16 02:02:01}
\Level{C}
\SubArea{Probability}
\IMG{001383_kostky2_1.png}
\LANGen{
\Question{
Two different dices (a white dice and a black dice) are rolled. We get the sum of the numbers
on both dices \(6\).
Find the probability that there is an even number on the black dice.}
{\(\frac{2}
{5}=0{.}4\)}{\(\frac{5}
{36}\,\dot{=}\,0{.}1389\)}{\(\frac{5}
{18}\,\dot{=}\,0{.}2778\)}{\(\frac{13}
{36}\,\dot{=}\,0{.}3611\)}{}{}}
\LANGpl{
\Question{
Rzucamy dwoma  kostkami (białą i czarną). Suma liczb na obu kostkach jest równa  \(6\).
Jakie jest prawdopodobieństwo, że na kostce czarnej zostanie wyrzucona liczba parzysta?}
{\(\frac{2}
{5}=0{,}4\)}{\(\frac{5}
{36}\,\dot{=}\,0{,}1389\)}{\(\frac{5}
{18}\,\dot{=}\,0{,}2778\)}{\(\frac{13}
{36}\,\dot{=}\,0{,}3611\)}{}{}}
\LANGcs{
\Question{
Házíme dvěma kostkami, bílou a černou. Jaká je pravděpodobnost, že na
černé kostce padne sudé číslo, víte-li, že součet na obou kostkách je
\(6\)?}
{\(\frac{2}
{5}=0{,}4\)}{\(\frac{5}
{36}\,\dot{=}\,0{,}1389\)}{\(\frac{5}
{18}\,\dot{=}\,0{,}2778\)}{\(\frac{13}
{36}\dot{=}\,0{,}3611\)}{}{}}
\LANGsk{
\Question{
Hodíme dvoma kockami, bielou a čiernou. Súčet na kockách je \(6\). Aká je pravdepodobnosť, že na čiernej kocke padlo párne číslo?}
{\(\frac{2}
{5}=0{,}4\)}{\(\frac{5}
{36}\,\dot{=}\,0{,}1389\)}{\(\frac{5}
{18}\,\dot{=}\,0{,}2778\)}{\(\frac{13}
{36}\,\dot{=}\,0{,}3611\)}{}{}}
\LANGes{
\Question{
Vamos a tirar dos dados, un negro y un blanco. ¿Qué probabilidad hay de que en el dado negro ha caído un número par si la suma de ambos dados es
\(6\)?}
{\(\frac{2}
{5}=0{,}4\)}{\(\frac{5}
{36}\,\dot{=}\,0{,}1389\)}{\(\frac{5}
{18}\,\dot{=}\,0{,}2778\)}{\(\frac{13}
{36}\dot{=}\,0{,}3611\)}{}{}}
\End

\Data\ID{9000138308}
\Updated{2021-02-03 03:02:38}
\Level{C}
\SubArea{Probability}
\IMG{001383_kostky2_2.png}
\LANGen{
\Question{
Two different dices (a white dice and a black dice) are rolled. The sum of the numbers on both dices is
\(8\). Find the probability
that there is \(4\) 
on the black dice.}
{\(\frac{1}
{5}=0{.}2\)}{\(\frac{1}
{4}=0{.}25\)}{\(\frac{6}
{36}\,\dot{=}\,0{.}1667\)}{\(\frac{11}
{36}\,\dot{=}\,0{.}3056\)}{}{}}
\LANGpl{
\Question{
Rzucamy dwoma  kostkami (białą i czarną). Suma liczb na obu kostkach jest równa 
\(8\). Jakie jest prawdopodobieństwo, że na kostce czarnej zostanie wyrzucona liczba  \(4\)?}
{\(\frac{1}
{5}=0{,}2\)}{\(\frac{1}
{4}=0{,}25\)}{\(\frac{6}
{36}\,\dot{=}\,0{,}1667\)}{\(\frac{11}
{36}\,\dot{=}\,0{,}3056\)}{}{}}
\LANGcs{
\Question{
Hodíme dvěma kostkami, bílou a černou. Jaká je pravděpodobnost, že na černé kostce
padne \(4\) za předpokladu,
že součet bude \(8\)?}
{\(\frac{1}
{5}=0{,}2\)}{\(\frac{1}
{4}=0{,}25\)}{\(\frac{6}
{36}\,\dot{=}\,0{,}1667\)}{\(\frac{11}
{36}\,\dot{=}\,0{,}3056\)}{}{}}
\LANGsk{
\Question{
Hodíme dvoma kockami, bielou a čiernou. Aká je pravdepodobnosť, že na čiernej kocke padne \(4\), ak súčet bude 
\(8\)?}
{\(\frac{1}
{5}=0{,}2\)}{\(\frac{1}
{4}=0{,}25\)}{\(\frac{6}
{36}\,\dot{=}\,0{,}1667\)}{\(\frac{11}
{36}\,\dot{=}\,0{,}3056\)}{}{}}
\LANGes{
\Question{
Vamos a tirar dos dados, un blanco y un negro. ¿Qué es la probabilidad de que el dado negro lleve \(4\) si la suma es \(8\)?}
{\(\frac{1}
{5}=0{,}2\)}{\(\frac{1}
{4}=0{,}25\)}{\(\frac{6}
{36}\,\dot{=}\,0{,}1667\)}{\(\frac{11}
{36}\,\dot{=}\,0{,}3056\)}{}{}}
\End

\Data\ID{9000138310}
\Updated{2021-02-03 03:02:08}
\Level{C}
\SubArea{Probability}
\LANGen{
\Question{
Two dices are rolled. The sum of the numbers on both dices is
\(9\). Find the probability
there is the number \(4\) 
on one of the dices.}
{\(\frac{2}
{4}=0{.}5\)}{\(\frac{12}
{36}\,\dot{=}\,0{.}3333\)}{\(\frac{10}
{36}\,\dot{=}\,0{.}2778\)}{\(\frac{8}
{12}\,\dot{=}\,0{.}6667\)}{}{}}
\LANGpl{
\Question{
Rzucamy dwoma  kostkami. Suma liczb na obu kostkach jest równa 
\(9\). Jakie jest prawdopodobieństwo, że wyrzucimy liczbę \(4\) 
na jednej z kostek?}
{\(\frac{2}
{4}=0{,}5\)}{\(\frac{12}
{36}\,\dot{=}\,0{,}3333\)}{\(\frac{10}
{36}\,\dot{=}\,0{,}2778\)}{\(\frac{8}
{12}\,\dot{=}\,0{,}6667\)}{}{}}
\LANGcs{
\Question{
Hodíme dvěma kostkami. Jaká je pravděpodobnost, že na jedné kostce padne
\(4\) za předpokladu,
že součet bude \(9\)?}
{\(\frac{2}
{4}=0{,}5\)}{\(\frac{12}
{36}\,\dot{=}\,0{,}3333\)}{\(\frac{10}
{36}\,\dot{=}\,0{,}2778\)}{\(\frac{8}
{12}\,\dot{=}\,0{,}6667\)}{}{}}
\LANGsk{
\Question{
Hodíme dvoma kockami. Aká je pravdepodobnosť, že na jednej kocke padne
 \(4\), ak súčet bude \(9\) ?}
{\(\frac{2}
{4}=0{,}5\)}{\(\frac{12}
{36}\,\dot{=}\,0{,}3333\)}{\(\frac{10}
{36}\,\dot{=}\,0{,}2778\)}{\(\frac{8}
{12}\,\dot{=}\,0{,}6667\)}{}{}}
\LANGes{
\Question{
Tiramos dos dedos y la suma de los puntos es \(9\). ¿Qué es la probabilidad de que un dedo tiene \(4\) puntos?}
{\(\frac{2}
{4}=0{,}5\)}{\(\frac{12}
{36}\,\dot{=}\,0{,}3333\)}{\(\frac{10}
{36}\,\dot{=}\,0{,}2778\)}{\(\frac{8}
{12}\,\dot{=}\,0{,}6667\)}{}{}}
\End

\Data\ID{9000154801}
\Updated{2021-02-17 01:02:40}
\Level{C}
\SubArea{Probability}
\LANGen{
\Question{
There are six money transports through the Sherwood forest. Robin Hood knows
that two of the transports are secured by soldiers. Find the respective probabilities
that if Robin's band attacks two random transports, then none, one and both
transports will be secured by the soldiers.}
{\(\frac{6}
{15};\, \frac{8} 
{15};\, \frac{1} 
{15}\)}{\(\frac{3}
{9};\, \frac{5} 
{9};\, \frac{1} 
{9}\)}{\(\frac{1}
{3};\, \frac{2} 
{3};\, \frac{2} 
{3}\)}{\(\frac{1}
{2};\, \frac{1} 
{4};\, \frac{1} 
{4}\)}{}{}}
\LANGpl{
\Question{
W lesie Sherwood jest sześć transportów  pieniężnych. Robin Hood wie,
że dwa transporty są zabezpieczone przez żołnierzy. Jakie jest prawdopodobieństwo tego, że jeśli  
 Robin Hood zaatakuje  dwa losowo wybrane transporty to  żaden, jeden  lub oba transporty 
będą zabezpieczone przez żołnierzy?}
{\(\frac{6}
{15};\, \frac{8} 
{15};\, \frac{1} 
{15}\)}{\(\frac{3}
{9};\, \frac{5} 
{9};\, \frac{1} 
{9}\)}{\(\frac{1}
{3};\, \frac{2} 
{3};\, \frac{2} 
{3}\)}{\(\frac{1}
{2};\, \frac{1} 
{4};\, \frac{1} 
{4}\)}{}{}}
\LANGcs{
\Question{
Robin Hood zná cestu šesti vozů s penězi. Ví, že dva jsou hlídané
vojáky. Jaké jsou postupně pravděpodobnosti, že ze dvou vozů, které
přepadne, nebude hlídaný žádný, bude hlídán právě jeden, resp. budou
hlídány vojáky oba přepadené vozy?}
{\(\frac{6}
{15};\, \frac{8} 
{15};\, \frac{1} 
{15}\)}{\(\frac{3}
{9};\, \frac{5} 
{9};\, \frac{1} 
{9}\)}{\(\frac{1}
{3};\, \frac{2} 
{3};\, \frac{2} 
{3}\)}{\(\frac{1}
{2};\, \frac{1} 
{4};\, \frac{1} 
{4}\)}{}{}}
\LANGsk{
\Question{
Robin Hood pozná cestu šiestich vozov s peniazmi. Vie, že dva vozy sú strážené vojakmi. Aké sú jednotlivé pravdepodobnosti, že z dvoch vozov, ktoré prepadne, nebude strážený žiaden, bude strážený práve jeden, resp. budú strážené vojakmi oba prepadnuté vozy?}
{\(\frac{6}
{15};\, \frac{8} 
{15};\, \frac{1} 
{15}\)}{\(\frac{3}
{9};\, \frac{5} 
{9};\, \frac{1} 
{9}\)}{\(\frac{1}
{3};\, \frac{2} 
{3};\, \frac{2} 
{3}\)}{\(\frac{1}
{2};\, \frac{1} 
{4};\, \frac{1} 
{4}\)}{}{}}
\LANGes{
\Question{
Robin Hood conoce las rutas de seis carros con dinero y sabe que dos de ellos están protegidos por soldados. ¿Qué probabilidades hay de que de los dos carros que va a robar no va a estar protegido ninguno, estará protegido exactamente uno o estarán protegidos ambos, respectivamente?}
{\(\frac{6}
{15};\, \frac{8} 
{15};\, \frac{1} 
{15}\)}{\(\frac{3}
{9};\, \frac{5} 
{9};\, \frac{1} 
{9}\)}{\(\frac{1}
{3};\, \frac{2} 
{3};\, \frac{2} 
{3}\)}{\(\frac{1}
{2};\, \frac{1} 
{4};\, \frac{1} 
{4}\)}{}{}}
\End

\Data\ID{9000154802}
\Updated{2021-02-17 02:02:01}
\Level{C}
\SubArea{Probability}
\LANGen{
\Question{
Three hundred soldiers know details related to the weapon transport to Nottingham.
The probability that a soldier betrays the sheriff and tells the details to Robin Hood
is \(0.01\).
This probability is fixed for all soldiers. Robin tries to find out the details on the
transport by asking each soldier. Find the probability that Robin will find out details
(i.e. at least one soldier tells the secret to Robin). Round your answer to three
decimal places.}
{\(0.951\)}{\(0.049\)}{\(0.827\)}{\(0.173\)}{}{}}
\LANGpl{
\Question{
Trzystu żołnierzy zna szczegóły dotyczące  transportu broni do  Nottingham.
Prawdopodobieństwo zdrady żołnierza i wyjawienie szczegółów  transportu   Robinowi  
wynosi  \(0{,}01\). Prawdopodobieństwo jest takie same dla wszystkich żołnierzy. Robin próbuje poznać szczegóły 
transportu, pytając każdego żołnierza. Jakie jest  prawdopodobieństwo tego,  że Robin pozna  szczegóły
(tzn. przynajmniej jeden żołnierz zdradzi tajemnicę)? Zaokrągli odpowiedź do trzech miejsc po przecinku.}
{\(0{,}951\)}{\(0{,}049\)}{\(0{,}827\)}{\(0{,}173\)}{}{}}
\LANGcs{
\Question{
O dodávce zbraní do Nottinghamu ví celkem \(300\) 
vojáků. Pravděpodobnost, že kterýkoliv voják šerifa zradí a prozradí Robinu Hoodovi
trasu, je \(0{,}01\).
Jaká je pravděpodobnost, že se Robinovi podaří zjistit
trasu dodávky alespoň od jednoho vojáka? Výsledek zaokrouhlete na
\(3\) 
desetinná místa.}
{\(0{,}951\)}{\(0{,}049\)}{\(0{,}827\)}{\(0{,}173\)}{}{}}
\LANGsk{
\Question{
O dodávke zbraní vie celkom tristo vojakov.
Pravdepodobnosť, že ktorýkoľvek vojak zradí šerifa a prezradí Robinovi trasu, je \(0{,}01\).
Táto pravdepodobnosť je rovnaká pre všetkých vojakov. Aká je pravdepodobnosť, že sa Robinovi  podarí zistiť trasu dodávky aspoň od jedného vojaka?
 Výsledok zaokrúhlite na tri desatinné miesta.}
{\(0{,}951\)}{\(0{,}049\)}{\(0{,}827\)}{\(0{,}173\)}{}{}}
\LANGes{
\Question{
\(300\) soldados saben sobre el importe de armas a Nottingham. La probabilidad de que cualquier soldado traiciona al sherif y le dice la ruta a Robin Hood es \(0{,}01\).
¿Qué probabilidad hay de que Robin Hood se entera de la ruta de por lo mínimo un soldado? Aproxima el resultado a
\(3\) 
cifras decimales.}
{\(0{,}951\)}{\(0{,}049\)}{\(0{,}827\)}{\(0{,}173\)}{}{}}
\End

\Data\ID{9000154804}
\Updated{2021-02-17 10:02:34}
\Level{C}
\SubArea{Probability}
\LANGen{
\Question{
Robin Hood wants to have \(6\) 
children with his love Maid Marian. Find the probability that they will have
\(2\) girls
and \(4\) 
boys. The probability that one child will be a girl is
\(48.79\%\) and the probability
of a boy is \(51.21\%\).
Round your answer to three decimal places.}
{\(0.246\)}{\(0.222\)}{\(0.015\)}{\(0.016\)}{}{}}
\LANGpl{
\Question{
Robin Hood chce mieć  \(6\) 
dzieci z   Marian. Jakie jest prawdopodobieństwo tego, że będą mieć 
\(2\) dziewczynki
i \(4\) 
chłopców? Prawdopodobieństwo urodzenia dziewczynki wynosi
\(48{,}79\%\) a chłopaka
\(51{,}21\%\).
Zaokrągli odpowiedź do trzech miejsc po przecinku.}
{\(0{,}246\)}{\(0{,}222\)}{\(0{,}015\)}{\(0{,}016\)}{}{}}
\LANGcs{
\Question{
Robin Hood má s Marian celkem \(6\) 
dětí. Jaká je pravděpodobnost, že se jim narodily
\(2\) dívky
a \(4\) 
chlapci? Uvažujte pravděpodobnost narození dívky
\(48{,}79\%\) a chlapce
\(51{,}21\%\). Výsledek
zaokrouhlete na \(3\) 
desetinná místa.}
{\(0{,}246\)}{\(0{,}222\)}{\(0{,}015\)}{\(0{,}016\)}{}{}}
\LANGsk{
\Question{
Robin Hood chce mať \(6\) 
detí so slečnou Mariannou. Aká je pravdepodobnosť, že sa im narodia
\(2\) dievčatá
a \(4\) 
chlapci? Pravdepodobnosť narodenia dievčaťa je
\(48{,}79\%\) a pravdepodobnosť narodenia chlapca je \(51{,}21\%\).
Výsledok zaokrúhlite na 3 desatinné miesta.}
{\(0{,}246\)}{\(0{,}222\)}{\(0{,}015\)}{\(0{,}016\)}{}{}}
\LANGes{
\Question{
Robin Hood y Marian tienen \(6\) 
hijos. ¿Qué probabilidad hay de que tengan  
\(2\) chicas
y \(4\) 
chicos? Vamos a tomar la probabilidad de tener una chica como 
\(48{,}79\%\) y de tener un chico como
\(51{,}21\%\). Aproxima el resultado a \(3\) cifras decimales.}
{\(0{,}246\)}{\(0{,}222\)}{\(0{,}015\)}{\(0{,}016\)}{}{}}
\End

\Data\ID{9000154805}
\Updated{2021-02-17 10:02:41}
\Level{C}
\SubArea{Probability}
\LANGen{
\Question{
A boy plays Monopoly game. He is in the jail and has to roll three times a pair of
dices. To escape from the jail he needs the number six on both dices. Find the
probability that he succeeds to escape the jail. Round your answer to three decimal
places.}
{\(0.081\)}{\(0.919\)}{\(0.028\)}{\(0.095\)}{}{}}
\LANGpl{
\Question{
Jaś gra w  Monopoly. Jest w więzieniu i musi rzucić  trzy razy parą kostek. Wyjście z więzienia zapewni mu wyrzucenie liczby sześć na obu kostkach. Ile wynosi prawdopodobieństwo opuszczenia przez niego więzienia?
Zaokrągli odpowiedź do trzech miejsc po przecinku.}
{\(0{,}081\)}{\(0{,}919\)}{\(0{,}028\)}{\(0{,}095\)}{}{}}
\LANGcs{
\Question{
Robin Hood hraje Monopoly. Je ve vězení a hází třikrát dvěma
kostkami. Aby se z vězení dostal, musí mu alespoň jednou padnout dvě
šestky. Jaká je šance, že se toto stane? Výsledek zaokrouhlete na
\(3\) 
desetinná místa.}
{\(0{,}081\)}{\(0{,}919\)}{\(0{,}028\)}{\(0{,}095\)}{}{}}
\LANGsk{
\Question{
Robin hrá Monopoly. Je vo väzení a hádže trikrát dvoma kockami. Aby sa z väzenia dostal, musia mu aspoň raz padnúť dve šestky. Aká je
pravdepodobnosť, že sa mu podarí dostať sa z väzenia? Výsledok zaokrúhlite na tri desatinné miesta.}
{\(0{,}081\)}{\(0{,}919\)}{\(0{,}028\)}{\(0{,}095\)}{}{}}
\LANGes{
\Question{
Robin Hood juega a Monopoly. Está en la prisión y tiene que tirar dos dados. Para salir de la prisión tiene que tirar dos seis. ¿Qué probabilidad hay de que tenga éxito? 
Aproxima el resultado a \(3\) cifras decimales.}
{\(0{,}081\)}{\(0{,}919\)}{\(0{,}028\)}{\(0{,}095\)}{}{}}
\End

\Data\ID{9000154809}
\Updated{2021-02-03 03:02:51}
\Level{C}
\SubArea{Probability}
\LANGen{
\Question{
The soldiers play a card game with a pack of
\(32\) cards.
Find the probability that there is at least one ace among three randomly selected
cards. (There are four aces in the pack.) Round your answer to two decimal places.}
{\(0.34\)}{\(0.66\)}{\(0.57\)}{\(0.43\)}{}{}}
\LANGpl{
\Question{
Jakie jest prawdopodobieństwo wystąpienia co najmniej jednego asa spośród trzech  losowo wybranych kart z \(32\)? Zaokrągli wynik do dwóch  miejsc po przecinku.}
{\(0{,}34\)}{\(0{,}66\)}{\(0{,}57\)}{\(0{,}43\)}{}{}}
\LANGcs{
\Question{
Jaká je pravděpodobnost, že mezi \(3\) náhodně vybranými kartami z balíčku \(32\) karet bude alespoň jedno eso? Výsledek zaokrouhlete na \(2\) desetinná místa.}
{\(0{,}34\)}{\(0{,}66\)}{\(0{,}57\)}{\(0{,}43\)}{}{}}
\LANGsk{
\Question{
Aká je pravdepodobnosť, že medzi troma náhodne vybranými kartami z balíčku
\(32\) kariet bude aspoň jedno eso?
Výsledok zaokrúhlite na dve desatinné miesta.}
{\(0{,}34\)}{\(0{,}66\)}{\(0{,}57\)}{\(0{,}43\)}{}{}}
\LANGes{
\Question{
Vamos a sacar \(3\) cartas al azar de un paquete de \(32\) cartas. ¿Qué es la probabilidad de sacar por lo mínimo un as? Aproxima el resultado a dos cifras decimales.}
{\(0{,}34\)}{\(0{,}66\)}{\(0{,}57\)}{\(0{,}43\)}{}{}}
\End
